% Accidents en zones montagneuses ou liés aux mouvements de terrain — SVTEEHB 3ème — Cours signé OnBuch+
% Programme officiel MINESEC. Compiler avec : tectonic N13-accidents-zones-montagneuses-ou-lies-mouvements-terrain.tex
\newcommand{\DOCMATIERE}{SVTEEHB}
\newcommand{\DOCNIVEAU}{3ème}
\newcommand{\DOCMODULE}{SVTEEHB – classe de 3ème}
\newcommand{\DOCLECON}{N13}
\newcommand{\DOCTITRE}{Accidents en zones montagneuses ou liés aux mouvements de terrain}
% OnBuch+ — préambule « cours signés » ENRICHI (compilé avec tectonic / XeLaTeX)
% Système visuel « Pop » d'OnBuch+ : fond crème (jamais blanc), bordures encre
% épaisses, ombres « dures » décalées, titres Archivo Black en capitales.
% Version MATHÉMATIQUES : graphes (pgfplots/tikz), boîtes pédagogiques
% (définition, propriété, méthode, attention, le savais-tu, exercice résolu,
% exercice, TP, bilan), page de garde et signature OnBuch+ ancrée en bas.
%
% Macros attendues dans main.tex AVANT \input{preamble} :
%   \DOCMATIERE  \DOCNIVEAU  \DOCMODULE  \DOCLECON (numéro)  \DOCTITRE
\documentclass[11pt]{article}

\usepackage{fontspec}
\usepackage[french]{babel}
\usepackage[a4paper,margin=2cm,top=3.4cm,bottom=2.8cm,headheight=1.4cm,footskip=1.3cm]{geometry}
\usepackage{amsmath,amssymb,amsfonts}
\usepackage{mathtools} % \xmapsto, \coloneqq... souvent utilisés par les modèles
\usepackage{cancel} % simplifications barrées (bilans de forces)
% \abs{z} (module, valeur absolue) et \norm{u} (norme), utilisés par les modèles
\providecommand{\abs}{}\let\abs\relax\DeclarePairedDelimiter{\abs}{\lvert}{\rvert}
\providecommand{\norm}{}\let\norm\relax\DeclarePairedDelimiter{\norm}{\lVert}{\rVert}
\usepackage[table]{xcolor} % [table] : fournit aussi \rowcolors (alternance de lignes)
\usepackage{pagecolor}
\usepackage{tikz}
\usetikzlibrary{arrows.meta,positioning,calc,shapes.geometric,shapes.misc,shapes.arrows,decorations.pathmorphing,decorations.pathreplacing,decorations.markings,patterns,shadows,mindmap,trees,fit,backgrounds,matrix,intersections,angles,quotes}
\usepackage{pgfplots}
\usepackage[version=4]{mhchem} % équations nucléaires \ce{^{235}_{92}U}
\usepackage[european,siunitx]{circuitikz} % schémas électriques
\pgfplotsset{compat=1.18}
\usepgfplotslibrary{fillbetween,groupplots} % aires entre courbes (calcul intégral)
\usepackage{enumitem}
\usepackage{fancyhdr}
% [most] charge toutes les bibliothèques tcolorbox (listings, minted, raster,
% documentation...) et gonfle inutilement la mémoire TeX sur un document long
% (~300 boîtes) : on ne charge que ce qui est réellement utilisé.
\usepackage[skins,breakable]{tcolorbox}
\usepackage{graphicx}
\usepackage{colortbl}
\usepackage{booktabs}
\usepackage{tabularx}
\usepackage{multirow}
\usepackage{diagbox} % case d'en-tête diagonale des tableaux à double entrée
% Colonnes extensibles centrée / gauche / droite (C, L, R), souvent utilisées par les modèles
\newcolumntype{C}{>{\centering\arraybackslash}X}
\newcolumntype{L}{>{\raggedright\arraybackslash}X}
\newcolumntype{R}{>{\raggedleft\arraybackslash}X}
\usepackage{array}
\usepackage{multicol}
\usepackage{siunitx}
% parse-numbers=false : accepte aussi \SI{2,5 \times 10^{-3}}{...}
\sisetup{locale=FR,output-decimal-marker={,},group-minimum-digits=4,parse-numbers=false}
\DeclareSIUnit\ohm{\ensuremath{\symup{\Omega}}} % U+2126 absent de la police
\DeclareSIUnit\division{div} % réglages d'oscilloscope (ms/div, V/div)
\DeclareSIUnit\tr{tr} % tours (vitesse de rotation, ex. \SI{1500}{\tr\per\minute})
\DeclareSIUnit\molar{M} % concentration molaire (ex. \SI{3}{\molar}, \SI{1}{\micro\molar})
\DeclareSIUnit\an{an} % année (le modèle confond parfois avec \year, un compteur TeX) — ex. \SI{5}{\kilo\meter\per\an}
\DeclareSIUnit\FCFA{FCFA} % franc CFA (ex. \SI{500}{\FCFA\per\kilo\gram})
\DeclareSIUnit\degreeBrix{\textdegree Brix} % teneur en sucre d'un sirop/jus (réfractométrie)
\DeclareSIUnit\month{mois} % le modèle utilise parfois \month, un compteur TeX, comme unité de durée
\DeclareSIUnit\microliter{\micro\liter} % le modèle écrit parfois \microliter en un seul token
\DeclareSIUnit\million{millions} % dénombrement (ex. \SI{15}{\million\per\mL} spermatozoïdes)
\usepackage{qrcode}
% \degree : commande fréquemment utilisée par les modèles pour un angle ou une
% température en dehors de \SI{}{} (non fournie par les paquets chargés).
\providecommand{\degree}{^{\circ}}
% \qed (fin de démonstration), utilisé par les modèles sans amsthm
\providecommand{\qed}{\hfill$\square$}
% \barre : double barre des valeurs interdites dans un tableau de variations (tkz-tab)
\providecommand{\barre}{\ensuremath{\|}}
% environnement proof (amsthm non chargé) utilisé par les modèles
\ifdefined\proof\else\newenvironment{proof}[1][Démonstration]{\par\noindent\textit{#1.}\ }{\qed\par}\fi
\usepackage{lastpage}
\usepackage{hyperref}
\hypersetup{hidelinks}

% ============================================================
% Palette « Pop » OnBuch+ — mêmes hex que la classe P (plus_ui.dart)
% ============================================================
\definecolor{poporange}{HTML}{F59321}
\definecolor{popdark}{HTML}{E07A0C}
\definecolor{popcream}{HTML}{FFF6E8}
\definecolor{popink}{HTML}{1C1714}
\definecolor{poppurple}{HTML}{7A5AE0}
\definecolor{popgreen}{HTML}{2E9E6A}
\definecolor{poppink}{HTML}{E0457B}
\definecolor{popblue}{HTML}{2F7DE1}
\definecolor{popgold}{HTML}{C98A12}
\definecolor{popmuted}{HTML}{8A7F76}
\definecolor{popline}{HTML}{EADFD2}
\definecolor{popgray}{HTML}{8A7F76}
\colorlet{popgrey}{popgray}
% Alias tolérants (le modèle nomme parfois les couleurs différemment)
\colorlet{popwhite}{white}
\colorlet{popblack}{popink}
\colorlet{popred}{poppink}
\colorlet{popyellow}{popgold}
% Teintes claires (fonds de tableaux/encadrés) — le modèle nomme parfois
% les couleurs avec un suffixe L (ex. popblueL) sans les avoir définies.
\colorlet{poporangeL}{poporange!20}
\colorlet{popdarkL}{popdark!20}
\colorlet{poppurpleL}{poppurple!20}
\colorlet{popgreenL}{popgreen!20}
\colorlet{poppinkL}{poppink!20}
\colorlet{popblueL}{popblue!20}
\colorlet{popgoldL}{popgold!20}
\colorlet{popredL}{popred!20}
\colorlet{popgrayL}{popgray!20}
% Teintes claires pour fonds de tableaux / zones
\colorlet{popblueL}{popblue!10!white}
\colorlet{poporangeL}{poporange!14!white}
\colorlet{popgreenL}{popgreen!12!white}
\colorlet{poppurpleL}{poppurple!12!white}
\colorlet{poppinkL}{poppink!12!white}
\colorlet{popgoldL}{popgold!14!white}

\pagecolor{popcream}

% ============================================================
% Typographie
% ============================================================
\defaultfontfeatures{Ligatures=TeX}
\setmainfont{PlusJakartaSans-Medium.ttf}[Path=../../../fonts/,BoldFont=PlusJakartaSans-Bold.ttf]
% Maths en sans-serif (Fira Math) avec chiffres et lettres droites de Plus
% Jakarta, pour que les formules se fondent dans le texte.
\usepackage[math-style=ISO,bold-style=ISO]{unicode-math}
\setmathfont{FiraMath-Regular.otf}[Path=../../../fonts/]
\setmathfont{PlusJakartaSans-Medium.ttf}[Path=../../../fonts/,range={up/num,up/{Latin,latin},\mathup}]
\usepackage{icomma} % virgule décimale sans espace en mode math
% Caractères mathématiques Unicode absents des polices de texte (Plus Jakarta,
% Archivo Black) : rendus par leur équivalent mathématique, en texte comme en
% formule — sinon ils s'affichent en carré vide (ex. « Arithmétique dans ℤ »).
\usepackage{newunicodechar}
\newunicodechar{ℝ}{\ensuremath{\mathbb{R}}}
\newunicodechar{ℤ}{\ensuremath{\mathbb{Z}}}
\newunicodechar{ℂ}{\ensuremath{\mathbb{C}}}
\newunicodechar{ℕ}{\ensuremath{\mathbb{N}}}
\newunicodechar{ℚ}{\ensuremath{\mathbb{Q}}}
\newunicodechar{𝔻}{\ensuremath{\mathbb{D}}}
\newunicodechar{α}{\ensuremath{\alpha}}
\newunicodechar{β}{\ensuremath{\beta}}
\newunicodechar{γ}{\ensuremath{\gamma}}
\newunicodechar{δ}{\ensuremath{\delta}}
\newunicodechar{ε}{\ensuremath{\varepsilon}}
\newunicodechar{θ}{\ensuremath{\theta}}
\newunicodechar{λ}{\ensuremath{\lambda}}
\newunicodechar{μ}{\ensuremath{\mu}}
\newunicodechar{σ}{\ensuremath{\sigma}}
\newunicodechar{τ}{\ensuremath{\tau}}
\newunicodechar{φ}{\ensuremath{\varphi}}
\newunicodechar{ψ}{\ensuremath{\psi}}
\newunicodechar{ω}{\ensuremath{\omega}}
\newunicodechar{Σ}{\ensuremath{\Sigma}}
% majuscules produites par \MakeUppercase dans les titres (α-aminés -> Α)
\newunicodechar{Α}{\ensuremath{\alpha}}
\newunicodechar{Β}{\ensuremath{\beta}}
\newunicodechar{Γ}{\ensuremath{\Gamma}}
\newunicodechar{Δ}{\ensuremath{\Delta}}
\newunicodechar{Ε}{\ensuremath{\varepsilon}}
\newunicodechar{Θ}{\ensuremath{\Theta}}
\newunicodechar{Λ}{\ensuremath{\Lambda}}
\newunicodechar{Μ}{\ensuremath{\mu}}
\newunicodechar{Τ}{\ensuremath{\tau}}
\newunicodechar{Φ}{\ensuremath{\Phi}}
\newunicodechar{Ψ}{\ensuremath{\Psi}}
\newunicodechar{Ω}{\ensuremath{\Omega}}
\newunicodechar{↦}{\ensuremath{\mapsto}}
\newunicodechar{∈}{\ensuremath{\in}}
\newunicodechar{∉}{\ensuremath{\notin}}
\newunicodechar{∧}{\ensuremath{\wedge}}
\newunicodechar{⇔}{\ensuremath{\Leftrightarrow}}
\newunicodechar{⟺}{\ensuremath{\Longleftrightarrow}}
\newunicodechar{⇒}{\ensuremath{\Rightarrow}}
\newunicodechar{⟹}{\ensuremath{\Longrightarrow}}
\newunicodechar{∘}{\ensuremath{\circ}}
\newunicodechar{∪}{\ensuremath{\cup}}
\newunicodechar{∩}{\ensuremath{\cap}}
\newunicodechar{≡}{\ensuremath{\equiv}}
\newunicodechar{⊂}{\ensuremath{\subset}}
\newunicodechar{⊥}{\ensuremath{\perp}}
\newunicodechar{∃}{\ensuremath{\exists}}
\newunicodechar{∀}{\ensuremath{\forall}}
\newunicodechar{∛}{\ensuremath{\sqrt[3]{\,}}}
\newunicodechar{∜}{\ensuremath{\sqrt[4]{\,}}}
\newunicodechar{⃗}{\ensuremath{{}^{\to}}}
\newunicodechar{ⁿ}{\ifmmode^{n}\else\textsuperscript{\ensuremath{n}}\fi}
\newunicodechar{ˣ}{\ifmmode^{x}\else\textsuperscript{\ensuremath{x}}\fi}
\newunicodechar{ᵇ}{\ifmmode^{b}\else\textsuperscript{\ensuremath{b}}\fi}
\newunicodechar{ʳ}{\ifmmode^{r}\else\textsuperscript{\ensuremath{r}}\fi}
\newunicodechar{ᵘ}{\ifmmode^{u}\else\textsuperscript{\ensuremath{u}}\fi}
\newunicodechar{ʸ}{\ifmmode^{y}\else\textsuperscript{\ensuremath{y}}\fi}
\newunicodechar{ᵃ}{\ifmmode^{a}\else\textsuperscript{\ensuremath{a}}\fi}
\newunicodechar{ᵉ}{\ifmmode^{e}\else\textsuperscript{\ensuremath{e}}\fi}
\newunicodechar{ᵏ}{\ifmmode^{k}\else\textsuperscript{\ensuremath{k}}\fi}
\newunicodechar{ᵖ}{\ifmmode^{p}\else\textsuperscript{\ensuremath{p}}\fi}
\newunicodechar{ᴺ}{\ifmmode^{N}\else\textsuperscript{\ensuremath{N}}\fi}
\newunicodechar{ᶜ}{\ifmmode^{c}\else\textsuperscript{\ensuremath{c}}\fi}
\newunicodechar{ʰ}{\ifmmode^{h}\else\textsuperscript{\ensuremath{h}}\fi}
\newunicodechar{ᵠ}{\ifmmode^{q}\else\textsuperscript{\ensuremath{q}}\fi}
\newunicodechar{ₙ}{\ifmmode_{n}\else\textsubscript{\ensuremath{n}}\fi}
\newunicodechar{ₐ}{\ifmmode_{a}\else\textsubscript{\ensuremath{a}}\fi}
\newunicodechar{ᵢ}{\ifmmode_{i}\else\textsubscript{\ensuremath{i}}\fi}
\newunicodechar{ᵣ}{\ifmmode_{r}\else\textsubscript{\ensuremath{r}}\fi}
\newunicodechar{ₖ}{\ifmmode_{k}\else\textsubscript{\ensuremath{k}}\fi}
\newunicodechar{ᵦ}{\ifmmode_{\beta}\else\textsubscript{\ensuremath{\beta}}\fi}
\newunicodechar{ₓ}{\ifmmode_{x}\else\textsubscript{\ensuremath{x}}\fi}
\newfontfamily\popdisplay{ArchivoBlack-Regular.ttf}[Path=../../../fonts/]
\newfontfamily\poplogo{SpaceGrotesk-Bold.ttf}[Path=../../../fonts/]
\newfontfamily\popsemi{PlusJakartaSans-SemiBold.ttf}[Path=../../../fonts/]
\newfontfamily\popxbold{PlusJakartaSans-ExtraBold.ttf}[Path=../../../fonts/]
\newfontfamily\popxboldls{PlusJakartaSans-ExtraBold.ttf}[Path=../../../fonts/,LetterSpace=3.0]

\setlist[enumerate]{leftmargin=1.7em,itemsep=5pt,topsep=4pt}
\setlist[itemize]{leftmargin=1.7em,itemsep=4pt,topsep=4pt,label={\color{poporange}\textbullet}}
\setlength{\parindent}{0pt}
\setlength{\parskip}{5pt}
\renewcommand{\arraystretch}{1.25}

% pgfplots : style Pop par défaut
\pgfplotsset{
  every axis/.append style={axis line style={popink,line width=1pt},tick style={popink},
    grid style={popline,line width=0.5pt},label style={font=\small},tick label style={font=\footnotesize}},
  popcurve/.style={poporange,line width=1.8pt,no markers,smooth},
}
% Même style disponible dans un \draw TikZ simple (hors pgfplots)
\tikzset{popcurve/.style={poporange,line width=1.8pt}}
\tikzset{popgrid/.style={popline,very thin}}
\colorlet{popcurve}{poporange} % le nom du style est parfois utilisé comme couleur

% ============================================================
% Pill / squiggle
% ============================================================
\newcommand{\poppill}[3]{%
  \tikz[baseline=-0.55ex]\node[draw=popink,line width=1.2pt,rounded corners=2.6pt,%
    fill=#2,inner xsep=6.5pt,inner ysep=3pt]{%
    \popxboldls\fontsize{7}{7}\selectfont\color{#3}\MakeUppercase{#1}};%
}
\newcommand{\popsquiggle}[1]{%
  \tikz[baseline=-0.4ex]{
    \draw[#1,line width=2.6pt,line cap=round]
      plot [smooth] coordinates {(0,0.09) (0.34,0.01) (0.68,0.21) (1.02,0.01) (1.36,0.10)};
  }%
}
% Mot-clé surligné (marqueur orange sous le texte)
\newcommand{\cle}[1]{{\popxbold\color{popdark}#1}}

% ============================================================
% En-tête / pied
% ============================================================
\pagestyle{fancy}
\fancyhf{}
\renewcommand{\headrulewidth}{0pt}
\renewcommand{\footrulewidth}{0pt}
\lhead{\raisebox{-0.15ex}{\poplogo\fontsize{13}{13}\selectfont\color{popink}OnBuch\color{poporange}+}}
\rhead{\poppill{\DOCMATIERE\ \textbullet\ \DOCNIVEAU\ \textbullet\ Leçon \DOCLECON}{white}{popink}}
\lfoot{\popxbold\fontsize{7.6}{7.6}\selectfont\color{popmuted}\MakeUppercase{Cours signé OnBuch+}\ \textbullet\ {\popsemi\DOCTITRE}}
\rfoot{\tikz[baseline=-0.6ex]\node[rounded corners=8.5pt,fill=popink,minimum height=17pt,minimum width=17pt,inner xsep=5pt,inner sep=0]{\popxbold\fontsize{8}{8}\selectfont\color{popcream}\thepage\,/\,\pageref*{LastPage}};}

% ============================================================
% Titres
% ============================================================
\newcommand{\chaptitre}[2]{%
  \begin{tcolorbox}[enhanced,colback=poporange,colframe=popink,boxrule=2.6pt,arc=11pt,
      drop shadow=popink,left=15pt,right=15pt,top=12pt,bottom=11pt,before skip=3pt,after skip=18pt]
    \poppill{#2}{popink}{popcream}\par\vspace{9pt}
    {\raggedright\hyphenpenalty=10000\exhyphenpenalty=10000\popdisplay\fontsize{22}{24}\selectfont\color{popcream}\MakeUppercase{#1}\par}
  \end{tcolorbox}
}
% Section numérotée : pastille orange avec le numéro + titre Archivo
\newcounter{popsec}
\newcommand{\coursec}[1]{%
  \stepcounter{popsec}\setcounter{popsub}{0}%
  \par\vspace{16pt}\noindent
  \begin{minipage}{\linewidth}
  \tikz[baseline=-0.7ex]\node[draw=popink,line width=1.6pt,fill=poporange,rounded corners=4pt,minimum size=22pt,inner sep=0]{\popdisplay\fontsize{12}{12}\selectfont\color{popcream}\thepopsec};%
  \hspace{8pt}{\popdisplay\fontsize{15}{17}\selectfont\color{popink}\MakeUppercase{#1}}\par
  \vspace{4pt}\noindent{\color{poporange}\rule{36pt}{3.6pt}}
  \end{minipage}\par\nopagebreak\vspace{8pt}
}
\newcounter{popsub}
\newcommand{\courssub}[1]{%
  \stepcounter{popsub}%
  \par\vspace{9pt}\noindent{\popxbold\fontsize{11.5}{13}\selectfont\color{popink}{\color{poporange}\thepopsec.\thepopsub}\hspace{6pt}#1}\par\nopagebreak\vspace{4pt}
}

% ============================================================
% Boîtes pédagogiques — même recette : fond blanc, bordure épaisse colorée,
% ombre dure encre, badge plein. Titre optionnel [#1].
% ============================================================
\tcbset{popbase/.style={breakable,enhanced,colback=white,boxrule=2.3pt,arc=9pt,
  shadow={3pt}{-3pt}{0pt}{popink},left=14pt,right=14pt,top=11pt,bottom=11pt,before skip=11pt,after skip=15pt,
  fontupper=\popsemi\fontsize{10.6}{14.5}\selectfont\color{popink},
  fontlower=\popsemi\fontsize{10.6}{14.5}\selectfont\color{popink}}}
% Coche dessinée (le glyphe ✓ n'existe pas dans les polices du document)
\newcommand{\popcheck}[1]{\tikz[baseline=-0.5ex]\draw[#1,line width=1.6pt,line cap=round,line join=round](-0.13,0.0)--(-0.04,-0.09)--(0.13,0.1);}
\providecommand{\coche}{\popcheck{popgreen}}
\providecommand{\resultat}[1]{\par\smallskip\noindent\fcolorbox{popgreen}{popgreenL}{\parbox{\dimexpr\linewidth-2\fboxsep-2\fboxrule\relax}{\textbf{Résultat.} #1}}\par\smallskip}
\newcommand{\popboxhead}[3]{\poppill{#1}{#2}{white}\ifx\relax#3\relax\else\hspace{6pt}{\popxbold\fontsize{10.6}{12}\selectfont\color{popink}#3}\fi\par\vspace{7pt}}

\newtcolorbox{aretenir}[1][]{popbase,colframe=popgreen,before upper={\popboxhead{À retenir}{popgreen}{#1}}}
\newtcolorbox{exemplebox}[1][]{popbase,colframe=poporange,before upper={\popboxhead{Exemple}{poporange}{#1}}}
\newtcolorbox{definition}[1][]{popbase,colframe=popblue,before upper={\popboxhead{Définition}{popblue}{#1}}}
\newtcolorbox{propriete}[1][]{popbase,colframe=poppurple,before upper={\popboxhead{Propriété}{poppurple}{#1}}}
\newtcolorbox{methode}[1][]{popbase,colframe=popgold,before upper={\popboxhead{Méthode}{popgold}{#1}}}
\newtcolorbox{attention}[1][]{popbase,colframe=poppink,before upper={\popboxhead{Attention}{poppink}{#1}}}
\newtcolorbox{experience}[1][]{popbase,colframe=popblue,colback=popblueL,before upper={\popboxhead{Expérience}{popblue}{#1}}}
% « Le savais-tu ? » : carte violette pleine (contexte, culture, Cameroun)
\newtcolorbox{savaistu}[1][]{popbase,colback=poppurple,colframe=popink,
  fontupper=\popsemi\fontsize{10.4}{14.2}\selectfont\color{white},
  before upper={\poppill{Le savais-tu\,?}{white}{poppurple}\ifx\relax#1\relax\else\hspace{6pt}{\popxbold\color{white}#1}\fi\par\vspace{7pt}}}
% Exercice résolu : énoncé en haut, \tcblower puis la solution sur fond crème
\newcounter{popexo}\newcounter{popexr}
\newtcolorbox{exoresolu}[1][]{popbase,colframe=poporange,colbacklower=popcream,
  segmentation style={popink,line width=1.2pt,dashed},
  before upper={\stepcounter{popexr}\popboxhead{Exercice résolu \thepopexr}{poporange}{#1}},
  before lower={\poppill{Solution}{popink}{popcream}\par\vspace{6pt}}}
% Exercice (non résolu en place) — la correction est dans la partie Corrigés
\newtcolorbox{exercice}[1][]{popbase,colframe=popink,
  before upper={\stepcounter{popexo}\popboxhead{Exercice \thepopexo}{popink}{#1}}}
% Corrigé
\newtcolorbox{corrige}[1][]{popbase,colframe=popgreen,colback=popgreenL,
  before upper={\popboxhead{Corrigé}{popgreen}{#1}}}
% Objectifs / prérequis / bilan
\newtcolorbox{objectifs}{popbase,colframe=popink,colback=poporangeL,
  before upper={\popboxhead{Objectifs de la leçon}{popink}{}}}
\newtcolorbox{prerequis}{popbase,colframe=popink,colback=popblueL,
  before upper={\popboxhead{Prérequis}{popblue}{}}}
\newtcolorbox{bilan}{popbase,colframe=popink,colback=white,boxrule=2.8pt,
  before upper={\popboxhead{Fiche bilan}{poporange}{L'essentiel à connaître pour l'examen}}}
% Figure encadrée avec légende
\newtcolorbox{popfigure}[1][]{enhanced,breakable=false,colback=white,colframe=popline,boxrule=1.4pt,arc=8pt,
  left=10pt,right=10pt,top=10pt,bottom=8pt,before skip=10pt,after skip=12pt,halign=center,
  lower separated=false,halign lower=center,
  fontlower=\popsemi\fontsize{9}{11.5}\selectfont\color{popmuted},
  #1}
\newcommand{\legende}[1]{\par\vspace{6pt}{\popsemi\fontsize{9}{11.5}\selectfont\color{popmuted}#1}}

% ============================================================
% Signature OnBuch+
% ============================================================
\newcommand{\poplogoinline}[1]{{\poplogo\fontsize{#1}{#1}\selectfont\color{popink}OnBuch\color{poporange}+}}
\newcommand{\signatureonbuch}{%
  \begin{tcolorbox}[enhanced,colback=popink,colframe=popink,boxrule=2.2pt,arc=10pt,
      drop shadow=poporange,left=14pt,right=14pt,top=10pt,bottom=10pt,before skip=0pt,after skip=0pt]
    \tikz[baseline=-0.7ex]\node[circle,fill=popgreen,draw=popcream,line width=1.3pt,minimum size=22pt,inner sep=0]{\popcheck{white}};%
    \hspace{9pt}\begin{minipage}[c]{0.62\linewidth}
      {\popxbold\fontsize{12}{13}\selectfont\color{popcream}Signé\ \textbullet\ {\poplogo OnBuch\color{poporange}+}}\par\vspace{2pt}
      {\raggedright\popsemi\fontsize{8}{10.5}\selectfont\color{popline}\MakeUppercase{Équipe pédagogique OnBuch+}\\\MakeUppercase{Conforme au programme officiel MINESEC}\par}
    \end{minipage}\hfill
    {\poplogo\fontsize{22}{22}\selectfont\color{popcream}OnBuch\color{poporange}+}
  \end{tcolorbox}%
}

% ============================================================
% Page de garde
% #1 titre · #2 matière · #3 niveau · #4 module · #5 n° leçon · #6 durée ·
% #7 description / objectifs (texte ou itemize)
% ============================================================
\newcommand{\pagegarde}[7]{%
  \thispagestyle{empty}
  \newgeometry{a4paper,margin=1.7cm,top=1.6cm,bottom=1.4cm}%
  \begin{tikzpicture}[remember picture,overlay]
    \fill[poporange!18] ([xshift=-3.2cm,yshift=-3.6cm]current page.north east) circle (4.6cm);
    \fill[poppurple!14] ([xshift=2.4cm,yshift=7.6cm]current page.south west) circle (3.8cm);
  \end{tikzpicture}%
  {\poplogo\fontsize{30}{30}\selectfont\color{popink}OnBuch\color{poporange}+}\hfill
  \raisebox{0.6ex}{\poppill{Cours signé \textbullet\ Premium}{popink}{popcream}}\par
  \vspace{1pt}\popsquiggle{poppurple}\par
  \vspace{8pt}
  \begin{tcolorbox}[enhanced,colback=poporange,colframe=popink,boxrule=2.8pt,arc=13pt,
      drop shadow=popink,left=18pt,right=18pt,top=12pt,bottom=13pt,after skip=10pt]
    \poppill{#2 \textbullet\ #3}{popink}{popcream}\hspace{5pt}\poppill{#4}{white}{popink}\par\vspace{8pt}
    {\popdisplay\fontsize{14}{14}\selectfont\color{popink}LEÇON #5}\par\vspace{4pt}
    {\raggedright\hyphenpenalty=10000\exhyphenpenalty=10000\popdisplay\fontsize{25}{27}\selectfont\color{popcream}\MakeUppercase{#1}\par}
    \vspace{8pt}
    {\popxbold\fontsize{9.5}{9.5}\selectfont\color{popink}\MakeUppercase{Durée conseillée : #6}}
  \end{tcolorbox}
  \begin{tcolorbox}[enhanced,colback=white,colframe=popink,boxrule=2.2pt,arc=10pt,
      drop shadow=popink,left=16pt,right=16pt,top=10pt,bottom=10pt,after skip=8pt]
    \poppill{Dans cette leçon}{poppurple}{white}\par\vspace{5pt}
    \popsemi\fontsize{9.3}{12.6}\selectfont\color{popink}#7
  \end{tcolorbox}
  \noindent\begin{minipage}[t]{0.487\linewidth}
    \begin{tcolorbox}[enhanced,colback=popgreen,colframe=popink,boxrule=2.2pt,arc=10pt,
        drop shadow=popink,left=11pt,right=11pt,top=10pt,bottom=10pt,height=3.1cm,valign=center]
      \tcbox[colback=white,colframe=popink,boxrule=1.3pt,arc=5pt,boxsep=0pt,nobeforeafter,
        left=3pt,right=3pt,top=3pt,bottom=3pt,valign=center]{\qrcode[height=1.9cm]{https://play.google.com/store/apps/details?id=com.luvvix.onbuch&hl=fr}}%
      \hspace{8pt}\begin{minipage}[c]{0.5\linewidth}
        {\popdisplay\fontsize{11}{12.5}\selectfont\color{white}\MakeUppercase{Télécharge OnBuch}}\par\vspace{4pt}
        {\raggedright\popsemi\fontsize{8.4}{11}\selectfont\color{white}Gratuit \textbullet\ Google Play\\Tuteur IA, annales, concours, résultats\par}
      \end{minipage}
    \end{tcolorbox}
  \end{minipage}\hfill
  \begin{minipage}[t]{0.487\linewidth}
    \begin{tcolorbox}[enhanced,colback=poppurple,colframe=popink,boxrule=2.2pt,arc=10pt,
        drop shadow=popink,left=11pt,right=11pt,top=10pt,bottom=10pt,height=3.1cm,valign=center]
      \tikz[baseline=-0.7ex]\node[circle,fill=white,draw=popink,line width=1.3pt,minimum size=1.6cm,inner sep=0]{\popdisplay\fontsize{24}{24}\selectfont\color{poppurple}+};%
      \hspace{8pt}\begin{minipage}[c]{0.6\linewidth}
        {\popdisplay\fontsize{11}{12.5}\selectfont\color{white}\MakeUppercase{Passe à OnBuch+}}\par\vspace{4pt}
        {\raggedright\popsemi\fontsize{8.4}{11}\selectfont\color{white}Tous les cours signés, Tuteur IA illimité, fiches PDF — onglet OnBuch+ de l'app\par}
      \end{minipage}
    \end{tcolorbox}
  \end{minipage}
  \vspace{8pt}
  \popcontenu
  \vfill
  \signatureonbuch
  \restoregeometry
  \newpage
}

% Rangée « Ce cours contient » (page de garde)
\newcommand{\poptile}[3]{% #1 couleur #2 gros texte #3 légende
  \begin{tcolorbox}[enhanced,colback=#1,colframe=popink,boxrule=2pt,arc=9pt,shadow={2.5pt}{-2.5pt}{0pt}{popink},
      left=6pt,right=6pt,top=9pt,bottom=9pt,halign=center,valign=center,height=1.75cm,width=\linewidth,nobeforeafter]
    {\popdisplay\fontsize{12.5}{13}\selectfont\color{white}\MakeUppercase{#2}}\par\vspace{3pt}
    {\centering\popsemi\fontsize{7.8}{9.5}\selectfont\color{white}#3\par}
  \end{tcolorbox}}
\newcommand{\popcontenu}{%
  \par\noindent{\popxboldls\fontsize{7.5}{7.5}\selectfont\color{popmuted}CE COURS SIGNÉ CONTIENT}\par\vspace{7pt}
  \noindent\begin{minipage}[t]{0.235\linewidth}\poptile{popblue}{Cours}{complet, illustré, pas à pas}\end{minipage}\hfill
  \begin{minipage}[t]{0.235\linewidth}\poptile{poporange}{Méthodes}{et exercices résolus}\end{minipage}\hfill
  \begin{minipage}[t]{0.235\linewidth}\poptile{poppink}{Exercices}{type examen + corrigés}\end{minipage}\hfill
  \begin{minipage}[t]{0.235\linewidth}\poptile{popgreen}{Fiche bilan}{l'essentiel à retenir}\end{minipage}\par}

% Sommaire
\newtcolorbox{sommaire}{enhanced,colback=white,colframe=popink,boxrule=2.2pt,arc=10pt,
  drop shadow=popink,left=18pt,right=18pt,top=13pt,bottom=13pt,before skip=6pt,after skip=20pt,
  before upper={\poppill{Sommaire}{poppurple}{white}\par\vspace{9pt}\popsemi\fontsize{10.6}{14.5}\selectfont\color{popink}}}

% Signature de fin de cours — poussée tout en bas de la dernière page
\newcommand{\findecours}{%
  \par\vfill\nopagebreak
  \begin{center}\popsquiggle{poporange}\end{center}\vspace{4pt}
  \signatureonbuch
}
\AtBeginDocument{\def\llbracket{\mathopen{[\mkern-3.5mu[}}\def\rrbracket{\mathclose{]\mkern-3.5mu]}}} % intervalles d'entiers (glyphe ⟦ absent de la police)
% Glyphes absents de Fira Math : redéfinis à partir de glyphes disponibles
\AtBeginDocument{%
  \def\setminus{\mathbin{\backslash}}\let\smallsetminus\setminus
  \def\vdots{\mathord{\vbox{\baselineskip3.2pt\lineskiplimit0pt\kern4pt\hbox{.}\hbox{.}\hbox{.}}}}%
  \def\ddots{\mathinner{\mkern1mu\raise6.5pt\hbox{.}\mkern2mu\raise3.6pt\hbox{.}\mkern2mu\raise0.7pt\hbox{.}\mkern1mu}}%
  \def\triangle{\mathord{\scalebox{0.9}{$\symup{\Delta}$}}}%
  \def\nabla{\mathord{\raisebox{\depth}{\scalebox{1}[-1]{$\symup{\Delta}$}}}}%
  \def\checkmark{\popcheck{popdark}}%
  \def\star{\mathbin{\ast}}\def\bigstar{\mathord{\ast}}%
  \def\diamond{\mathbin{\scalebox{0.6}{$\Diamond$}}}%
  \def\bigcup{\mathop{\mathchoice{\raisebox{-0.3em}{\scalebox{1.7}{$\cup$}}}{\scalebox{1.25}{$\cup$}}{\cup}{\cup}}\displaylimits}%
  \def\bigcap{\mathop{\mathchoice{\raisebox{-0.3em}{\scalebox{1.7}{$\cap$}}}{\scalebox{1.25}{$\cap$}}{\cap}{\cap}}\displaylimits}%
  \def\lesseqqgtr{\mathrel{\lessgtr}}\def\bigoplus{\mathop{\oplus}\displaylimits}%
}
\providecommand{\ssi}{si et seulement si} % abréviation fréquente
\AtBeginDocument{\providecommand{\wideparen}{\overparen}} % arc de cercle

%% FIGFIX-BEGIN v5 — correctifs globaux des figures (docs/onbuch-generation/tools/figfix)
\usepackage{adjustbox}\usepackage{etoolbox}\tcbuselibrary{raster}
% toute figure TikZ/pgfplots plus large que la ligne est réduite à la largeur disponible
\BeforeBeginEnvironment{tikzpicture}{\begin{adjustbox}{max width=\linewidth}}
\AfterEndEnvironment{tikzpicture}{\end{adjustbox}}
% légendes pgfplots lisibles par-dessus les courbes
\pgfplotsset{every axis/.append style={legend style={fill=white,fill opacity=0.92,text opacity=1,draw=black!15,inner sep=3pt}}}
% étiquettes d'arêtes (syntaxe edge["..."]) avec fond blanc
\tikzset{every edge quotes/.append style={fill=white,inner sep=1.2pt}}
%% FIGFIX-END
\begin{document}

\pagegarde{Accidents en zones montagneuses ou liés aux mouvements de terrain}{SVTEEHB}{3ème}{SVTEEHB – classe de 3ème}{N13}{4 séances (fiche de progression harmonisée)}{Cette leçon étudie les accidents liés aux mouvements de terrain (glissements, éboulements, coulées de boue) fréquents dans les zones montagneuses du Cameroun. Elle explique leurs causes (déformations des roches, action de l'eau, activités humaines) et présente les techniques de prévention et de protection des populations.
\vspace{4pt}
\begin{multicols}{2}\small
\begin{itemize}
\item À la fin de cette leçon, je sais définir un mouvement de terrain et citer des exemples camerounais.
\item Je sais distinguer une déformation souple d'une déformation cassante des roches.
\item Je sais expliquer comment l'action mécanique et chimique de l'eau fragilise les terrains.
\item Je sais montrer le rôle des activités humaines (déboisement, constructions, cultures) dans le déclenchement des mouvements de terrain.
\item Je sais interpréter des documents (photographies, coupes de terrain, données pluviométriques) pour identifier les causes d'un glissement de terrain.
\item Je sais citer et expliquer les techniques de prévention et de protection contre les mouvements de terrain.
\end{itemize}\end{multicols}}

\begin{sommaire}
\begin{multicols}{2}
\begin{enumerate}
\item Les mouvements de terrain : définition et exemples
\item Les déformations des roches : souples et cassantes
\item Le rôle de l'eau dans les mouvements de terrain
\item Le rôle de l'action de l'homme
\item Prévention et protection contre les mouvements de terrain
\item Synthèse et démarche d'analyse d'un accident
\item Activité d'intégration
\item Méthodes et exercices résolus
\item Exercices
\item Corrigés des exercices
\item Fiche bilan
\end{enumerate}
\end{multicols}
\end{sommaire}

\begin{objectifs}
\begin{itemize}
\item À la fin de cette leçon, je sais définir un mouvement de terrain et citer des exemples camerounais.
\item Je sais distinguer une déformation souple d'une déformation cassante des roches.
\item Je sais expliquer comment l'action mécanique et chimique de l'eau fragilise les terrains.
\item Je sais montrer le rôle des activités humaines (déboisement, constructions, cultures) dans le déclenchement des mouvements de terrain.
\item Je sais interpréter des documents (photographies, coupes de terrain, données pluviométriques) pour identifier les causes d'un glissement de terrain.
\item Je sais citer et expliquer les techniques de prévention et de protection contre les mouvements de terrain.
\item Je sais rédiger et justifier une démarche ou une conclusion à partir d'une situation concrète.
\item Je sais appliquer ces notions pour proposer des mesures de protection d'un quartier menacé.
\end{itemize}
\end{objectifs}

\begin{prerequis}
\begin{itemize}
\item Notion de roche et de sol (altération des roches, vue en 4ème).
\item Le cycle de l'eau et l'érosion (4ème).
\item Notion de pente et de relief (géographie et SVT du cycle).
\item Lecture et interprétation de graphiques et de tableaux de données.
\end{itemize}
\end{prerequis}

\newpage

\coursec{Les mouvements de terrain : définition, classification et exemples camerounais}

\courssub{Situation d'accroche et problématique}

En octobre 2019, après plusieurs jours de pluies diluviennes, un pan de colline s'est détaché à Bafoussam, dans la région de l'Ouest du Cameroun. Des tonnes de terre et de roches ont dévalé la pente à grande vitesse, ensevelissant des habitations du quartier Gouache et faisant plus de 40 victimes. Ce drame n'est pas isolé : chaque saison des pluies, les quartiers de Yaoundé bâtis sur les versants (Messa, Nkolbisson, Ngoa-Ekélé) subissent des affaissements et des coulées de boue ; l'axe routier Yaoundé–Bafoussam, vital pour l'économie, est régulièrement coupé par des éboulements ; sur les flancs du mont Cameroun, les coulées de lave et les lahars menacent les populations riveraines.

\begin{center}
\textbf{Pourquoi le sol, apparemment stable, se met-il brutalement en mouvement ?}\par
\textbf{Quels sont les facteurs naturels (pente, roches, eau) et humains (construction, déforestation) qui favorisent ces accidents ?}\par
\textbf{Comment identifier les signes avant-coureurs et protéger les populations ?}
\end{center}

Cette leçon pose les bases : définition rigoureuse, classification des mouvements, exemples locaux, analyse des causes (déformations, eau, homme), prévention et démarche d'analyse d'un accident. Nous adopterons une démarche d'investigation : observation de terrain $\rightarrow$ hypothèses $\rightarrow$ analyse des causes $\rightarrow$ modélisation schématique.

\courssub{Définition et vocabulaire de base}

\begin{definition}[Mouvement de terrain]
Un \cle{mouvement de terrain} (ou mouvement de masse) est le \textbf{déplacement vers le bas de matériaux géologiques} (roches meubles ou dures, terres, débris, boue) \textbf{sous l'effet de la gravité}, le long d'une surface de rupture ou dans un chenal, sans agent de transport liquide (rivière, glace, vent) dominant.
\end{definition}

\noindent Trois éléments sont indispensables pour caractériser un mouvement de terrain :
\begin{enumerate}
    \item \textbf{Le matériau :} roches altérées (saprolite), sols latéritiques, argiles, roches dures fracturées, déchets d'exploitation.
    \item \textbf{La pente :} une déclivité (angle $\alpha$ par rapport à l'horizontale) non nulle. La composante du poids $P\sin\alpha$ est le moteur.
    \item \textbf{La gravité :} force motrice universelle. Le mouvement se produit quand la \textbf{contrainte de cisaillement} (force tangentielle au plan de glissement) dépasse la \textbf{résistance au cisaillement} du matériau (cohésion + frottement interne).
\end{enumerate}

\begin{aretenir}[Bilan des forces sur une pente]
Sur un plan incliné d'angle $\alpha$, un bloc de masse $m$ (poids $\vec{P}=m\vec{g}$) subit :
\begin{itemize}
    \item Une composante normale $P_n = P\cos\alpha$ (écrase le bloc sur le plan, favorise le frottement).
    \item Une composante tangente $P_t = P\sin\alpha$ (entraîne le bloc vers le bas).
\end{itemize}
Le bloc glisse si $P_t > F_{\text{frottement}} = \mu P_n$, soit $\tan\alpha > \mu$ (où $\mu$ est le coefficient de frottement). Pour un sol meuble, on raisonne en \textbf{contraintes} : contrainte de cisaillement $\tau$ vs résistance $\tau_f = c' + \sigma'\tan\varphi'$ (critère de Mohr-Coulomb), où $\sigma' = \sigma - u$ est la contrainte effective ($u$ = pression interstitielle).
\end{aretenir}

\courssub{Classification des principaux types de mouvements}

La classification internationale (Varnes, 1978 ; adaptée programme 3ème) distingue les mouvements selon la \textbf{cinématique} (type de déplacement) et le \textbf{matériau}. Voici les quatre types majeurs à connaître pour le BEPC.

\begin{popfigure}
\centering
\begin{tikzpicture}[scale=1.1, every node/.style={font=\small}]
    % Profil topographique initial (ligne pointillée)
    \draw[dashed, thick, popmuted] (0,3) -- (2,2.5) -- (5,1.5) -- (8,0.5) -- (10,0);
    % Profil après glissement (solide)
    \draw[very thick, popdark] (0,3) -- (2,2.5) -- (3.5,2.2) -- (5.5,1.8) -- (7.5,1.2) -- (9.5,0.5) -- (10,0);
    % Surface de rupture (courbe concave)
    \draw[very thick, poporange, decorate, decoration={snake, amplitude=1.5pt, segment length=8pt}] (1.5,2.6) .. controls (3,2.0) and (5,1.6) .. (7,1.2);
    \node[poporange, right] at (7.2,1.2) {\textbf{Surface de rupture}};
    % Masse glissée (hachurée)
    \pattern[pattern=north west lines, pattern color=popblue!50] (1.5,2.6) .. controls (3,2.0) and (5,1.6) .. (7,1.2) -- (7.5,1.2) -- (5.5,1.8) -- (3.5,2.2) -- cycle;
    \node[popblue, align=center] at (4.5,2.4) {\textbf{Masse glissée}\\(bloc cohérent)};
    % Bourrelet (pied)
    \pattern[pattern=north east lines, pattern color=popgreen!60] (7,1.2) -- (7.5,1.2) -- (9.5,0.5) -- (9,0.3) -- cycle;
    \node[popgreen!70!black, align=center] at (8.5,0.6) {\textbf{Bourrelet}\\(accumulation)};
    % Scarp (tête)
    \draw[thick, poppurple] (1.5,2.6) -- (1.5,2.9);
    \node[poppurple, left] at (1.5,2.9) {\textbf{Arrachement (scarp)}};
    % Flèche gravité
    \draw[->, very thick, poppink] (4.5,3.2) -- (4.5,2.6) node[midway, right] {\textbf{$\vec{g}$ (Gravité)}};
    % Maison en haut (risque)
    \draw[popdark, thick] (0.5,2.9) rectangle (1.2,3.4);
    \draw[popdark, thick] (0.85,3.4) -- (0.85,3.7) -- (1.15,3.4);
    \node[above, popdark] at (0.85,3.7) {Habitat};
    % Route en bas
    \draw[popline, thick, double distance=2pt] (8,0) -- (10,0);
    \node[below, popmuted] at (9,-0.3) {Route coupée};
    % Échelle
    \draw[|<->|, thick] (0,-0.8) -- (2,-0.8) node[midway, below] {$\sim$ 100 m};
    % Légende motifs
    \begin{scope}[yshift=-1.8cm, xshift=0.5cm]
        \fill[pattern=north west lines, pattern color=popblue!50] (0,0) rectangle (0.6,0.3) node[right, xshift=0.2cm] {Masse glissée};
        \fill[pattern=north east lines, pattern color=popgreen!60] (0,-0.5) rectangle (0.6,-0.2) node[right, xshift=0.2cm] {Bourrelet};
        \draw[poporange, decorate, decoration={snake, amplitude=1pt, segment length=5pt}] (0,-1) -- (0.6,-1) node[right, xshift=0.2cm] {Surface de rupture};
    \end{scope}
\end{tikzpicture}
\legende{Schéma-type d'un \textbf{glissement de terrain rotationnel} : la surface de rupture est courbe (concave vers le haut), le bloc bascule en arrière (rotation), formant un bourrelet au pied et un scarp (arrachement) à la tête. L'habitat sur la tête est menacé, la route au pied est obstruée.}
\end{popfigure}

\begin{propriete}[Les quatre types majeurs (Programme 3ème)]
\begin{center}
\begin{tabularx}{\linewidth}{|l|X|X|X|}\hline
\rowcolor{popblueL}
\textbf{Type} & \textbf{Cinématique (mouvement)} & \textbf{Matériau typique} & \textbf{Vitesse / Dangerosité} \\ \hline
\textbf{Glissement de terrain} & Translation ou rotation le long d'une surface de rupture \textbf{bien individualisée} (plan ou courbe). Le massif se déplace « en bloc ». & Argiles, marnes, roches altérées (saprolite), couches tendres intercalées dans du dur. & Lente (mm/an) à rapide (m/jour). \textbf{Très destructeur} pour les infrastructures (routes, maisons). \\ \hline
\textbf{Éboulement (Chute de blocs)} & Chute libre, rebonds, roulement de \textbf{blocs individuels} détachés d'une falaise ou d'une paroi rocheuse. Pas de surface de rupture continue. & Roches dures fracturées (granite, basalte, grès), corniches calcaires. & \textbf{Extrêmement rapide} (secondes). Danger mortel immédiat, portée loin au pied de la falaise. \\ \hline
\textbf{Coulée de boue (Coulée de débris)} & Écoulement \textbf{fluide} dans un chenal (thalweg, ravin), mélange eau + sédiments (argiles, sable, blocs). Comportement rhéologique (visqueux). & Sols latéritiques saturés, cendres volcaniques, matériaux meubles en zone de forte pluie. & Rapide à très rapide (m/s à 10 m/s). \textbf{Portée longue} (km en fond de vallée), enterre tout sur son passage. \\ \hline
\textbf{Affaissement (Tassement / Effondrement)} & Enfoncement \textbf{vertical} ou sub-vertical du sol, sans déplacement latéral majeur. Perte de volume (porosité) ou vidange de cavité. & Sols compressibles (tourbe, argiles), terrains minés (carrières), karst (grottes), remblais mal compactés. & Lent (tassement différentiel) ou brutal (effondrement de carrière, doline). Dégrade bâtiments et réseaux. \\ \hline
\end{tabularx}
\end{center}
\end{propriete}

\begin{attention}[Pièges de vocabulaire fréquents au BEPC]
\begin{itemize}
    \item \textbf{Glissement $\neq$ Coulée :} Le glissement garde une structure bloc (surface de rupture nette) ; la coulée se déforme comme un fluide (pas de surface de rupture unique, écoulement en chenal).
    \item \textbf{Éboulement $\neq$ Glissement :} L'éboulement concerne des \textbf{blocs détachés} qui chutent/roulent (paroi verticale ou sous-cavée) ; le glissement concerne une \textbf{masse continue} qui glisse sur un plan incliné.
    \item \textbf{Affaissement $\neq$ Tassement fondationnel :} L'affaissement géologique touche une zone large (karst, minage) ; le tassement foundationnel est local sous un bâtiment (mauvaise portance).
    \item \textbf{Mouvement de terrain $\neq$ Séisme :} Le séisme est une \textbf{secousse vibratoire} d'origine tectonique (libération d'énergie élastique). Le mouvement de terrain est un \textbf{déplacement gravitaire} de masse. \textit{Toutefois, un séisme peut \textbf{déclencher} des glissements/éboulements (site amplificateur).}
\end{itemize}
\end{attention}

\courssub{Exemples camerounais : réalités locales et données chiffrées}

Le Cameroun, par sa géologie (socle ancien, volcans, bassins sédimentaires), son climat (fortes pluies tropicales) et son relief (montagnes de l'Ouest, Mont Cameroun, Adamawa), est un pays à \textbf{fort aléa mouvements de terrain}.

\begin{exemplebox}[Le glissement de Bafoussam, octobre 2019 (Région de l'Ouest)]
\textbf{Contexte :} Le quartier Gouache est construit sur les versants d'une colline constituée de \textbf{granites altérés en saprolite} (arène granitique, sable grossier, faible cohésion) sur plusieurs dizaines de mètres d'épaisseur. La pente naturelle est forte ($>30^\circ$).
\textbf{Déclencheur :} Pluies exceptionnelles ($>200$ mm en 72 h, selon la météo nationale). L'infiltration massive a \textbf{saturé} la saprolite : pression interstitielle $\uparrow$, frottement effectif $\downarrow$, cohésion apparente $\downarrow$.
\textbf{Phénomène :} \textbf{Glissement rotationnel} soudain d'un volume estimé $\sim 50\,000$ m$^3$. Surface de rupture concave à $\sim 15$ m de profondeur. Vitesse estimée $> 5$ m/s (coulée de débris terminale).
\textbf{Bilan humain :} \textbf{43 morts} (bilan officiel), dizaines de blessés, 50+ maisons détruites.
\textbf{Leçon :} L'urbanisation non réglementée sur pente forte + absence de drainage + pluies extrêmes = catastrophe annoncée.
\end{exemplebox}

\begin{savaistu}[Bafoussam 2019 : le rôle de l'urbanisation anarchique]
Avant 2019, le quartier Gouache était une zone agricole. L'extension urbaine non planifiée a entraîné : (1) \textbf{Décapage de la végétation} (perte d'évapotranspiration, érosion superficielle) ; (2) \textbf{Coupes de pente} pour terrasser les maisons (sur-creusement de la tête, surcharge au pied) ; (3) \textbf{Absence de réseau d'assainissement} (eaux pluviales et usées infiltrées directement dans le sol). Le Plan d'Occupation des Sols (POS) de Bafoussam classait cette zone en \textbf{zone inconstructible (aléa fort)}. Le non-respect du POS est une cause majeure.
\end{savaistu}

\begin{exemplebox}[Yaoundé : les quartiers de pente (Messa, Nkolbisson, Ngoa-Ekélé, Mvog-Ada)]
La capitale est bâtie sur un \textbf{socle granitique et gneissique} profondément altéré (latérite, saprolite $> 20$ m). Les pentes sont douces à modérées ($10^\circ$--$25^\circ$) mais très urbanisées.
\begin{itemize}
    \item \textbf{Phénomène dominant :} \textbf{Glissements lents} (fluage, solifluxion) et \textbf{affaissements différentiels} (fissures maisons, routes bombées).
    \item \textbf{Déclencheur annuel :} Saison des pluies (sept-oct, mars-mai). Saturation de l'horizon latéritique argileux (horizon B).
    \item \textbf{Exemple 2020/2021 :} Glissement avenue Kennedy (Messa), effondrement mur de soutènement Nkolbisson, coupure route Nkolbisson–Emana.
    \item \textbf{Facteur aggravant :} \textbf{Remblais non compactés} pour élargir les parcelles, fuites réseaux d'eau potable (Camwater) et assainissement.
\end{itemize}
\end{exemplebox}

\begin{exemplebox}[Axe lourd Yaoundé–Bafoussam (RN 6) et routes de l'Ouest]
La RN 6 traverse les \textbf{montagnes de l'Ouest} (Bamileke) : roches volcaniques (basaltes, trachytes) et granites altérés, pentes très fortes ($> 35^\circ$), falaises de bord de route.
\begin{itemize}
    \item \textbf{Éboulements fréquents :} Blocs basaltiques (colonnes de désagrégation) détachés par gel/dégel (altitude $> 1500$ m) et pression racinaire. Coupures régulières (ex. pk 120, pk 145).
    \item \textbf{Glissements de versant :} Argiles volcaniques (smectites) gonflantes/retractiles. Route elle-même agit comme drain (remblai) ou barrage (tranchée).
    \item \textbf{Coût économique :} Pertes de temps, usure véhicules, coûts de déblaiement (MINTP : plusieurs milliards FCFA/an).
\end{itemize}
\end{exemplebox}

\begin{exemplebox}[Mont Cameroun (Région du Sud-Ouest) : spécificité volcanique]
Volcan actif (dernière éruption 2000). Pentes douces en bas ($5^\circ$), fortes au sommet ($30^\circ$).
\begin{itemize}
    \item \textbf{Coulées de lave :} Mouvement de terrain \textit{sui generis} (fluide chaud), détruit tout sur passage (ex. 1999, 2000 : route Limbe–Buea, plantations CDC).
    \item \textbf{Lahars (coulées de boue volcaniques) :} Mélange eau de pluie + cendres/ponces non consolidées. Déclenchés par fortes pluies \textbf{sans éruption}. Vitesse $> 20$ km/h, portée $> 20$ km. Menacent Buea, Limbe, Idenau.
    \item \textbf{Glissements de flanc :} Instabilité structurelle de l'édifice volcanique (fractures, hydrothermalisme). Risque d'effondrement de secteur (type Mont St Helens).
\end{itemize}
\end{exemplebox}

\begin{popfigure}
\centering
\begin{tikzpicture}[scale=1.3]
    % Contour Cameroun simplifié
    \draw[very thick, popdark] 
        (0,0) -- (0.5,0.2) -- (1,0.5) -- (1.5,1) -- (2,1.5) -- (2.5,2) -- (2.8,2.5) % Nord
        -- (3,3) -- (3.2,3.5) -- (3,4) -- (2.5,4.2) -- (2,4) -- (1.5,3.8) -- (1,3.5) -- (0.5,3) -- (0,2.5) % Extrême-Nord / Nord
        -- (-0.5,2) -- (-1,1.5) -- (-1.2,1) -- (-1.3,0.5) -- (-1.2,0) -- (-1, -0.3) -- (-0.5, -0.5) -- (0,0); % Ouest / Sud-Ouest
    % Zones à risque (hachurées)
    % 1. Ouest / Nord-Ouest (Bafoussam, Bamenda)
    \fill[pattern=north west lines, pattern color=poporange!60, opacity=0.7] 
        (0.5,1.5) ellipse (0.6 and 0.4);
    \node[poporange, font=\bfseries\small, align=center] at (0.5,1.5) {OUEST\\NORD-OUEST\\(Bafoussam,\\Bamenda)};
    % 2. Mont Cameroun (Sud-Ouest)
    \fill[pattern=north east lines, pattern color=poppink!60, opacity=0.7] 
        (-0.8,-0.2) circle (0.35);
    \node[poppink, font=\bfseries\small, align=center] at (-0.8,-0.2) {MONT\\CAMEROUN\\(Lahars,\\Coulées lave)};
    % 3. Yaoundé (Centre)
    \fill[pattern=crosshatch, pattern color=popblue!60, opacity=0.7] 
        (1.2,1.0) circle (0.25);
    \node[popblue, font=\bfseries\small, align=center] at (1.2,1.0) {YAOUNDÉ\\(Glissements\\lents, Affaissements)};
    % 4. Adamawa (Ngaoundéré) - aléa modéré
    \fill[pattern=dots, pattern color=popgreen!60, opacity=0.5] 
        (2.2,2.8) ellipse (0.5 and 0.3);
    \node[popgreen!70!black, font=\small, align=center] at (2.2,2.8) {ADAMAWA\\(Éboulements\\falaises)};
    % Villes repères
    \foreach \x/\y/\name/\col in {
        0.5/1.5/Bafoussam/poporange,
        -0.8/-0.2/Buea/poppink,
        1.2/1.0/Yaoundé/popblue,
        2.2/2.8/Ngaoundéré/popgreen!70!black,
        2.5/1.0/Garoua/popmuted,
        -0.5/0.5/Douala/popmuted} {
        \fill[\col] (\x,\y) circle (1.5pt);
        \node[\col, font=\tiny, anchor=west] at (\x+0.05,\y) {\name};
    }
    % Légende
    \begin{scope}[xshift=4cm, yshift=-1cm, font=\small]
        \fill[pattern=north west lines, pattern color=poporange!60] (0,0) rectangle (0.5,0.25) node[right, xshift=0.2cm] {Glissements / Éboulements (Ouest, NO)};
        \fill[pattern=north east lines, pattern color=poppink!60] (0,-0.4) rectangle (0.5,-0.15) node[right, xshift=0.2cm] {Risques volcaniques (Mont Cameroun)};
        \fill[pattern=crosshatch, pattern color=popblue!60] (0,-0.8) rectangle (0.5,-0.55) node[right, xshift=0.2cm] {Glissements lents / Affaissements (Yaoundé)};
        \fill[pattern=dots, pattern color=popgreen!60] (0,-1.2) rectangle (0.5,-0.95) node[right, xshift=0.2cm] {Éboulements falaises (Adamawa)};
    \end{scope}
    % Flèche Nord
    \draw[->, thick] (3.5,3.5) -- (3.5,4.2) node[above] {N};
\end{tikzpicture}
\legende{Carte simplifiée du Cameroun localisant les principales zones à risque de mouvements de terrain. \textbf{Orange (Ouest/Nord-Ouest)} : glissements rotationnels dans saprolite granitique, éboulements basaltiques (Bafoussam, Bamenda, axe Yaoundé–Bafoussam). \textbf{Rose (Mont Cameroun)} : lahars, coulées de lave, glissements de flanc volcanique. \textbf{Bleu (Yaoundé)} : glissements lents, affaissements différentiels sur latérite argileuse. \textbf{Vert (Adamawa)} : éboulements de falaises gréseuses/basaltiques.}
\end{popfigure}

\coursec{Les déformations des roches : souples et cassantes}

\courssub{Comportement mécanique des matériaux géologiques}

La réponse d'une roche ou d'un sol à une contrainte (poids, charge tectonique, pression interstitielle) dépend de sa \textbf{rhéologie}. On distingue deux modes de déformation fondamentaux qui conditionnent le type de mouvement de terrain.

\begin{definition}[Déformation souple (ductile) et déformation cassante (fragile)]
\begin{itemize}
    \item \textbf{Déformation souple (ductile) :} Le matériau se déforme \textbf{de manière continue et permanente} sans rupture brutale. La déformation s'accumule (fluage, solifluxion). Elle est favorisée par : température/pression élevées (profondeur), roches tendres (argiles, sels, gypses), vitesse de chargement lente, présence d'eau (lubrification).
    \item \textbf{Déformation cassante (fragile) :} Le matériau accumule de l'énergie élastique puis \textbf{romp brutalement} par fracturation (faille, diaclase, surface de rupture). Elle est favorisée par : roches dures (granite, basalte, quartzite), faible confinement (proximité surface), vitesse de chargement rapide (séisme, choc), faible teneur en eau.
\end{itemize}
\end{definition}

\begin{propriete}[Lien déformation -- Type de mouvement de terrain]
\begin{center}
\begin{tabularx}{\linewidth}{|l|X|X|}\hline
\rowcolor{popblueL}
\textbf{Type de déformation} & \textbf{Mouvement de terrain associé} & \textbf{Exemple} \\ \hline
\textbf{Souple (ductile)} & \textbf{Glissements lents}, fluage (creep), solifluxion, coulées de boue lentes, affaissements progressifs. & Versants de Yaoundé (latérite argileuse) : fissures lentes, murs déversés sur années. \\ \hline
\textbf{Cassante (fragile)} & \textbf{Glissements rapides} (rupture nette), \textbf{éboulements}, chutes de blocs, effondrements brutaux (dolines, carrières). & Bafoussam 2019 (rupture nette saprolite), Éboulements RN6 (blocs basaltiques). \\ \hline
\textbf{Transition souple $\rightarrow$ cassante} & Un mouvement lent (souple) peut \textbf{basculer} en mouvement rapide (cassante) si un déclencheur (pluie extrême, séisme) fait dépasser le pic de résistance. & Scenario catastrophe : fluage ancien + pluie exceptionnelle = glissement catastrophique. \\ \hline
\end{tabularx}
\end{center}
\end{propriete}

\begin{experience}[Modélisation en classe : Déformation souple vs cassante]
\textbf{Matériel :} Barre de pâte à modeler (argile), barre de craie ou biscuit sec, deux supports, masses (poids), règle.
\textbf{Protocole :}
\begin{enumerate}
    \item Placer chaque barre sur deux supports (portée 15 cm). Mesurer la flèche initiale $f_0$.
    \item Ajouter des masses progressives (ex. 50 g par 50 g) au centre. Noter la flèche $f$ à chaque étape.
    \item Observer le comportement jusqu'à la rupture (craie) ou la déformation plastique majeure (pâte).
\end{enumerate}
\textbf{Observations / Interprétation :}
\begin{itemize}
    \item \textbf{Pâte à modeler (argile) :} Flèche augmente \textbf{régulièrement}, pas de rupture nette $\rightarrow$ \textbf{Déformation souple}. Analogue aux glissements lents / coulées.
    \item \textbf{Craie (roche dure) :} Flèche quasi nulle puis \textbf{rupture soudaine} sans avertissement $\rightarrow$ \textbf{Déformation cassante}. Analogue aux éboulements / glissements rapides.
\end{itemize}
\textbf{Conclusion :} La nature du matériau (cohésion, frottement interne) dicte le mode de rupture et donc la dangerosité (prévisible vs soudain).
\end{experience}

\coursec{Le rôle de l'eau dans les mouvements de terrain}

\courssub{Action mécanique de l'eau : le facteur déclencheur n°1}

L'eau est le principal \textbf{facteur déclencheur} des mouvements de terrain au Cameroun (climat tropical humide). Elle agit par trois mécanismes couplés :

\begin{propriete}[Mécanismes mécaniques de l'eau]
\begin{enumerate}
    \item \textbf{Augmentation du poids (charge) :} Saturation des pores $\rightarrow$ masse volumique $\rho_{\text{sat}} > \rho_{\text{sec}}$. Le poids $P = mg$ augmente $\rightarrow$ composante tangente $P_t = P\sin\alpha$ augmente. \textit{Effet :} Moteur du glissement renforcé.
    \item \textbf{Pression interstitielle (poussée d'Archimède interne) :} L'eau dans les pores exerce une pression $u$ (pression de pore). Elle \textbf{soustraite} la contrainte effective $\sigma' = \sigma - u$ (Terzaghi). La résistance au cisaillement $\tau_f = c' + \sigma'\tan\varphi'$ \textbf{diminue}. \textit{Effet :} Lubrification de la surface de rupture.
    \item \textbf{Forces de suintement (seepage forces) :} Écoulement d'eau dans le sol $\rightarrow$ forces de frottement visqueux sur les grains. Si l'écoulement est \textbf{ascendant} (source, perchage), il \textbf{allège} le sol efficace (risque de liquéfaction / souplesse). Si \textbf{descendant} (infiltration), il \textbf{alourdit} et pousse vers le bas.
\end{enumerate}
\end{propriete}

\begin{aretenir}[Le concept de "Perchage" (cas fréquent au Cameroun)]
Sur un versant à deux couches : \textbf{Horizon A perméable} (sable, latérite meuble) sur \textbf{Horizon B imperméable} (argile, roche fraîche).
\begin{enumerate}
    \item Pluie intense $\rightarrow$ infiltration rapide dans A.
    \item Eau bloquée par B $\rightarrow$ \textbf{nappe perchée} se forme à l'interface A/B.
    \item Pression interstitielle $u$ monte $\rightarrow$ contrainte effective $\sigma'$ chute $\rightarrow$ résistance $\tau_f$ chute.
    \item L'interface A/B devient une \textbf{surface de glissement privilégiée} (plan de faiblesse lubrifié).
\end{enumerate}
\textit{Exemple :} Glissements translationnels superficiels sur latérite argileuse (Yaoundé, Bafoussam).
\end{aretenir}

\courssub{Action chimique de l'eau : préparation à long terme}

L'eau agit aussi comme agent d'\textbf{altération chimique} (météorisation), affaiblissant la roche \textit{avant} toute sollicitation mécanique.

\begin{propriete}[Altération chimique et sensibilité à l'eau]
\begin{itemize}
    \item \textbf{Hydrolyse des silicates :} Feldspaths $\rightarrow$ Argiles (kaolinite, smectite). Perte de cohésion, gonflement/retrait.
    \item \textbf{Dissolution :} Calcaire, gypse, sel $\rightarrow$ karst, cavités, affaissements brutaux (dolines).
    \item \textbf{Gonflement / Retrait des argiles (smectites) :} Variation volume $\uparrow \downarrow$ selon humidité $\rightarrow$ fissuration du sol (macropores) $\rightarrow$ infiltration plus rapide $\rightarrow$ cycle vicieux. \textit{Très présent sur roches volcaniques basaltiques (Ouest, Mont Cameroun).}
    \item \textbf{Lessivage / Élessivage :} Migration des fines (argiles, oxydes) $\rightarrow$ perte de cohésion, formation de pipes (tuyaux d'érosion interne) $\rightarrow$ affaissements, effondrements.
\end{itemize}
\end{propriete}

\begin{exemplebox}[Argiles gonflantes (smectites) sur basaltes de l'Ouest]
Les basaltes altérés donnent des argiles à \textbf{smectite} (gonflement jusqu'à 15\% volume).
\begin{itemize}
    \item \textbf{Saison sèche :} Retrait $\rightarrow$ fissures profondes (polygones de dessiccation). L'eau de la pluie suivante s'infiltre \textbf{instantanément} par ces fissures (court-circuitage).
    \item \textbf{Saison des pluies :} Gonflement $\rightarrow$ pression de gonflement $\uparrow$ (pousse murs, fissure dalles) + saturation $\rightarrow$ perte résistance.
\end{itemize}
C'est un facteur majeur des dégâts sur routes (RN6) et maisons (Bafoussam, Dschang).
\end{exemplebox}

\coursec{L'action de l'homme : facteur aggravant et déclencheur}

L'activité humaine modifie l'équilibre naturel des versants. Elle agit sur les \textbf{facteurs de prédisposition} (géologie, pente, végétation) et peut être un \textbf{facteur déclencheur} direct.

\begin{methode}[Analyse des impacts anthropiques sur la stabilité des versants]
\begin{enumerate}
    \item \textbf{Déforestation / Débroussaillage :}
    \begin{itemize}
        \item Perte \textbf{cohésion racinaire} (racines = "clous" traversant surface rupture potentielle, apport 5--20 kPa).
        \item Perte \textbf{évapotranspiration} $\rightarrow$ sol plus humide en permanence $\rightarrow$ pression interstitielle de base $\uparrow$.
        \item Érosion de surface $\rightarrow$ minceur du sol, affleurement roche altérée.
    \end{itemize}
    \item \textbf{Terrassements / Coupes et Remblais (Urbanisation, Routes) :}
    \begin{itemize}
        \item \textbf{Surcreusement en tête de versant :} Supprime le contrefort (butée) $\rightarrow$ moment stabilisateur $\downarrow$.
        \item \textbf{Surcharge en pied de versant :} Remblais, déchets, constructions $\rightarrow$ moment renversant $\uparrow$ + charge additionnelle.
        \item \textbf{Création de pentes artificielles} plus fortes que l'angle de repos naturel du matériau.
    \end{itemize}
    \item \textbf{Modification de l'hydrologie :}
    \begin{itemize}
        \item \textbf{Imperméabilisation} (toits, routes, béton) $\rightarrow$ ruissellement concentré $\rightarrow$ érosion + infiltration localisée massive.
        \item \textbf{Fuites réseaux} (eau potable, assainissement) $\rightarrow$ recharge artificielle continue de la nappe perchée. \textit{Ex. Yaoundé : fuites Camwater = facteur majeur glissements lents.}
        \item \textbf{Drainage mal conçu} : Rejet eaux pluviales sur versant sans exutoire $\rightarrow$ saturation localisée.
    \end{itemize}
    \item \textbf{Vibrations / Chocs :} Traffic poids lourds, carrières (explosifs), machines travaux publics $\rightarrow$ fatigue du matériau, tassement soudain, déclenchement éboulements.
    \item \textbf{Exploitation minière / Carrières :} Création de vides (galeries, fronts de taille) $\rightarrow$ affaissements, éboulements, modification contrainte.
\end{enumerate}
\end{methode}

\begin{savaistu}[Yaoundé : le paradoxe de l'urbanisation]
Yaoundé "ville aux sept collines" : l'urbanisation galopante (croissance 4--5\%/an) a conduit à l'occupation de \textbf{100\% des versants constructibles}, dont beaucoup classés "aléa fort" dans le Plan Directeur d'Urbanisme (PDU). Les \textbf{quartiers précaires} (non lotis, sans assainissement) sont les plus exposés. La mairie et le MINHDU tentent la \textbf{relocalisation} (ex. site Nkolbisson 2) mais l'enjeu foncier et social est immense.
\end{savaistu}

\coursec{Prévention et protection contre les mouvements de terrain}

La gestion du risque suit la chaîne : \textbf{Connaissance (Aléa) $\rightarrow$ Prévention (Enjeu/Vulnérabilité) $\rightarrow$ Protection (Parade) $\rightarrow$ Gestion de crise}.

\courssub{Prévention : agir avant l'événement (Réduction de l'aléa et de la vulnérabilité)}

\begin{propriete}[Mesures de prévention (Non structurelles et structurelles douces)]
\begin{itemize}
    \item \textbf{Cartographie de l'aléa / Zonage réglementaire :} Cartes d'aléa (faible/moyen/fort) $\rightarrow$ \textbf{Plans d'Occupation des Sols (POS)} / \textbf{Plans d'Aménagement Urbain (PAU)}. \textbf{Interdiction de construire} en zone rouge (aléa fort). \textit{Ex. POS Bafoussam, PDU Yaoundé.}
    \item \textbf{Surveillance / Monitoring :} 
    \begin{itemize}
        \item \textit{Topographie / GNSS :} Repères fixes, mesure déplacements (mm/an).
        \item \textit{Inclinomètres :} Forages profonds, mesure déformation en profondeur (surface rupture).
        \item \textit{Piézomètres :} Niveau nappe / pression interstitielle (alerte pluie).
        \item \textit{Fissuromètres / Extensomètres :} Suivi ouverture fissures bâtiments/sol.
    \end{itemize}
    \item \textbf{Gestion des eaux (Hydrologie douce) :}
    \begin{itemize}
        \item \textbf{Drainage superficiel :} Rigoles, caniveaux, fossés de garde, bassins de rétention $\rightarrow$ évacuer l'eau \textbf{avant} qu'elle n'infiltre.
        \item \textbf{Drainage profond :} Drains français (tranchées drainantes), drains horizontaux (forages drainants) $\rightarrow$ abaisser nappe perchée, réduire $u$.
        \item \textbf{Étanchéité des réseaux :} Réparation fuites eau potable/assainissement (priorité Yaoundé).
    \end{itemize}
    \item \textbf{Revégétalisation / Génie végétal :} Reboisement essences à racines profondes/pivotantes (Eucalyptus, Acacia, Vetiver \textit{Chrysopogon zizanioides}). \textbf{Bio-ingénierie} : fascines, palissades, plantations en godets sur géotextile. Renforce cohésion ($c \uparrow$), évapotranspiration ($\uparrow$), protection surface.
    \item \textbf{Information / Éducation / Savoir-être :} Sensibilisation populations (signes avant-coureurs : fissures, arbres penchés, sources nouvelles), exercices d'évacuation, respect interdictions construction.
\end{itemize}
\end{propriete}

\courssub{Protection : ouvrages de génie civil (Parades structurelles)}

Utilisés quand la prévention (déplacement population) est impossible ou en complément. Coût élevé, maintenance nécessaire.

\begin{propriete}[Principaux ouvrages de protection (à identifier sur schéma)]
\begin{center}
\begin{tabularx}{\linewidth}{|l|X|X|}\hline
\rowcolor{popblueL}
\textbf{Ouvrage} & \textbf{Principe / Fonction} & \textbf{Type de mouvement ciblé} \\ \hline
\textbf{Mur de soutènement} (poids, en console, cloué) & S'oppose à la poussée du sol (butée au pied). Nécessite \textbf{drainage arrière} (impératif). & Glissements rotationnels / translationnels, Affaissements. \\ \hline
\textbf{Parois clouées / Berlinoises} & Barre d'armature (clous) injectées dans le sol + grillage + béton projeté. Renforce la masse. & Glissements profonds, talus routiers (RN6). \\ \hline
\textbf{Drains profonds / Forages drainants} & Tubages perforés + gravier + géotextile. Abaissent la nappe $\rightarrow$ $u \downarrow \rightarrow \tau_f \uparrow$. & Tous mouvements liés à l'eau (glissements, coulées). \\ \hline
\textbf{Filets / Pare-blocs / Méridiennes} & Filets câbles acier (haute énergie) plaqués sur falaise ou barrières en pied de pente. Interceptent blocs. & \textbf{Éboulements}, chutes de blocs (RN6, carrières). \\ \hline
\textbf{Ouvrages de torrent (Seuils, Batardeaux)} & Ralentissent vitesse coulées, piègent matériaux solides, protègent aval. & \textbf{Coulées de boue / Lahars} (Mont Cameroun, vallées). \\ \hline
\textbf{Déviation / Endiguement} & Canaux de dérivation, digues pour écarter le flux des enjeux. & Coulées de lave (Mont Cameroun), lahars. \\ \hline
\end{tabularx}
\end{center}
\end{propriete}

\begin{attention}[Erreurs fatales dans la conception des ouvrages]
\begin{itemize}
    \item \textbf{Mur de soutènement SANS drainage :} L'eau s'accumule derrière $\rightarrow$ poussée hydrostatique + pression interstitielle $\rightarrow$ \textbf{renversement ou glissement du mur lui-même}. (Cause n°1 de rupture des murs au Cameroun).
    \item \textbf{Ouvrage sous-dimensionné :} Calcul sur contrainte statique seulement, sans prise en charge séisme / pluie centennale / surcharge future.
    \item \textbf{Entretien négligé :} Drains colmatés, filets déchirés, végétation invasive sur méridiennes $\rightarrow$ perte d'efficacité en 2--3 ans.
\end{itemize}
\end{attention}

\coursec{Synthèse et démarche d'analyse d'un accident}

\courssub{Méthode d'analyse d'un accident (Démarche scientifique BEPC)}

\begin{methode}[Démarche type pour analyser un accident de mouvement de terrain (Sujet BEPC)]
\begin{enumerate}
    \item \textbf{Décrire le phénomène :} Type de mouvement (glissement, éboulement, coulée, affaissement), cinématique, volume estimé, vitesse.
    \item \textbf{Caractériser le site (Facteurs de prédisposition) :}
    \begin{itemize}
        \item Topographie : pente ($\alpha$), forme (concave/convexe), hauteur.
        \item Géologie / Pédologie : nature roches/sols (argile, sable, roche dure, altération), structure (stratification, fracturation, schistosité \emph{versant} ou \emph{revers}).
        \item Hydrogéologie : niveau piézométrique, perméabilité, existence nappe perchée.
        \item Végétation / Occupation du sol : forêt, culture, habitat, nu, remblais.
    \end{itemize}
    \item \textbf{Identifier le(s) facteur(s) déclencheur(s) :} Pluie (intensité, durée, cumul), séisme, action homme (coupe, fuite, vibration), érosion pied.
    \item \textbf{Expliquer le mécanisme (Physique) :} 
    \begin{itemize}
        \item Bilan forces/contraintes : Moteur ($P\sin\alpha$, poussée eau) vs Résistance ($c', \varphi', \sigma'$).
        \item Rôle de l'eau : $u \uparrow \rightarrow \sigma' \downarrow \rightarrow \tau_f \downarrow$.
        \item Type de déformation : Souple (lent, signes précurseurs) ou Cassante (brutal).
    \end{itemize}
    \item \textbf{Bilan des conséquences :} Humain, matériel, économique, environnemental.
    \item \textbf{Proposer des mesures adaptées :} Prévention (zonage, drainage, végétalisation) et/ou Protection (mur, filet, monitoring), avec justification technique.
\end{enumerate}
\end{methode}

\courssub{Application résolue : Analyse complète d'un cas pratique}

\begin{exoresolu}[Cas : Lotissement "Espoir" sur versant argilo-sableux (région du Centre)]
\textbf{Énoncé :} Un promoteur a loti un versant de pente $25^\circ$ constitué de sables argilo-feldspathiques (arène granitique) sur 10 m d'épaisseur, reposant sur granite frais. La forêt a été rasée. Les parcelles sont terrassées (coupes/remblais). Aucun drainage n'est prévu. En août (pleine saison des pluies), des fissures en arc de cercle apparaissent en tête de versant, des arbres penchent vers l'aval, et la route en pied se soulève. Analyser la situation et proposer des solutions.

\tcblower

\textbf{Solution détaillée :}

\textbf{1. Type de mouvement :} \textbf{Glissement de terrain rotationnel} (fissures en arc = scarp, arbres penchés = rotation bloc, bourrelet route = pied). Évolution possible vers \textbf{coulée de débris} si saturation totale et rupture nette.

\textbf{2. Facteurs de prédisposition (Conditions nécessaires) :}
\begin{itemize}
    \item \textit{Géotechnique :} Arène granitique (sable + argile) = \textbf{cohésion $c'$ faible} ($\sim 5$--$10$ kPa), frottement $\varphi' \sim 25^\circ$--$30^\circ$. Perméabilité moyenne $\rightarrow$ risque perchage sur granite frais (imperméable).
    \item \textit{Topographie :} Pente $25^\circ$ = \textbf{Forte}. $P_t = P\sin25^\circ \approx 0,42 P$. Moteur gravitaire important.
    \item \textit{Structure :} Interface arène/granite frais = \textbf{surface de rupture potentielle} plane/ondulée, pentée vers l'aval (versant).
    \item \textit{Végétation :} Rasure = \textbf{Perte cohésion racinaire} ($\Delta c \approx -5$ à $-15$ kPa) + \textbf{Perte évapotranspiration} (sol plus humide).
    \item \textit{Aménagement :} \textbf{Surcreusement tête} (allègement contrefort) + \textbf{Surcharge pied} (remblais) = \textbf{Déstabilisation mécanique majeure} (diminue moment stabilisateur, augmente moment renversant).
\end{itemize}

\textbf{3. Facteur déclencheur (Condition suffisante) :}
\textbf{Pluies intenses et prolongées d'août.} Infiltration massive dans arène $\rightarrow$ saturation $\rightarrow$ \textbf{perchage eau sur granite frais} $\rightarrow$ pression interstitielle $u$ maximale à l'interface $\rightarrow$ contrainte effective $\sigma' = \sigma - u$ minimale $\rightarrow$ résistance $\tau_f = c' + \sigma'\tan\varphi'$ < contrainte cisaillement $\tau$. \textbf{Rupture cassante} sur surface préexistante (transition souple $\rightarrow$ cassante).

\textbf{4. Preuves observées (Signes avant-coureurs non écoutés) :}
\begin{itemize}
    \item Fissures en arc de cercle (traction en tête) = déformation souple initiale (fluage).
    \item Arbres "pistolet" (tronc courbe) = mouvement lent ancien (solifluxion/fluage).
    \item Route bombée/fissurée = bourrelet actif (pied en compression).
\end{itemize}

\textbf{5. Mesures d'urgence et pérennes :}
\begin{itemize}
    \item \textbf{Urgence :} Évacuation habitants zone menacée (tête + corps glissement), interdiction d'accès, monitoring (topographie, fissuromètres).
    \item \textbf{Drainage (Clé) :} \textbf{Drains profonds} (forages horizontaux ou verticaux) pour abaisser nappe perchée sur granite. \textbf{Rigoles étanches} en tête pour détourner eaux de ruissellement.
    \item \textbf{Soutènement :} \textbf{Mur de pied} (en console ou cloué) \textbf{AVEC drainage arrière} pour buter le bourrelet. Reprofilage pente (allègement tête, déchargement).
    \item \textbf{Végétalisation :} \textbf{Vetiver} en lignes de niveau sur la masse glissée stabilisée + arbres à racines pivotantes en tête.
    \item \textbf{Réglementation :} Arrêt terrassements, respect POS, relocation si aléa résiduel fort.
\end{itemize}
\end{exoresolu}

\courssub{Synthèse des connaissances pour le BEPC}

\begin{aretenir}[Essentiel à retenir pour l'examen]
\begin{itemize}
    \item \textbf{Définition :} Mouvement de terrain = déplacement gravitaire de masses rocheuses/terreuses le long d'une pente.
    \item \textbf{4 types majeurs :} \textbf{Glissement} (bloc, surface rupture nette), \textbf{Éboulement} (blocs chutent/roulent), \textbf{Coulée de boue} (écoulement fluide en chenal), \textbf{Affaissement} (enfoncement vertical).
    \item \textbf{Déformations :} \textbf{Souple} (lente, ductile, signes précurseurs) vs \textbf{Cassante} (brutale, fragile, rupture nette). Un mouvement lent peut basculer en rapide.
    \item \textbf{Rôle de l'eau (Mécanique) :} Poids $\uparrow$, Pression interstitielle $u \uparrow \rightarrow \sigma' \downarrow \rightarrow \tau_f \downarrow$ (lubrification), Forces de suintement. \textbf{Perchage} = mécanisme clé sur sols à deux couches.
    \item \textbf{Rôle de l'eau (Chimique) :} Altération (argiles, dissolution), Gonflement/Retrait (smectites), Érosion interne (pipes).
    \item \textbf{Rôle de l'homme :} Déforestation (perte racines, évapotranspiration), Terrassements (surcreusement/surcharge), Fuites réseaux, Vibrations, Imperméabilisation.
    \item \textbf{Prévention :} Zonage (POS/PAU), Surveillance (inclinomètres, piézomètres), Drainage (superficiel/profond), Revégétalisation (Vetiver), Information.
    \item \textbf{Protection :} Murs de soutènement (\textbf{avec drainage}), Parois clouées, Filets pare-blocs, Ouvrages de torrent (seuils, batardeaux).
    \item \textbf{Contexte Cameroun :} Ouest/NO (glissements saprolite, éboulements basaltes), Yaoundé (glissements lents latérite), Mont Cameroun (lahars, laves), Adamawa (éboulements falaises).
    \item \textbf{Distinction clé :} Mouvement de terrain (gravité, masse) $\neq$ Séisme (tectonique, onde), mais séisme peut déclencher mouvement de terrain.
\end{itemize}
\end{aretenir}

\begin{exercice}[Entraînement BEPC : Analyse de situation et rédaction]
\begin{enumerate}
    \item \textbf{Lecture de paysage :} On observe sur un versant concave : des arbres au tronc courbé vers l'amont ("pistolet"), une fissure en arc de cercle en haut du versant, une route en bas dont le revêtement est bombé et fissuré. 
    \begin{enumerate}
        \item Identifier le type de mouvement en cours.
        \item Préciser s'il s'agit d'une déformation souple ou cassante. Justifier.
        \item Citer deux erreurs d'aménagement visibles si des maisons sont construites juste au-dessus de la fissure.
    \end{enumerate}
    \item \textbf{Rôle de l'eau :} Expliquer, en utilisant la notion de contrainte effective ($\sigma' = \sigma - u$), pourquoi une pluie intense déclenche un glissement sur une pente constituée de sable perméable reposant sur une argile imperméable.
    \item \textbf{Action de l'homme :} Un lotissement est créé sur une pente de $20^\circ$ en latérite argileuse. Le promoteur rase la forêt, fait des coupes/remblais pour aplanir les parcelles, et ne prévoit pas de caniveaux. Lister 3 actions anthropiques déstabilisantes et expliquer leur mécanisme pour chacune.
    \item \textbf{Prévention/Protection :} Compléter le tableau :
    \begin{center}
    \begin{tabular}{|l|l|}\hline
    \textbf{Mesure} & \textbf{Type (Prévention / Protection) \& Action principale} \\ \hline
    Carte d'aléa / POS & \\ \hline
    Mur de soutènement drainé & \\ \hline
    Plantation de Vetiver & \\ \hline
    Filets pare-blocs & \\ \hline
    Inclinomètres / Piézomètres & \\ \hline
    Réparation fuites Camwater & \\ \hline
    \end{tabular}
    \end{center}
\end{enumerate}
\end{exercice}

\begin{corrige}[Exercice n°1]
\begin{enumerate}
    \item \textbf{(a) Glissement de terrain rotationnel lent} (versant concave, arbres penchés = rotation/fluage, fissure arc = scarp, route bombée = bourrelet actif). 
    \textbf{(b) Déformation \textbf{souple} (ductile).} Justification : \textit{Indices progressifs (arbres pistolet = fluage sur années, fissures lentes, route bombée lente). Pas de rupture soudaine. Vitesse mm/jour à cm/jour.}
    \textbf{(c) Erreurs :} \textit{1. Maisons en tête de glissement (zone d'arrachement, perte contrefort). 2. Absence drainage visible (eau s'infiltre). 3. Terrassement apparent (coupes/remblais non végétalisés, surcharge pied).}
    \item \textbf{Mécanisme :} 
    \begin{itemize}
        \item Pluie $\rightarrow$ infiltration rapide dans sable perméable.
        \item Eau bloquée par argile imperméable $\rightarrow$ \textbf{Nappe perchée} à l'interface.
        \item Pression interstitielle $u$ \textbf{augmente} fortement à l'interface.
        \item Contrainte effective $\sigma' = \sigma - u$ \textbf{diminue} (Terzaghi).
        \item Résistance au cisaillement $\tau_f = c' + \sigma'\tan\varphi'$ \textbf{diminue}.
        \item Contrainte de cisaillement $\tau$ (due au poids $P\sin\alpha$) devient $> \tau_f$ $\rightarrow$ \textbf{Rupture / Glissement}.
    \end{itemize}
    \item \textbf{3 actions déstabilisantes :}
    \begin{enumerate}
        \item \textbf{Déforestation :} Perte cohésion racinaire ($c \downarrow$) + Perte évapotranspiration (humidité $\uparrow \rightarrow u \uparrow$).
        \item \textbf{Coupes/remblais (Terrassement) :} Surcreusement tête (allègement butée, moment stabilisateur $\downarrow$) + Surcharge pied (remblai, moment renversant $\uparrow$).
        \item \textbf{Absence drainage (Caniveaux) :} Ruissellement concentré $\rightarrow$ infiltration massive locale $\rightarrow$ saturation rapide $\rightarrow$ $u \uparrow \rightarrow \tau_f \downarrow$.
    \end{enumerate}
    \item \textbf{Tableau complété :}
    \begin{center}
    \begin{tabular}{|l|l|}\hline
    \textbf{Mesure} & \textbf{Type \& Action} \\ \hline
    Carte d'aléa / POS & \textbf{Prévention} : Zonage, interdiction construction zone fort aléa. \\ \hline
    Mur de soutènement drainé & \textbf{Protection} : Butée mécanique au pied + drainage ($u \downarrow$). \\ \hline
    Plantation de Vetiver & \textbf{Prévention} : Renforcement racinaire ($c \uparrow$), évapotranspiration ($u \downarrow$), protection surface. \\ \hline
    Filets pare-blocs & \textbf{Protection} : Interception blocs chutants (éboulements). \\ \hline
    Inclinomètres / Piézomètres & \textbf{Prévention} : Surveillance (détection mouvement lent, niveau eau), alerte précoce. \\ \hline
    Réparation fuites Camwater & \textbf{Prévention} : Suppression recharge artificielle nappe ($u \downarrow$). \\ \hline
    \end{tabular}
    \end{center}
\end{enumerate}
\end{corrige}

\coursec{Les déformations des roches : souples et cassantes}

Dans la section précédente, nous avons découvert ce qu'est un mouvement de terrain (glissement, éboulement, affaissement) et observé que ces accidents frappent souvent les zones montagneuses du Cameroun, comme les hauts plateaux de l'Ouest ou les flancs du Mont Cameroun. Mais une question fondamentale reste posée : \emph{pourquoi certaines roches cèdent-elles et se mettent-elles en mouvement ?} Pour y répondre, il faut comprendre comment les roches réagissent aux forces qui s'exercent sur elles. C'est l'objet de cette section : les roches ne restent pas éternellement intactes ; elles se \cle{déforment}, et cette déformation peut se faire de deux manières très différentes.

\courssub{Observation de faits : deux visages des roches déformées}

Commençons par observer, comme le ferait un géologue sur le terrain.

\textbf{Observation 1 : des roches pliées.} Sur les parois de certaines falaises ou dans les tranchées creusées pour les routes (par exemple sur les axes routiers de l'Ouest camerounais), on voit parfois des couches de roches qui forment des vagues, des ondulations régulières. Les couches sont \emph{courbées} mais elles restent \emph{continues} : aucune n'est coupée, aucune n'est décalée. Ces ondulations des couches rocheuses s'appellent des \cle{plis}.

\textbf{Observation 2 : des roches brisées.} Ailleurs, sur un autre affleurement, on observe le contraire : les couches sont nettement \emph{coupées} par une fissure, et les deux blocs situés de part et d'autre de cette fissure sont \emph{décalés} l'un par rapport à l'autre. La couche qui se trouvait à gauche en haut se retrouve à droite en bas. Cette fracture avec décalage s'appelle une \cle{faille}.

Ces deux observations semblent contradictoires : comment une roche, matériau que nous percevons comme dur et rigide, peut-elle tantôt se plier, tantôt se briser comme du verre ? C'est précisément ce que nous allons expliquer.

\begin{attention}[Ne pas se fier aux apparences quotidiennes]
À notre échelle, une roche paraît incassable et indéformable. Mais à l'échelle des temps géologiques (des millions d'années) et sous l'effet de forces immenses, les roches se comportent parfois comme des matériaux malléables. Le comportement d'une roche dépend des \emph{conditions} dans lesquelles elle se trouve, pas seulement de sa nature.
\end{attention}

\courssub{Expérience de modélisation en classe}

Pour comprendre ces deux comportements, réalisons une expérience simple de modélisation, réalisable dans n'importe quelle salle de classe.

\begin{experience}[Modéliser les déformations souple et cassante]
\textbf{Matériel :} une pâte molle (pâte à modeler, ou mieux : un empilement de couches de pâtes de couleurs différentes), une craie scolaire (ou un biscuit sec), les mains de l'expérimentateur.

\textbf{Protocole :}
\begin{enumerate}
\item Déposer l'empilement de couches de pâte à modeler sur la table. Exercer une pression \emph{lente et progressive} en poussant les deux extrémités l'une vers l'autre (compression lente).
\item Observer l'évolution de la forme des couches.
\item Prendre ensuite la craie et exercer une flexion \emph{rapide} en tirant ses deux extrémités vers le haut.
\item Observer ce qui se produit.
\end{enumerate}

\textbf{Observations :}
\begin{itemize}
\item La pâte à modeler se \emph{courbe} : les couches forment des ondulations (des plis) mais restent continues. Aucune rupture n'apparaît.
\item La craie, soumise à une contrainte rapide, se \emph{rompt nettement} en deux morceaux, le long d'une surface de fracture.
\end{itemize}

\textbf{Interprétation :} un même type de contrainte (une force exercée sur un matériau) produit deux résultats différents selon la \emph{manière} dont elle est appliquée et selon l'\emph{état} du matériau :
\begin{itemize}
\item contrainte lente sur un matériau déformable $\Rightarrow$ déformation \emph{sans rupture} : c'est la \cle{déformation souple} ;
\item contrainte brutale dépassant la résistance du matériau $\Rightarrow$ déformation \emph{avec rupture} : c'est la \cle{déformation cassante}.
\end{itemize}

\textbf{Conclusion :} les plis observés dans les falaises correspondent à des roches qui ont été déformées \emph{lentement}, dans des conditions où elles pouvaient se plier ; les failles correspondent à des roches qui ont \emph{cassé} lorsque la contrainte a dépassé leur résistance.
\end{experience}

\courssub{La déformation souple : les plis}

\begin{definition}[Déformation souple]
Une \cle{déformation souple} est une déformation qui modifie la forme d'une roche \emph{sans la rompre}. Les couches rocheuses se courbent et forment des \cle{plis}, tout en restant continues.
\end{definition}

Dans quelles conditions une roche peut-elle se plier sans casser ? L'expérience de la pâte à modeler nous donne la clé : il faut que la contrainte soit appliquée \emph{lentement} et que la roche soit dans un état qui la rend déformable. Dans la nature, ces conditions sont réunies \emph{en profondeur}, à l'intérieur de la Terre :

\begin{itemize}
\item la \cle{température} y est élevée : la chaleur ramollit les roches, comme la chaleur ramollit le beurre ou le plastique ;
\item la \cle{pression} y est forte et s'exerce de tous côtés : elle empêche la roche d'éclater et l'oblige à se déformer progressivement ;
\item les contraintes s'exercent pendant des durées considérables, de l'ordre de \emph{milliers à millions d'années}.
\end{itemize}

\begin{aretenir}[Conditions de la déformation souple]
La déformation souple (plissement) se produit lorsque la roche est soumise \emph{lentement} à des contraintes, généralement \emph{en profondeur}, où la chaleur et la pression la rendent déformable. Les couches se courbent mais restent \emph{continues}.
\end{aretenir}

Les plis que l'on observe aujourd'hui en surface, dans les falaises, sont donc d'anciennes roches profondes qui ont été plissées il y a très longtemps, puis ramenées à la surface par les mouvements de la Terre et dégagées par l'érosion.

\courssub{La déformation cassante : failles, fractures et diaclases}

\begin{definition}[Déformation cassante]
Une \cle{déformation cassante} est une déformation qui \emph{rompt} la roche. Elle se produit lorsque la contrainte appliquée dépasse la \emph{résistance} de la roche. On distingue :
\begin{itemize}
\item la \cle{faille} : fracture \emph{avec décalage} des deux blocs de part et d'autre de la cassure ;
\item la \cle{diaclase} : fracture \emph{sans décalage} des blocs (simple fissure) ;
\item on regroupe failles et diaclases sous le terme général de \cle{fractures}.
\end{itemize}
\end{definition}

La déformation cassante se produit typiquement \emph{près de la surface}, où les roches sont froides et rigides : comme la craie de notre expérience, elles ne peuvent pas se plier et elles cassent brutalement dès que la contrainte dépasse leur résistance.

\begin{attention}[Piège fréquent : faille ou diaclase ?]
Ne confonds pas faille et diaclase ! Les deux sont des fractures, mais :
\begin{itemize}
\item dans une \emph{faille}, les couches sont \emph{décalées} : un repère (par exemple une couche colorée) ne se prolonge pas en face ;
\item dans une \emph{diaclase}, il y a une fissure mais \emph{aucun décalage} : les couches restent en vis-à-vis.
\end{itemize}
Au BEPC, on te demandera souvent de justifier ton identification : cite toujours le critère du décalage.
\end{attention}

\begin{popfigure}
\begin{center}
\begin{tikzpicture}[scale=0.9]
% ===== Panneau gauche : plis (déformation souple) =====
\begin{scope}
\draw[fill=popblueL,draw=none] plot[smooth,domain=0:6,samples=60] (\x,{2.2+0.6*sin(120*\x)}) -- (6,0) -- (0,0) -- cycle;
\draw[popblue,thick] plot[smooth,domain=0:6,samples=60] (\x,{2.2+0.6*sin(120*\x)});
\draw[poporange,thick] plot[smooth,domain=0:6,samples=60] (\x,{1.6+0.6*sin(120*\x)});
\draw[popgreen,thick] plot[smooth,domain=0:6,samples=60] (\x,{1.0+0.6*sin(120*\x)});
\draw[popgold,thick] plot[smooth,domain=0:6,samples=60] (\x,{0.4+0.6*sin(120*\x)});
% flèches de compression
\draw[->,popink,very thick] (-0.9,1.4) -- (-0.1,1.4);
\draw[->,popink,very thick] (6.9,1.4) -- (6.1,1.4);
\node[popink] at (3,3.3) {\textbf{A. Déformation souple}};
\node[popmuted,align=center] at (3,-0.6) {couches courbées mais \textbf{continues} : plis};
\end{scope}
% ===== Panneau droit : faille (déformation cassante) =====
\begin{scope}[xshift=9.2cm]
% bloc gauche
\draw[fill=popblueL,draw=popline] (0,0) -- (2.6,0) -- (3.4,2.8) -- (0,2.8) -- cycle;
\draw[poporange,thick] (0,0.5) -- (2.72,0.5);
\draw[popgreen,thick] (0,1.1) -- (2.85,1.1);
\draw[popgold,thick] (0,1.7) -- (2.98,1.7);
\draw[popblue,thick] (0,2.3) -- (3.11,2.3);
% bloc droit (décalé vers le bas)
\draw[fill=popblueL,draw=popline] (3.4,2.8) -- (2.6,0) -- (6,0) -- (6,2.8) -- cycle;
\draw[poporange,thick] (3.5,1.5) -- (6,1.5);
\draw[popgreen,thick] (3.63,2.1) -- (6,2.1);
\draw[popgold,thick] (2.85,0.9) -- (6,0.9);
\draw[popblue,thick] (2.72,0.3) -- (6,0.3);
% plan de faille
\draw[popink,very thick] (2.6,0) -- (3.4,2.8);
\node[popink,rotate=65] at (2.75,1.6) {\small plan de faille};
% flèches de mouvement
\draw[->,poppurple,very thick] (1.5,2.0) -- (1.5,2.7);
\draw[->,poppurple,very thick] (4.8,1.6) -- (4.8,0.9);
\node[popink] at (3,3.3) {\textbf{B. Déformation cassante}};
\node[popmuted,align=center] at (3,-0.6) {couches \textbf{rompues et décalées} : faille};
\end{scope}
\end{tikzpicture}
\end{center}
\legende{Comparaison des deux types de déformation des roches. \textbf{A} : déformation souple — sous une contrainte lente (flèches noires), les couches se courbent en plis sans se rompre. \textbf{B} : déformation cassante — la roche s'est rompue le long d'un plan de faille ; les couches colorées, autrefois continues, sont décalées (flèches violettes : mouvement relatif des deux blocs).}
\end{popfigure}

\courssub{Comment distinguer les deux types de déformation sur le terrain ?}

\begin{methode}[Identifier une déformation souple ou cassante sur une photo ou un affleurement]
Face à un document (photographie de falaise, dessin d'affleurement), procède en trois étapes :
\begin{enumerate}
\item \textbf{Repère les couches} : cherche des bandes de couleurs ou de textures différentes qui matérialisent les couches rocheuses.
\item \textbf{Suis chaque couche des yeux} d'un bout à l'autre de l'affleurement :
\begin{itemize}
\item si la couche est \emph{continue} mais \emph{courbée} (elle monte et descend comme une vague) $\Rightarrow$ il s'agit d'un \emph{pli} : déformation \textbf{souple} ;
\item si la couche est \emph{interrompue} par une fissure et \emph{décalée} de part et d'autre $\Rightarrow$ il s'agit d'une \emph{faille} : déformation \textbf{cassante}.
\end{itemize}
\item \textbf{Justifie ta réponse} en citant le critère observé : « les couches sont continues et courbées » ou « les couches sont décalées le long d'une fracture ».
\end{enumerate}
\end{methode}

\begin{attention}[Une même roche, deux comportements possibles]
Erreur classique : croire qu'une roche est « souple » ou « cassante » par nature, une fois pour toutes. En réalité, \emph{la même roche} peut se plier en profondeur (où règnent chaleur et pression élevées) et casser en surface (où elle est froide et rigide). Le comportement dépend des \emph{conditions} (température, pression, vitesse d'application de la contrainte), pas seulement de la nature de la roche. Un caramel froid casse net ; le même caramel chaud s'étire sans rompre !
\end{attention}

\courssub{Le lien avec les mouvements de terrain}

Pourquoi cette étude est-elle si importante pour comprendre les accidents en zone montagneuse ? Parce que les fractures et les failles créent dans les roches des \cle{plans de faiblesse}. Une roche fracturée n'est plus un bloc solide et cohérent : elle est découpée en morceaux séparés par des surfaces le long desquelles l'accroche est faible. Sur un versant, ces morceaux peuvent alors \emph{glisser} le long des plans de faille ou de diaclase, sous l'effet de leur propre poids : c'est le déclenchement d'un \emph{éboulement} ou d'un \emph{glissement de terrain}.

\begin{popfigure}
\begin{center}
\begin{tikzpicture}[scale=1.0]
% versant
\draw[fill=popcream,draw=popline,thick] (0,0) -- (8,0) -- (8,1.2) -- (3.2,4.6) -- (0,4.6) -- cycle;
% surface du versant
\draw[popdark,very thick] (0,4.6) -- (3.2,4.6) -- (8,1.2);
% plans de fracture dans le versant
\draw[popink,thick,dashed] (2.2,4.6) -- (4.6,1.9);
\draw[popink,thick,dashed] (3.0,4.6) -- (5.4,1.75);
\draw[popink,thick,dashed] (1.4,4.6) -- (3.9,2.15);
\node[popink,align=center,rotate=-52] at (4.35,3.6) {\small plans de fracture};
% bloc glissant
\begin{scope}[shift={(0.55,-0.35)}]
\draw[fill=poporangeL,draw=poporange,very thick] (3.0,4.6) -- (5.4,1.75) -- (6.15,2.35) -- (3.85,4.6) -- cycle;
\end{scope}
\node[poporange,align=center] at (5.6,3.6) {\small bloc en glissement};
% flèche de glissement
\draw[->,poppink,very thick] (5.3,3.15) -- (6.3,2.35);
\node[poppink] at (6.9,2.9) {\small direction du glissement};
% gravité
\draw[->,popblue,very thick] (4.9,4.1) -- (4.9,3.1);
\node[popblue] at (4.9,4.45) {\small poids};
% annotations
\node[popmuted] at (1.2,0.5) {\small roche saine};
\node[popmuted,align=center] at (1.6,3.2) {\small versant fracturé\\ (failles, diaclases)};
\end{tikzpicture}
\end{center}
\legende{Versant fracturé et glissement d'un bloc rocheux. Les fractures (traits noirs pointillés) découpent la roche du versant en blocs. Sous l'effet de son poids (flèche bleue), un bloc séparé du reste du versant glisse le long d'un plan de fracture (flèche rose) : c'est un éboulement. Les fractures héritées des déformations cassantes constituent des plans de faiblesse qui favorisent les mouvements de terrain.}
\end{popfigure}

\begin{savaistu}[Les diaclases du granite de l'Ouest camerounais]
Les hauts plateaux de l'Ouest du Cameroun (régions de Bafoussam, Dschang, Mbouda) sont constitués de roches granitiques très anciennes. Ces granites sont parcourus de nombreuses \emph{diaclases}, fractures héritées des contraintes subies il y a des millions d'années. L'eau de pluie s'infiltre dans ces fissures, altère la roche et réduit encore l'accroche entre les blocs. Résultat : après les fortes pluies, d'impressionnants blocs granitiques se détachent et dévalent les pentes, menaçant les habitations et les routes. Comprendre les fractures des roches, c'est donc comprendre pourquoi certains quartiers ou certains axes routiers sont particulièrement exposés aux éboulements.
\end{savaistu}

\courssub{D'où viennent les contraintes ? (savoir admis)}

Nous avons parlé à plusieurs reprises des « contraintes » qui déforment les roches. D'où viennent-elles ? Ces forces, dites \cle{contraintes tectoniques}, sont liées aux mouvements internes de la Terre. À ton niveau, ce résultat est \emph{admis} : nous ne cherchons pas à le démontrer. Retiens simplement que l'intérieur de la Terre est en mouvement et que ces mouvements exercent des forces capables de plisser ou de fracturer les roches sur de vastes régions. L'origine précise de ces contraintes sera étudiée dans les classes supérieures.

\begin{exoresolu}[Analyser un affleurement]
\textbf{Énoncé.} Sur une photographie de falaise prise le long d'une route de montagne, un élève observe des couches de roches qui dessinent des ondulations régulières, comme des vagues. Aucune couche n'est interrompue. Un peu plus loin, il observe une autre zone où les couches sont coupées par une fissure nette, la couche sombre située à gauche de la fissure se retrouvant 1 m plus bas à droite de la fissure.

\begin{enumerate}
\item Identifie le type de déformation observé dans chaque zone, en justifiant ta réponse.
\item Laquelle de ces deux zones présente, a priori, le plus grand risque d'éboulement ? Justifie.
\end{enumerate}

\tcblower
\textbf{Solution détaillée.}

\textbf{1.} Dans la première zone, les couches sont \emph{courbées mais continues} : il s'agit de \emph{plis}, donc d'une \textbf{déformation souple}. La roche s'est déformée sans se rompre, ce qui indique une contrainte appliquée lentement, probablement en profondeur.

Dans la seconde zone, les couches sont \emph{interrompues par une fissure et décalées} (la couche sombre se retrouve 1 m plus bas de l'autre côté) : il s'agit d'une \emph{faille}, donc d'une \textbf{déformation cassante}. La contrainte a dépassé la résistance de la roche, qui a rompu.

\textbf{2.} La zone fracturée par la faille présente le plus grand risque d'éboulement. En effet, le plan de faille constitue un \emph{plan de faiblesse} : les blocs de roche situés au-dessus ne sont plus solidaires du reste de la falaise et peuvent glisser le long de ce plan sous l'effet de leur poids, surtout si l'eau de pluie s'y infiltre. La zone plissée, en revanche, est constituée de couches continues qui restent cohérentes entre elles.
\end{exoresolu}

\begin{exercice}[À toi de jouer]
Sur un dessin d'affleurement, tu observes trois couches de roches (une couche claire entre deux couches sombres). Les trois couches forment un large arc de cercle, sans aucune interruption.
\begin{enumerate}
\item Nomme la structure observée et le type de déformation correspondant.
\item Rappelle les conditions (température, pression, vitesse de la contrainte) qui permettent ce type de déformation.
\item Cite le critère qui te permet d'affirmer qu'il ne s'agit pas d'une faille.
\end{enumerate}
\end{exercice}

\begin{corrige}[Exercice : À toi de jouer]
\textbf{1.} La structure observée est un \emph{pli} : il s'agit d'une \emph{déformation souple}.

\textbf{2.} Cette déformation se produit lorsque la roche est soumise \emph{lentement} à des contraintes, dans des conditions de \emph{température élevée} et de \emph{forte pression}, c'est-à-dire généralement en profondeur.

\textbf{3.} Les couches sont \emph{continues} : aucune n'est interrompue ni décalée de part et d'autre d'une fracture. Or une faille se reconnaît précisément au décalage des couches le long d'un plan de fracture. Ce critère étant absent, il ne s'agit pas d'une faille.
\end{corrige}

\begin{aretenir}[L'essentiel de la section]
\begin{itemize}
\item Soumises à des contraintes, les roches se \emph{déforment} de deux manières :
\item \textbf{Déformation souple} : la roche se plie \emph{sans se rompre} $\Rightarrow$ formation de \emph{plis} (couches continues et courbées). Conditions : contrainte lente, chaleur et pression élevées (en profondeur).
\item \textbf{Déformation cassante} : la roche \emph{se rompt} quand la contrainte dépasse sa résistance $\Rightarrow$ formation de \emph{fractures} : \emph{failles} (avec décalage) et \emph{diaclases} (sans décalage). Conditions : roche froide et rigide (près de la surface), contrainte brutale ou excessive.
\item Le comportement d'une roche dépend des \emph{conditions}, pas seulement de sa nature : une même roche peut se plier en profondeur et casser en surface.
\item Les fractures créent des \emph{plans de faiblesse} le long desquels les blocs rocheux peuvent glisser : c'est l'une des causes des mouvements de terrain.
\end{itemize}
\end{aretenir}

Nous savons désormais que les roches fracturées sont fragilisées et prêtes à céder. Mais une fracture seule ne suffit pas toujours à déclencher un mouvement de terrain : il faut souvent un élément déclencheur ou aggravant. Le plus important d'entre eux est l'\emph{eau}, dont l'action mécanique et chimique fera l'objet de la section suivante.

\coursec{Le rôle de l'eau dans les mouvements de terrain}

Dans la section précédente, nous avons vu que les roches et les terrains peuvent se déformer de manière \cle{souple} (plis) ou \cle{cassante} (failles, fractures). Un terrain fracturé ou incliné est un terrain \emph{fragilisé} : il ne demande qu'un élément déclencheur pour se mettre en mouvement. Or, dans la plupart des catastrophes recensées au Cameroun — glissements de Bafoussam (2019), effondrements dans les Monts Bamboutos, éboulements sur les routes de l'Ouest —, cet élément déclencheur porte un nom bien connu : \textbf{l'eau}. Comment un liquide aussi banal que l'eau de pluie peut-il faire s'effondrer des pans entiers de collines ? C'est ce que nous allons découvrir en distinguant deux types d'action : l'\cle{action mécanique} et l'\cle{action chimique} de l'eau.

\courssub{Observation et mise en évidence expérimentale}

\textbf{Situation-problème.} Au marché, tu as peut-être remarqué qu'un tas de sable sec reste stable, avec des parois raides, tandis qu'après une forte pluie, ce même tas s'affaisse et s'étale. De même, les terrassiers qui creusent une fondation constatent que les parois tiennent par temps sec mais s'écroulent pendant la saison des pluies. Formulons l'hypothèse : \emph{l'eau diminue la stabilité des terrains}. Vérifions-la par une expérience simple.

\begin{experience}[Stabilité d'un talus de sable sec et d'un talus de sable arrosé]
\textbf{Matériel :} deux bacs identiques, du sable de même nature, un arrosoir, une pelle ou une truelle.

\textbf{Protocole :}
\begin{enumerate}
\item Dans le bac 1, confectionner un talus de sable \textbf{sec} en forme de petite colline à paroi raide.
\item Dans le bac 2, confectionner un talus identique, puis l'\textbf{arroser progressivement} jusqu'à saturation.
\item Observer l'évolution des deux talus pendant quelques minutes et noter les résultats dans un tableau.
\end{enumerate}

\textbf{Observations :}
\begin{center}
\begin{tabularx}{\linewidth}{|l|X|X|}
\hline
\rowcolor{popblueL} & \textbf{Talus sec (bac 1)} & \textbf{Talus arrosé (bac 2)} \\
\hline
Aspect initial & Paroi raide et stable & Paroi raide et stable \\
\hline
Après quelques minutes & Aucun changement & Le sable s'affaisse, coule et s'étale à la base \\
\hline
\end{tabularx}
\end{center}

\textbf{Interprétation :} le sable sec garde une certaine cohésion ; le sable gorgé d'eau perd sa cohésion et s'effondre. L'eau a donc \emph{déstabilisé} le talus.

\textbf{Conclusion :} l'hypothèse est vérifiée : l'eau réduit la stabilité d'un terrain. Reste à comprendre \emph{par quels mécanismes} elle agit.
\end{experience}

\courssub{L'action mécanique de l'eau}

\textbf{L'infiltration et l'alourdissement du sol}

Quand il pleut sur un versant, une partie de l'eau ruisselle en surface, mais une autre partie \textbf{s'infiltre} dans le sol à travers les pores (petits vides entre les grains) et les fissures des roches. Cette eau accumulée a deux conséquences mécaniques :

\begin{itemize}
\item \textbf{Elle alourdit le terrain.} L'eau possède une masse : un sol gorgé d'eau est beaucoup plus lourd qu'un sol sec. Le poids supplémentaire augmente la force qui « tire » le terrain vers le bas de la pente.
\item \textbf{Elle augmente la pression dans les pores.} L'eau emprisonnée dans les vides du sol exerce une poussée qui \emph{écarte les grains les uns des autres}. Les grains se touchent moins, s'accrochent moins : la \cle{cohésion} du matériau diminue fortement.
\end{itemize}

\begin{definition}[Cohésion d'un sol]
La \cle{cohésion} d'un sol est la force qui maintient ses particules (grains, fragments de roche) liées les unes aux autres. Plus la cohésion est grande, plus le terrain est stable ; plus elle est faible, plus le terrain risque de s'effondrer ou de glisser.
\end{definition}

\textbf{L'eau, un lubrifiant des surfaces de glissement}

Sous un sol perméable (sableux ou fissuré) se trouve souvent une couche \textbf{imperméable}, par exemple une couche d'\cle{argile}. L'eau infiltrée ne peut pas la traverser : elle s'accumule au-dessus et forme une \cle{nappe perchée}. L'argile, gorgée d'eau, devient molle et glissante, comme du savon mouillé. La couche de sol située au-dessus peut alors \textbf{glisser} le long de cette surface, comme une dalle posée sur un plan incliné savonné : c'est le \cle{glissement de terrain}.

\begin{propriete}[Rôle mécanique de l'eau — admis, issu des observations]
L'eau infiltrée dans un versant :
\begin{itemize}
\item alourdit le sol et augmente la pression dans les pores, ce qui \textbf{diminue la cohésion} des matériaux ;
\item joue le rôle de \textbf{lubrifiant} le long des surfaces de glissement (argiles gorgées d'eau).
\end{itemize}
Elle favorise ainsi les glissements et les écoulements de terrain.
\end{propriete}

\begin{popfigure}
\begin{tikzpicture}[scale=1.0]
% Couche argileuse (imperméable)
\fill[popgoldL] (0,0) -- (10,0) -- (10,1.2) -- (0,2.2) -- cycle;
\draw[popdark, thick] (0,2.2) -- (10,1.2);
\node[popdark] at (8.6,0.55) {\small Couche argileuse (imperméable)};
% Sol perméable au-dessus
\fill[popgreenL] (0,2.2) -- (10,1.2) -- (10,3.4) -- (0,5.2) -- cycle;
\draw[popdark, thick] (0,5.2) -- (10,3.4);
\node[popdark] at (3.2,4.6) {\small Sol perméable (sable, roches fissurées)};
% Nappe perchée
\fill[popblue!35] (0,2.2) -- (10,1.2) -- (10,1.7) -- (0,2.8) -- cycle;
\node[popblue] at (6.9,1.75) {\small Nappe perchée};
% Surface de glissement
\draw[popink, very thick, dashed] (0,2.2) -- (10,1.2);
\node[popink] at (1.4,1.55) {\small Surface de glissement};
% Flèches de pluie
\foreach \x in {1.5,3.5,5.5,7.5} {
  \draw[-{Stealth}, popblue, thick] (\x,7.2) -- (\x,6.0);
}
\node[popblue] at (9.1,6.7) {\small Pluie};
% Flèches d'infiltration
\foreach \x in {2.5,4.5,6.5} {
  \draw[-{Stealth}, popblue, thick] (\x,5.6) -- (\x+0.15,3.6);
}
\node[popblue] at (2.6,3.2) {\small Infiltration};
% Flèche de glissement
\draw[-{Stealth}, popink, very thick] (5.2,4.4) -- (7.6,3.7);
\node[popink] at (7.9,4.15) {\small Glissement};
\end{tikzpicture}
\legende{Coupe schématique d'un versant pendant les pluies : l'eau s'infiltre dans le sol perméable, s'accumule en nappe perchée au-dessus de la couche argileuse imperméable, qui devient une surface de glissement lubrifiée.}
\end{popfigure}

\textbf{Le ruissellement et le sapement du pied des pentes}

L'eau qui ne s'infiltre pas \textbf{ruisselle} à la surface. En descendant la pente, elle arrache des particules : c'est l'\cle{érosion}. En bas du versant, l'eau (ruissellement concentré, rivière, mer) attaque et creuse le \textbf{pied de la pente} : c'est le \cle{sapement}. Privé de son soutien, le versant situé au-dessus perd son équilibre et s'effondre. C'est le même principe qu'un mur dont on creuserait la base : il finit par tomber.

\begin{exemplebox}[Les falaises et les berges]
Sur les berges de certains fleuves camerounais, l'eau creuse la base des talus pendant la saison des pluies ; des pans de terre se détachent ensuite et tombent dans l'eau. Le même phénomène de sapement explique l'effondrement progressif des falaises maritimes de la côte atlantique.
\end{exemplebox}

\courssub{L'action chimique de l'eau}

L'eau n'agit pas seulement par son poids et son pouvoir lubrifiant : elle agit aussi \textbf{chimiquement} sur les roches, c'est-à-dire qu'elle les transforme et les détruit lentement.

\textbf{La dissolution des roches : le cas du calcaire}

Certaines roches, comme le \cle{calcaire}, sont lentement \textbf{dissoutes} par l'eau de pluie (légèrement chargée en dioxyde de carbone de l'air). Au fil du temps, cette dissolution creuse des fissures, des galeries et des grottes : on obtient un paysage appelé \cle{karst}. Lorsque le toit d'une cavité souterraine devient trop mince, il s'effondre brutalement : c'est un \cle{affaissement} (ou fontis), qui peut engloutir des maisons ou des routes construites à la surface.

\textbf{L'altération des minéraux : la formation des argiles}

L'eau transforme aussi chimiquement certains minéraux des roches en \textbf{argiles} : c'est l'\cle{altération}. Or nous avons vu que l'argile, une fois gorgée d'eau, devient molle et glissante. L'action chimique de l'eau prépare donc le terrain aux glissements en fabriquant, à long terme, des couches argileuses qui serviront de surfaces de glissement. Action chimique et action mécanique se combinent ainsi dans le temps.

\begin{aretenir}[Les deux actions de l'eau]
\begin{itemize}
\item \textbf{Action mécanique :} infiltration $\rightarrow$ alourdissement du sol + pression dans les pores $\rightarrow$ perte de cohésion ; argiles gorgées d'eau $\rightarrow$ lubrification des surfaces de glissement ; ruissellement $\rightarrow$ érosion et sapement du pied des pentes.
\item \textbf{Action chimique :} dissolution des roches (calcaire $\rightarrow$ karst $\rightarrow$ affaissements) et altération des minéraux en argiles.
\end{itemize}
\end{aretenir}

\courssub{Exploitation de données : pluies et glissements de terrain}

Si l'eau déclenche les glissements, on doit observer une \textbf{corrélation} entre les fortes pluies et la fréquence des glissements. Vérifions-le sur des données (fictives mais réalistes) recueillies dans la région de l'Ouest du Cameroun en 2019.

\begin{popfigure}
\begin{tikzpicture}
\begin{axis}[
  width=0.95\linewidth, height=6.5cm,
  xlabel={Mois de l'année 2019},
  ylabel={Hauteur de pluie (mm)},
  axis y line*=left, axis x line*=bottom,
  xtick={1,2,3,4,5,6,7,8,9,10,11,12},
  xticklabels={J,F,M,A,M,J,J,A,S,O,N,D},
  ymin=0, ymax=450,
  legend style={at={(0.02,0.98)}, anchor=north west, font=\small},
]
\addplot[popcurve, color=popblue, mark=*, thick] coordinates {
(1,20) (2,35) (3,90) (4,160) (5,210) (6,180) (7,120) (8,140) (9,260) (10,410) (11,150) (12,30)
};
\legend{Pluie (mm)}
\end{axis}
\begin{axis}[
  width=0.95\linewidth, height=6.5cm,
  ylabel={Glissements recensés},
  axis y line*=right, axis x line=none,
  xtick=\empty, ymin=0, ymax=9,
  legend style={at={(0.98,0.98)}, anchor=north east, font=\small},
]
\addplot[color=popink, mark=square*, thick, dashed] coordinates {
(1,0) (2,0) (3,1) (4,1) (5,2) (6,2) (7,1) (8,1) (9,3) (10,8) (11,2) (12,0)
};
\legend{Glissements}
\end{axis}
\end{tikzpicture}
\legende{Évolution mensuelle de la hauteur de pluie (courbe bleue, axe de gauche) et du nombre de glissements de terrain recensés (courbe noire, axe de droite) dans la région de l'Ouest en 2019. Les deux pics coïncident en octobre.}
\end{popfigure}

\begin{methode}[Interpréter un graphique pluie / glissements]
\begin{enumerate}
\item Repérer le ou les \textbf{pics} de la courbe des pluies (ici : octobre, environ \SI{410}{mm}).
\item Repérer le ou les pics de la courbe des glissements (ici : octobre, 8 glissements).
\item \textbf{Comparer} les positions des pics : s'ils coïncident (ou si le pic des glissements suit de peu celui des pluies), il y a corrélation.
\item \textbf{Conclure} : les glissements sont d'autant plus fréquents que les pluies sont abondantes, ce qui confirme le rôle déclencheur de l'eau.
\end{enumerate}
\end{methode}

\begin{exemplebox}[Le glissement de Bafoussam (octobre 2019)]
Dans la nuit du 28 au 29 octobre 2019, un quartier de Bafoussam (région de l'Ouest) a été enseveli par un glissement de terrain qui a fait de nombreuses victimes. Les relevés pluviométriques montrent que plus de \SI{200}{mm} de pluie étaient tombés en quelques jours avant la catastrophe, sur un versant raide et densément habité. Le sol gorgé d'eau, alourdi et privé de cohésion, a glissé le long du versant.
\end{exemplebox}

\begin{exoresolu}[Exploiter le graphique pluie / glissements]
\textbf{Énoncé.} À l'aide du graphique ci-dessus : 1) Indiquer le mois le plus pluvieux et la hauteur de pluie correspondante. 2) Indiquer le mois où l'on recense le plus de glissements. 3) Conclure sur la relation entre pluies et glissements.
\tcblower
\textbf{Solution détaillée.}
\begin{enumerate}
\item Le pic de la courbe bleue se situe en \textbf{octobre}, avec environ \SI{410}{mm} de pluie : c'est le mois le plus pluvieux.
\item Le pic de la courbe noire se situe également en \textbf{octobre}, avec \textbf{8 glissements} recensés.
\item Les deux pics coïncident : il existe une \textbf{corrélation} entre l'abondance des pluies et la fréquence des glissements. Conformément au rôle mécanique de l'eau (perte de cohésion, lubrification des surfaces de glissement), les pluies exceptionnelles d'octobre ont déclenché davantage de mouvements de terrain.
\end{enumerate}
\end{exoresolu}

\begin{attention}[L'eau n'est pas la seule cause !]
Ne conclus jamais qu'il « suffit qu'il pleuve » pour qu'un versant s'effondre. L'eau agit le plus souvent comme un \textbf{élément déclencheur} sur un terrain \emph{déjà fragilisé} : pente raide, roches fracturées (déformations cassantes vues à la section précédente), couches argileuses, déboisement qui supprime les racines retenant le sol. Au BEPC, une bonne réponse cite toujours la fragilité préalable du terrain \emph{et} le rôle déclencheur de l'eau.
\end{attention}

\begin{savaistu}[L'eau, sculpteur et destructeur des paysages camerounais]
Dans les hautes terres de l'Ouest et du Nord-Ouest, l'altération chimique des roches volcaniques par l'eau a produit, sur des millénaires, des sols argileux très épais — fertiles pour l'agriculture, mais redoutablement glissants en saison des pluies. C'est pourquoi les mêmes régions qui offrent les meilleures terres agricoles du pays comptent aussi parmi les plus exposées aux glissements de terrain.
\end{savaistu}

En résumé, l'eau intervient à tous les étages du déclenchement d'un mouvement de terrain : elle fragilise les roches par son action chimique (dissolution, altération en argiles), elle déstabilise les versants par son action mécanique (alourdissement, pression dans les pores, lubrification, sapement), et les données pluviométriques confirment son rôle de déclencheur. Mais l'eau n'est pas le seul facteur aggravant : dans la section suivante, nous verrons que les activités humaines peuvent elles aussi préparer — ou précipiter — ces catastrophes.

\coursec{Le rôle de l'action de l'homme}

Dans la section précédente, nous avons vu que l'eau est un acteur majeur des mouvements de terrain : elle alourdit les sols, diminue la cohésion entre les particules et peut transformer un versant stable en coulée de boue. Mais l'eau n'agit pas seule. En observant les cartes des catastrophes au Cameroun, on constate un fait troublant : les glissements de terrain touchent beaucoup plus souvent les quartiers urbanisés et les collines déboisées que les forêts voisines, alors que la pente et la pluie y sont les mêmes. Cette observation pose une question simple : \cle{l'action de l'homme} peut-elle, à elle seule, déclencher ou aggraver les mouvements de terrain ? C'est ce que nous allons démontrer dans cette section, en suivant la démarche habituelle : observation, hypothèse, vérification par comparaison de documents, conclusion.

\courssub{Le déboisement : un versant privé de ses « racines protectrices »}

\textbf{Observation de départ.} Après de fortes pluies dans les hauts plateaux de l'Ouest, on remarque que les collines encore couvertes de forêt restent intactes, tandis que les collines voisines, coupées pour le bois ou les champs, présentent des cicatrices de glissement. Formulons une hypothèse : la végétation, et en particulier les racines des arbres, jouerait un rôle dans la stabilité du sol.

\begin{propriete}[Rôle stabilisateur des racines (admise, issue d'observations de terrain)]
Les \cle{racines} des arbres et des arbustes stabilisent les versants par une double action :
\begin{itemize}
\item \textbf{action mécanique} : les racines forment dans le sol un véritable « grillage » naturel qui retient les particules de terre et augmente la \cle{cohésion} du sol ;
\item \textbf{action sur l'eau} : les racines absorbent (pompent) l'eau du sol et les feuilles interceptent une partie de la pluie, ce qui limite la saturation du sol en eau.
\end{itemize}
Un versant boisé est donc plus stable qu'un versant nu de pente identique.
\end{propriete}

\begin{definition}[Déboisement]
Le \cle{déboisement} est la suppression de la couverture végétale (forêt, brousse) d'un terrain, par coupe du bois, défrichement pour l'agriculture ou incendies de brousse.
\end{definition}

Lorsque l'homme déboise un versant, il supprime le grillage de racines qui retenait le sol. Deux conséquences en découlent :
\begin{enumerate}
\item l'\cle{érosion} s'accélère : l'eau de pluie, qui n'est plus interceptée par les feuilles ni pompée par les racines, ruisselle en surface et emporte les particules de terre (on parle de \emph{ruissellement érosif}) ;
\item les \cle{glissements de terrain} deviennent plus fréquents : le sol, moins cohérent et plus chargé en eau, finit par céder sous son propre poids.
\end{enumerate}

\begin{attention}[Une forêt « lourde » est-elle dangereuse ?]
On pourrait croire que le poids des arbres surcharge le versant. C'est faux : le rôle stabilisateur des racines (cohésion + pompage de l'eau) l'emporte largement sur le poids des troncs. Toutes les observations de terrain montrent qu'un versant boisé glisse beaucoup moins qu'un versant nu.
\end{attention}

\courssub{Les constructions sur les pentes : la surcharge du versant}

\textbf{Intuition.} Pose une brique sur un tas de sable incliné : le sable s'écroule. Le raisonnement est identique pour une maison construite sur un versant.

\begin{definition}[Surcharge]
On appelle \cle{surcharge} le poids supplémentaire apporté à un versant par les constructions (maisons, immeubles, routes, remblais). Cette surcharge augmente les forces qui « tirent » le sol vers le bas de la pente.
\end{definition}

Une maison en dur peut peser plusieurs dizaines de tonnes. Construite sur une pente, elle ajoute une force dirigée vers le bas qui s'ajoute au poids du sol lui-même. Si le versant était déjà proche de la limite de stabilité (pente raide, sol argileux, forte pluie), cette surcharge peut suffire à déclencher le glissement. De plus, les constructions modifient le \cle{drainage} des eaux : les toitures, les cours cimentées et les routes rendent le sol imperméable ; l'eau de pluie n'y pénètre plus et se concentre en certains points (pied des murs, bordures de route), où elle s'infiltre alors en grande quantité et fragilise localement le sol.

\begin{attention}[Le haut de la pente n'est pas un refuge !]
Erreur fréquente : croire que construire \emph{en haut} d'une pente protège du glissement. C'est le contraire : la surcharge exercée en haut du versant augmente la force motrice qui pousse tout le sol vers le bas. Une maison construite au sommet peut ainsi \emph{déclencher} le glissement qui l'emportera, elle et les maisons situées en contrebas.
\end{attention}

\begin{exemplebox}[Les quartiers sur pentes de Yaoundé]
Yaoundé, la « ville aux sept collines », connaît chaque saison des pluies des glissements de terrain et des effondrements de maisons. De nombreuses habitations ont été construites sur des pentes raides, sans étude de sol préalable et sans système d'évacuation des eaux. La surcharge des bâtiments, ajoutée à la saturation du sol par les pluies, provoque régulièrement des drames : maisons fissurées puis emportées, routes coupées, familles sinistrées.
\end{exemplebox}

\courssub{Les cultures sur pentes raides et le labour dans le sens de la pente}

Dans les régions montagneuses du Cameroun (Ouest, Nord-Ouest, Adamaoua), la pression démographique pousse les paysans à cultiver des pentes de plus en plus raides. Or ces pratiques fragilisent le sol de deux manières :
\begin{itemize}
\item \textbf{les cultures sur pentes raides} remplacent la forêt par des plantes aux racines courtes (maïs, haricot, manioc), qui retiennent mal le sol ; entre deux récoltes, le sol reste nu et exposé à l'érosion ;
\item \textbf{le labour dans le sens de la pente} crée des sillons verticaux qui deviennent de véritables « toboggans » pour l'eau de pluie : le ruissellement y gagne en vitesse, creuse des ravines et emporte la terre fertile.
\end{itemize}
La bonne pratique, à retenir, consiste à labourer \emph{perpendiculairement à la pente} (on dit « suivant les courbes de niveau ») et à planter des haies vives, ce qui freine l'eau et retient le sol. Nous retrouverons ces techniques dans la section consacrée à la prévention.

\courssub{Terrassements, routes et rejets d'eaux usées}

\textbf{Les terrassements et le creusement des routes.} Pour faire passer une route à flanc de colline, on creuse le versant : c'est le \cle{terrassement}. Cette opération a deux effets dangereux :
\begin{itemize}
\item elle crée un \cle{talus} (paroi raide) souvent trop vertical, donc instable, qui s'éboule en blocs ou en coulées ;
\item elle \emph{sape le pied} du versant : en retirant la base qui soutenait la masse de terre située au-dessus, on supprime le « contrefort » naturel, et tout le versant peut glisser.
\end{itemize}

\textbf{Les rejets d'eaux usées et les fuites de canalisations.} Dans les quartiers dépourvus d'assainissement, les eaux usées sont rejetées directement sur le sol ; de même, une canalisation d'eau qui fuit injecte en continu de l'eau dans le sous-sol. Il en résulte une \cle{saturation artificielle} du sol en eau, qui reproduit exactement les conditions d'un glissement de terrain étudiées dans la section précédente : alourdissement du sol, perte de cohésion, puis rupture.

\begin{popfigure}
\begin{tikzpicture}[scale=0.92]
% ===== Panneau gauche : versant boisé stable =====
\begin{scope}
\fill[popgreenL] (0,0) -- (0,2.6) -- (6,0.6) -- (6,0) -- cycle;
\draw[popdark,thick] (0,2.6) -- (6,0.6);
\draw[popdark,thick] (0,0) -- (0,2.6);
\draw[popdark,thick] (0,0) -- (6,0);
% arbres
\foreach \x/\y in {0.9/2.32, 2.2/1.96, 3.5/1.61, 4.8/1.25}{
  \draw[popdark,thick] (\x,\y) -- (\x,\y+0.55);
  \fill[popgreen!70!black] (\x,\y+0.8) circle (0.32);
  % racines
  \draw[poporange,thick] (\x,\y) -- (\x-0.25,\y-0.45);
  \draw[poporange,thick] (\x,\y) -- (\x+0.2,\y-0.5);
  \draw[poporange,thick] (\x,\y) -- (\x,\y-0.55);
}
% pluie interceptée
\foreach \x in {0.6,1.5,2.4,3.3,4.2,5.1}{
  \draw[popblue,->] (\x,3.6) -- (\x,3.1);
}
\node[popblue] at (3,3.9) {\small pluie};
\node[align=center] at (3,-0.7) {\small\textbf{Versant boisé : stable}};
\node[align=center,popmuted] at (3,-1.25) {\footnotesize racines = grillage retenant le sol};
\end{scope}
% ===== Panneau droit : versant déboisé, construit =====
\begin{scope}[xshift=7.5cm]
\fill[popgoldL] (0,0) -- (0,2.6) -- (6,0.6) -- (6,0) -- cycle;
\draw[popdark,thick] (0,2.6) -- (6,0.6);
\draw[popdark,thick] (0,0) -- (0,2.6);
\draw[popdark,thick] (0,0) -- (6,0);
% souches
\foreach \x/\y in {2.2/1.96, 3.5/1.61}{
  \draw[popdark,thick] (\x,\y) -- (\x,\y+0.18);
}
% maison
\fill[popink!80] (0.5,2.32) rectangle (1.5,3.02);
\draw[popdark,thick] (0.4,3.02) -- (1,3.42) -- (1.6,3.02);
\node[white] at (1,2.65) {\tiny maison};
% fissures
\draw[poppink,very thick] (1.8,2.05) -- (2.1,1.5) -- (1.9,1.0);
\draw[poppink,very thick] (3.2,1.68) -- (3.5,1.2) -- (3.3,0.75);
% glissement
\fill[poporange!60] (4.2,1.1) -- (6,0.55) -- (6,0) -- (4.6,0) -- cycle;
\draw[popink,->,very thick] (4.3,1.35) to[bend left=20] (5.3,0.55);
\node[popink] at (5.6,1.15) {\footnotesize glissement};
% pluie qui ruisselle
\foreach \x in {0.6,1.8,3.0,4.2,5.4}{
  \draw[popblue,->] (\x,3.6) -- (\x,3.1);
}
\draw[popblue,thick,->] (2.6,1.85) to[bend left=15] (4.6,0.9);
\node[popblue] at (3.4,2.35) {\footnotesize ruissellement};
\node[align=center] at (3,-0.7) {\small\textbf{Versant déboisé et construit : instable}};
\node[align=center,popmuted] at (3,-1.25) {\footnotesize surcharge + fissures + ruissellement};
\end{scope}
\end{tikzpicture}
\legende{Comparaison de deux versants voisins soumis à la même pluie. À gauche, le versant boisé : les racines (orange) retiennent le sol et pompent l'eau ; le versant reste stable. À droite, le versant déboisé et urbanisé : la surcharge de la maison, le ruissellement et la saturation en eau provoquent des fissures (rose) puis un glissement de terrain.}
\end{popfigure}

\begin{popfigure}
\begin{tikzpicture}[scale=0.95]
% versant
\fill[popgoldL] (0,0) -- (0,3.4) -- (5.2,3.4) -- (5.2,0) -- cycle;
\draw[popdark,thick] (0,3.4) -- (0,0) -- (5.2,0);
% route creusée
\fill[popmuted!60] (0,1.5) rectangle (3.4,1.95);
\draw[popdark,thick] (0,1.95) -- (3.4,1.95);
\draw[popdark,thick] (0,1.5) -- (3.4,1.5);
\node at (1.7,1.72) {\small route};
% talus
\draw[popdark,very thick] (3.4,1.5) -- (4.3,3.4);
\node[popink] at (4.55,2.6) {\small talus};
% éboulement au pied du talus
\fill[poporange!70] (3.4,1.5) -- (4.6,1.5) -- (3.9,2.3) -- cycle;
\draw[popink,->,thick] (4.1,2.9) -- (3.85,2.35);
\node[align=center,popink] at (5.3,1.35) {\footnotesize blocs};
\node[align=center,popink] at (5.3,1.0) {\footnotesize éboulés};
% voiture schématique
\fill[popblue] (1.2,1.95) rectangle (2.0,2.25);
\fill[popdark] (1.35,1.95) circle (0.09);
\fill[popdark] (1.85,1.95) circle (0.09);
% annotation sapement
\draw[popink,->,thick] (2.6,0.5) -- (3.35,1.35);
\node[align=center,popink] at (2.3,0.35) {\footnotesize pied du versant sapé};
\node[align=center] at (2.6,-0.6) {\small\textbf{Talus routier éboulé après le creusement de la route}};
\end{tikzpicture}
\legende{Schéma d'un talus routier éboulé. Le creusement de la route a supprimé le pied du versant (sapement) et créé un talus trop raide : des blocs se détachent et s'accumulent sur la chaussée, menaçant les usagers.}
\end{popfigure}

\courssub{Étude de cas documentaire : deux versants voisins, deux destins différents}

Pour vérifier notre hypothèse de départ (l'action de l'homme aggrave les mouvements de terrain), la méthode la plus rigoureuse consiste à comparer deux terrains \emph{identiques en tout point}, sauf pour l'action de l'homme. C'est le principe du \emph{témoin} que tu as déjà rencontré dans les expériences de ce cours.

\begin{experience}[Comparer pour comprendre : le versant boisé et le versant urbanisé]
\textbf{Document.} Deux versants voisins d'une même colline, de pente identique (environ 30°) et soumis à la même pluviométrie, sont observés pendant cinq saisons des pluies. Le versant A est resté boisé ; le versant B a été déboisé puis loti (maisons, chemin de terre, rejets d'eaux usées).

\begin{center}
\begin{tabularx}{\linewidth}{|l|X|X|}
\hline
\rowcolor{popblueL} \textbf{Critère observé} & \textbf{Versant A (boisé)} & \textbf{Versant B (déboisé, urbanisé)} \\
\hline
Couverture végétale & forêt dense, racines profondes & sol nu entre les maisons \\
\hline
Ruissellement de l'eau & faible, eau absorbée & fort, rigoles et ravines \\
\hline
Fissures du sol & aucune & nombreuses après chaque pluie \\
\hline
Glissements en 5 ans & 0 & 3 (dont 1 avec maisons détruites) \\
\hline
\end{tabularx}
\end{center}

\textbf{Interprétation.} La pente et la pluie étant identiques sur les deux versants, elles ne peuvent pas expliquer la différence observée. Les seuls facteurs qui changent sont liés à l'action de l'homme : déboisement (perte de la cohésion racinaire), surcharge des constructions, modification du drainage et saturation par les eaux usées.

\textbf{Conclusion.} L'action de l'homme a transformé un versant naturellement stable en versant instable : elle est bien une \emph{cause} des mouvements de terrain, et non une simple circonstance. L'hypothèse de départ est validée.
\end{experience}

\begin{exoresolu}[Analyser un glissement de terrain]
\textbf{Énoncé.} Dans un quartier de l'Ouest camerounais, une colline autrefois boisée a été défrichée, puis des maisons en dur ont été construites jusqu'au sommet de la pente. Une route a été creusée à mi-versant et les eaux usées sont rejetées sur le sol. Après trois jours de pluie, un glissement de terrain emporte quatre maisons. Dégage \textbf{quatre actions humaines} qui ont contribué à cet accident et explique le rôle de chacune.
\tcblower
\textbf{Solution détaillée.}
\begin{enumerate}
\item \textbf{Le déboisement} : la suppression des arbres a privé le sol du grillage de racines qui assurait sa cohésion et pompait l'eau ; le sol s'est fragilisé et s'est saturé plus vite.
\item \textbf{La surcharge par les constructions} : le poids des maisons en dur, y compris au sommet de la pente, a augmenté les forces qui poussent le sol vers le bas.
\item \textbf{Le creusement de la route} : il a sapé le pied du versant et créé un talus instable, supprimant le soutien naturel de la masse de terre située au-dessus.
\item \textbf{Les rejets d'eaux usées} : ils ont saturé artificiellement le sol en eau, ce qui l'a alourdi et a diminué la cohésion entre les particules.
\end{enumerate}
La pluie de trois jours a joué le rôle d'\emph{élément déclencheur}, mais ce sont les quatre actions humaines qui avaient rendu le versant instable. Sans elles, la même pluie sur la colline boisée n'aurait probablement causé aucun dégât.
\end{exoresolu}

\begin{savaistu}[Le reboisement, une solution officielle]
Le reboisement des versants est l'une des mesures recommandées par le MINEPDED (Ministère de l'Environnement, de la Protection de la Nature et du Développement Durable) dans les zones à risque du Cameroun. Planter des arbres sur une colline menacée, c'est reconstituer le « grillage » de racines qui retient le sol : une protection naturelle, peu coûteuse et durable contre les glissements de terrain.
\end{savaistu}

\begin{aretenir}[L'essentiel de la section]
L'action de l'homme aggrave ou déclenche les mouvements de terrain par cinq mécanismes principaux : le \cle{déboisement} (perte de la cohésion racinaire), les \cle{constructions sur les pentes} (surcharge et modification du drainage), les \cle{cultures et labours dans le sens de la pente} (érosion accélérée), les \cle{terrassements} (talus instables, pied du versant sapé) et les \cle{rejets d'eaux} (saturation artificielle du sol). Comprendre ces mécanismes, c'est déjà prévenir : c'est ce que nous verrons dans la section suivante, consacrée aux techniques de prévention et de protection.
\end{aretenir}

\coursec{Prévention et protection contre les mouvements de terrain}

Nous venons de voir que l'action de l'homme — déboisement, constructions mal placées, terrassements — peut aggraver les risques de mouvements de terrain. Mais l'homme n'est pas seulement un facteur d'aggravation : il peut aussi être un acteur de \cle{prévention} et de \cle{protection}. Comment, concrètement, peut-on empêcher un versant de glisser ou limiter les dégâts d'un éboulement ? C'est l'objet de cette section, qui te donnera une démarche complète d'analyse d'une situation à risque, très utile pour les exercices du BEPC.

\courssub{Connaître le terrain : la première étape de la prévention}

Avant de construire une maison, une route ou une école en zone montagneuse, il faut d'abord \textbf{connaître le terrain}. On ne peut pas se protéger d'un danger qu'on n'a pas identifié.

\begin{definition}[Prévention et protection]
La \cle{prévention} regroupe l'ensemble des mesures prises \emph{avant} un accident pour éviter qu'il ne se produise (études, aménagements, interdictions de construire). La \cle{protection} regroupe les mesures qui limitent les dégâts lorsque le mouvement de terrain se produit malgré tout (ouvrages de défense, évacuation des populations).
\end{definition}

La connaissance du terrain repose sur deux outils essentiels :

\begin{itemize}
\item \textbf{La cartographie des zones à risque} : les spécialistes observent les reliefs, repèrent les anciennes coulées de boue, les falaises éboulées, les pentes trop raides, puis dressent des \emph{cartes des zones à risque}. Ces cartes indiquent où il est dangereux de construire.
\item \textbf{Les études de sol avant construction} : avant tout chantier, on prélève des échantillons de sol et de roche pour connaître leur nature (argile, sable, rocher), leur profondeur et leur comportement face à l'eau. Un sol argileux sur une pente, par exemple, devient très glissant quand il est mouillé : c'est un signal d'alarme.
\end{itemize}

\begin{attention}[Ne jamais construire « au hasard »]
Beaucoup d'accidents en zone montagneuse viennent de constructions réalisées sans aucune étude de sol, sur d'anciennes coulées de boue ou en bordure de falaise. La première mesure de prévention ne coûte presque rien : c'est de \cle{choisir le bon endroit}.
\end{attention}

\courssub{Les techniques de stabilisation des versants}

Lorsqu'un versant menace de glisser ou de s'ébouler, on peut le \textbf{stabiliser} par des ouvrages ou des techniques mécaniques.

\begin{definition}[Mur de soutènement]
Un \cle{mur de soutènement} est un ouvrage (en béton, en pierres ou en gabions) qui retient les terres d'un versant pour empêcher leur glissement.
\end{definition}

Voici les principales techniques à connaître :

\begin{itemize}
\item \textbf{Les murs de soutènement} : placés au pied d'un talus, ils s'opposent au glissement des terres. On les rencontre le long des routes de montagne.
\item \textbf{Les gabions} : ce sont des cages métalliques remplies de pierres, empilées contre un talus. Lourds et perméables à l'eau, ils retiennent le sol tout en laissant l'eau s'écouler.
\item \textbf{Le clouage des versants} : on enfonce de longues tiges métalliques (les « clous ») dans la pente pour fixer la couche de terre meuble à la roche solide située en dessous.
\item \textbf{Les filets de protection} : tendus contre les falaises ou les parois rocheuses, ils retiennent les blocs rocheux et empêchent les éboulements d'atteindre les routes et les habitations.
\end{itemize}

\begin{exemplebox}[Les gabions sur les routes de l'Ouest camerounais]
Sur certains tronçons routiers de la région de l'Ouest, où les pluies sont abondantes et les reliefs accidentés, des \cle{gabions} (cages de pierres) sont installés le long des talus pour stabiliser les bords de route et éviter que les terres ne glissent sur la chaussée pendant la saison des pluies. C'est une technique simple, solide et peu coûteuse, bien adaptée au contexte local.
\end{exemplebox}

\begin{attention}[Un mur sans drainage peut aggraver le risque]
L'eau qui s'infiltre dans le sol s'accumule derrière un mur de soutènement. Si elle ne peut pas s'évacuer, elle exerce une forte pression sur le mur et alourdit la terre : le mur peut alors être renversé et le glissement devient encore plus violent. Il faut \textbf{toujours associer un mur de soutènement à un système de drainage} (trous d'évacuation dans le mur, drains dans le sol).
\end{attention}

\courssub{La gestion de l'eau : évacuer l'ennemi numéro un des pentes}

Nous l'avons vu dans les sections précédentes : l'eau est la principale cause des mouvements de terrain. Elle alourdit le sol, rend l'argile glissante et fragilise les roches. La solution logique consiste donc à \textbf{empêcher l'eau de s'accumuler dans la pente}.

\begin{itemize}
\item \textbf{Les fossés de crête} : creusés en haut de la pente, ils interceptent l'eau qui descend de la montagne avant qu'elle n'atteigne la zone fragile.
\item \textbf{Les caniveaux et canaux de drainage} : ils recueillent les eaux de pluie et les conduisent loin des pentes, vers des zones sûres.
\item \textbf{Les drains} : ce sont des conduites enterrées qui captent l'eau infiltrée dans le sol et l'évacuent, ce qui allège le terrain.
\end{itemize}

\begin{popfigure}
\begin{center}
\begin{tikzpicture}[scale=1.0, every node/.style={font=\small}]
% Sol / versant
\fill[popgoldL] (0,0) -- (10,0) -- (10,4.2) -- (6.5,4.2) -- (2.5,1.6) -- (0,1.6) -- cycle;
\draw[popdark, thick] (0,1.6) -- (2.5,1.6) -- (6.5,4.2) -- (10,4.2);
% Fossé de crête
\fill[popblue!50] (7.3,4.2) -- (7.7,3.7) -- (8.1,4.2) -- cycle;
\draw[popdark] (7.3,4.2) -- (7.7,3.7) -- (8.1,4.2);
\node[above, popblue] at (7.7,4.25) {fossé de crête};
% Arbres (reboisement)
\foreach \x/\y in {4.0/2.55, 5.0/3.05, 5.9/3.5}{
  \draw[popdark, thick] (\x,\y) -- (\x,\y+0.45);
  \fill[popgreen!70] (\x,\y+0.75) circle (0.32);
}
\node[popgreen!60!black, align=center] at (5.0,1.6) {reboisement\\ (racines profondes)};
% Racines schématiques
\foreach \x/\y in {4.0/2.55, 5.0/3.05, 5.9/3.5}{
  \draw[popdark] (\x,\y) -- (\x-0.2,\y-0.35);
  \draw[popdark] (\x,\y) -- (\x+0.2,\y-0.35);
}
% Mur de soutènement
\fill[popmuted!60] (2.2,0) rectangle (2.7,1.6);
\draw[popdark, thick] (2.2,0) rectangle (2.7,1.6);
\node[below, align=center] at (2.45,-0.05) {mur de\\ soutènement};
% Drain derrière le mur
\draw[popblue, thick, dashed] (3.2,1.4) -- (3.2,0.3) -- (2.75,0.15);
\node[popblue, align=left] at (4.6,0.55) {drain (évacuation\\ de l'eau infiltrée)};
% Caniveau au pied
\fill[popblue!50] (0.4,0) rectangle (1.6,0.25);
\draw[popdark] (0.4,0) rectangle (1.6,0.25);
\node[below, popblue] at (1.0,-0.05) {caniveau};
% Route
\draw[popdark, very thick] (0,0) -- (2.2,0);
\node[below left] at (0.2,0) {route};
% Flèche eau de pluie
\draw[->, popblue, thick] (8.6,4.9) -- (7.9,4.35);
\node[popblue, above] at (8.7,4.9) {eau de pluie};
\end{tikzpicture}
\end{center}
\legende{Schéma d'un versant aménagé contre les mouvements de terrain. Le fossé de crête intercepte l'eau de pluie en haut de la pente ; le reboisement fixe le sol grâce aux racines ; le mur de soutènement retient les terres au pied du talus ; le drain évacue l'eau infiltrée derrière le mur ; le caniveau protège la route.}
\end{popfigure}

\courssub{Les solutions végétales : la nature au service de la stabilité}

Les plantes sont des alliées précieuses contre les mouvements de terrain. Leurs \textbf{racines} agissent comme un véritable filet naturel qui retient les particules du sol, et leur feuillage absorbe une partie de l'eau de pluie.

\begin{itemize}
\item \textbf{Le reboisement des versants} : planter des arbres sur les pentes dénudées fixe le sol et réduit l'érosion. C'est l'inverse du déboisement, qui, lui, aggrave le risque.
\item \textbf{Les espèces à racines profondes} : on choisit de préférence des arbres dont les racines s'enfoncent loin dans le sol, car ils ancrent les couches profondes et pas seulement la surface.
\item \textbf{Les bandes boisées} : des rangées d'arbres plantées le long des pentes ou des routes freinent les coulées de boue et retiennent les blocs rocheux.
\end{itemize}

\begin{savaistu}[Des solutions à la portée des villages]
Le reboisement est une technique de prévention peu coûteuse que les communautés locales peuvent réaliser elles-mêmes. Dans plusieurs villages des régions montagneuses du Cameroun, des campagnes de plantation d'arbres sur les versants dénudés ont permis de réduire l'érosion des sols et de protéger les champs et les cases contre les coulées de boue de la saison des pluies.
\end{savaistu}

\courssub{Les mesures d'urbanisme et l'éducation des populations}

Les meilleurs ouvrages ne suffisent pas si les constructions sont mal placées ou si les populations ne savent pas réagir.

\textbf{Les mesures d'urbanisme} consistent à :
\begin{itemize}
\item \textbf{interdire la construction dans les zones à haut risque} (bas de falaises, anciennes coulées de boue, pentes très raides) ;
\item faire \textbf{respecter les normes de construction} : fondations adaptées, étude de sol obligatoire, limitation de la hauteur des talus creusés.
\end{itemize}

\textbf{L'éducation et l'alerte des populations} sont tout aussi importantes. Il faut apprendre à reconnaître les \cle{signes avant-coureurs} d'un mouvement de terrain :
\begin{itemize}
\item des \textbf{fissures} qui apparaissent dans les murs des maisons ou dans le sol ;
\item des \textbf{arbres ou poteaux qui penchent} anormalement ;
\item des \textbf{eaux troubles} dans les sources ou les ruisseaux (signe que la terre se déplace) ;
\item de petits éboulements répétés au pied d'une falaise.
\end{itemize}

Chaque communauté exposée doit disposer d'un \textbf{plan d'évacuation} : un itinéraire connu de tous, un lieu de rassemblement en sécurité et un responsable chargé de donner l'alerte.

\begin{savaistu}[Des fissures en haut de la pente : alerte immédiate !]
Des fissures apparaissant dans le sol en haut d'une pente sont un \cle{signe avant-coureur} de glissement de terrain : elles montrent que la masse de terre commence à se détacher. Si tu observes de telles fissures, il faut \textbf{alerter immédiatement} les adultes et les autorités, et s'éloigner de la zone. Quelques heures ou quelques jours peuvent séparer l'apparition des fissures du glissement : ce temps précieux peut sauver des vies.
\end{savaistu}

\textbf{Conduite à tenir en cas d'alerte} :
\begin{enumerate}
\item \textbf{Évacuer} immédiatement la zone menacée, en suivant le plan d'évacuation, sans perdre de temps à emporter des biens.
\item \textbf{Prévenir les autorités} (chef de quartier, mairie, gendarmerie, protection civile) et les voisins.
\item \textbf{Ne jamais traverser une coulée de boue} ou une zone d'éboulement, même à pied : la boue est profonde, instable et peut emporter une personne en quelques secondes.
\end{enumerate}

\courssub{Tableau de synthèse : à chaque cause, sa technique de prévention}

\begin{center}
\begin{tabularx}{\linewidth}{|X|X|}
\hline
\rowcolor{popblueL} \textbf{Cause du risque} & \textbf{Technique de prévention ou de protection} \\
\hline
Pente trop raide, terres qui glissent & Mur de soutènement, gabions, clouage du versant \\
\hline
Accumulation d'eau dans le sol & Fossé de crête, caniveaux, drains (drainage) \\
\hline
Sol dénudé, érosion par les pluies & Reboisement, plantes à racines profondes, bandes boisées \\
\hline
Chute de blocs rocheux (éboulements) & Filets de protection, murs pare-blocs \\
\hline
Constructions en zone dangereuse & Cartographie des risques, interdiction de construire, normes de construction \\
\hline
Populations non informées & Éducation, surveillance des signes avant-coureurs, plan d'évacuation \\
\hline
\end{tabularx}
\end{center}

\courssub{La démarche d'analyse d'une situation à risque}

Face à un exercice du BEPC ou à une situation concrète, il faut raisonner de façon méthodique.

\begin{methode}[Analyser une situation à risque en 3 étapes]
\begin{enumerate}
\item \textbf{Identifier la cause} du risque : pente raide ? eau ? déboisement ? construction mal placée ?
\item \textbf{Choisir la technique adaptée} à cette cause précise (voir le tableau de synthèse).
\item \textbf{Justifier le choix} : expliquer \emph{pourquoi} cette technique supprime ou réduit la cause identifiée.
\end{enumerate}
\end{methode}

\begin{exoresolu}[Un village menacé par les coulées de boue]
\textbf{Énoncé.} Un village est construit au pied d'une colline dont la pente a été entièrement déboisée pour faire des champs. Chaque saison des pluies, des coulées de boue atteignent les premières cases. (1) Identifie les causes de ces accidents. (2) Propose deux mesures de prévention adaptées et justifie ton choix.
\tcblower
\textbf{Solution détaillée.}
\begin{enumerate}
\item \emph{Causes} : (a) la pente est raide et le sol n'est plus retenu ; (b) le \textbf{déboisement} a supprimé les racines qui fixaient le sol ; (c) l'\textbf{eau de pluie} s'infiltre et ruisselle librement, alourdissant la terre qui glisse sous forme de boue.
\item \emph{Mesures proposées} :
\begin{itemize}
\item \textbf{Reboiser la pente} avec des arbres à racines profondes : les racines fixeront le sol et le feuillage réduira l'impact de la pluie. Cette mesure s'attaque directement à la cause (le déboisement).
\item \textbf{Creuser un fossé de crête et des caniveaux} pour évacuer l'eau de pluie loin de la pente : moins d'eau dans le sol, donc moins de risque de coulée de boue. On peut ajouter un mur de soutènement \emph{avec drainage} au pied de la colline pour protéger les cases.
\end{itemize}
\end{enumerate}
\emph{Conclusion} : la prévention efficace combine toujours plusieurs techniques qui agissent sur les différentes causes du risque.
\end{exoresolu}

\begin{aretenir}[L'essentiel de la section]
\begin{itemize}
\item La prévention commence par la \cle{connaissance du terrain} : cartographie des zones à risque et études de sol avant toute construction.
\item Les techniques de stabilisation : \cle{murs de soutènement}, \cle{gabions}, \cle{clouage}, \cle{filets de protection}.
\item L'eau doit être évacuée loin des pentes : \cle{fossés de crête}, \cle{caniveaux}, \cle{drains}.
\item Le \cle{reboisement} fixe les sols grâce aux racines des arbres.
\item Il faut interdire les constructions en zone à haut risque et éduquer les populations : reconnaître les \cle{signes avant-coureurs} (fissures, arbres penchés, eaux troubles) et appliquer le plan d'évacuation.
\item En cas d'alerte : évacuer, prévenir les autorités, ne jamais traverser une coulée de boue.
\end{itemize}
\end{aretenir}

\coursec{Synthèse et démarche d'analyse d'un accident}

Dans la section précédente, nous avons vu qu'il existe des techniques de \cle{prévention} (reboisement, drainage, murs de soutènement) et de \cle{protection} (évacuation, alerte, aménagement des constructions) contre les mouvements de terrain. Mais toutes ces mesures ne sont efficaces que si l'on a d'abord \textbf{compris les causes} de l'accident : on ne soigne bien un malade que si le diagnostic est juste ! Cette dernière section te donne la méthode complète pour analyser, comme un vrai scientifique, un accident lié à un mouvement de terrain, et pour rédiger une conclusion argumentée — une compétence directement évaluée au BEPC.

\courssub{Bilan des causes des mouvements de terrain}

Faisons le point sur tout ce que nous avons appris dans cette leçon. Un \cle{mouvement de terrain} (glissement, éboulement, coulée de boue, affaissement) résulte toujours d'une combinaison de facteurs que l'on peut classer en deux grandes familles.

\begin{aretenir}[Les deux familles de causes]
\begin{itemize}
\item \textbf{Causes naturelles :}
\begin{itemize}
\item la \cle{pente} : plus le versant est raide, plus les matériaux sont instables ;
\item la nature des \cle{roches} : des roches fracturées (déformations cassantes : failles, diaclases) ou altérées offrent des surfaces de glissement ;
\item les \cle{pluies} : l'eau alourdit le sol, réduit la cohésion des matériaux et fragilise les roches par action mécanique et chimique ;
\item les \cle{déformations souples} (plis) peuvent redresser des couches de roches et les rendre glissantes.
\end{itemize}
\item \textbf{Causes humaines :}
\begin{itemize}
\item le \cle{déboisement} : sans racines, le sol n'est plus retenu ;
\item les \cle{constructions} mal placées (sur versants raides, en haut de falaise) qui surchargent le sol ;
\item les terrassements, déblais routiers et carrières qui coupent le pied des versants.
\end{itemize}
\end{itemize}
\end{aretenir}

\begin{attention}[Un accident a rarement une seule cause]
Dans la réalité comme dans les exercices du BEPC, un glissement de terrain résulte presque toujours de la \textbf{combinaison} de plusieurs facteurs : une cause naturelle de fond (pente, roche fracturée) + un élément déclencheur (fortes pluies) + souvent une aggravation humaine (déboisement). Cherche systématiquement plusieurs causes dans les documents !
\end{attention}

\courssub{Schéma-bilan de la leçon}

La carte mentale ci-dessous résume l'ensemble de la leçon : apprends à la refaire de mémoire, c'est un excellent outil de révision.

\begin{popfigure}
\begin{center}
\begin{tikzpicture}[
  grow=right,
  level 1/.style={sibling distance=3.4cm, level distance=4.2cm},
  level 2/.style={sibling distance=1.15cm, level distance=3.6cm},
  every node/.style={font=\small},
  racine/.style={rectangle, rounded corners, draw=popink, fill=popink, text=white, font=\bfseries, align=center, inner sep=5pt},
  branche/.style={rectangle, rounded corners, draw=popblue, fill=popblueL, align=center, inner sep=4pt},
  feuille/.style={rectangle, rounded corners, draw=popmuted!60, fill=popcream, align=center, inner sep=3pt, font=\footnotesize}
]
\node[racine, align=center]{Mouvements\\de terrain}
  child { node[branche, align=center]{Prévention et\\protection}
    child { node[feuille, align=center]{reboisement, drainage,\\murs de soutènement} }
    child { node[feuille, align=center]{éviter de construire\\sur les versants raides} }
    child { node[feuille, align=center]{surveillance, alerte,\\évacuation} }
  }
  child { node[branche] {Causes}
    child { node[feuille, align=center]{Homme : déboisement,\\constructions, terrassements} }
    child { node[feuille, align=center]{Eau : action mécanique\\(cohésion) et chimique (altération)} }
    child { node[feuille, align=center]{Déformations : souples (plis)\\et cassantes (failles, diaclases)} }
  };
\end{tikzpicture}
\end{center}
\legende{Carte mentale-bilan de la leçon : les mouvements de terrain ont des causes (déformations des roches, action de l'eau, action de l'homme) contre lesquelles on lutte par des techniques de prévention et de protection.}
\end{popfigure}

\courssub{La démarche d'analyse d'un accident}

Au BEPC, on te donnera souvent des documents (photographie d'un versant effondré, tableau des pluies, carte d'une région, extrait de presse) et une consigne du type : « À partir de l'exploitation des documents, expliquez l'origine de cet accident et proposez des mesures de prévention. » Voici la démarche rigoureuse à suivre, en quatre étapes.

\begin{methode}[Analyser un accident de mouvement de terrain en 4 étapes]
\begin{enumerate}
\item \textbf{Observer les documents} : lire attentivement chaque document (photo, carte, tableau, texte) et relever les faits : pente raide ? arbres absents ? pluies abondantes ? constructions ? roches fissurées ?
\item \textbf{Identifier les facteurs} : classer chaque fait relevé en cause naturelle (pente, roche, pluie) ou cause humaine (déboisement, construction, terrassement).
\item \textbf{Proposer une explication justifiée} : relier chaque facteur au mécanisme du cours (ex. : « l'eau réduit la cohésion du sol », « les racines retiennent le sol ») pour expliquer comment l'accident s'est produit.
\item \textbf{Proposer des mesures} : pour chaque cause identifiée, proposer une mesure de prévention ou de protection qui la combat directement.
\end{enumerate}
\end{methode}

\begin{methode}[Rédiger une réponse justifiée]
Chaque idée de ta réponse doit suivre le schéma en trois temps :
\begin{enumerate}
\item une \textbf{affirmation} (ce que tu déduis) ;
\item une \textbf{preuve tirée du document} (ce que tu as observé, avec sa référence : « photo 1 », « tableau 2 »...) ;
\item le \textbf{lien avec le savoir du cours} (le mécanisme scientifique qui explique le fait).
\end{enumerate}
Formule type : « \emph{[affirmation]} car \emph{[preuve du document]} ; or \emph{[savoir du cours]}. »
\end{methode}

\begin{attention}[Le piège classique de la copie]
Ne te contente \textbf{jamais} de décrire le document (« on voit une colline, des maisons, de la boue ») : une description ne rapporte aucun point. Il faut \textbf{exploiter} le document pour répondre à la question posée : chaque observation doit servir à identifier une cause et à construire ton explication. Pose-toi toujours la question : « en quoi ce que je vois m'aide-t-il à répondre ? »
\end{attention}

\begin{exemplebox}[Exemple de phrase correctement justifiée]
« Le versant a glissé car il était déboisé (photo 1 : absence d'arbres) et gorgé d'eau après 3~jours de pluie (tableau 2) : or les racines retiennent le sol et l'eau réduit sa cohésion. »
\par\smallskip
On y retrouve bien : l'affirmation (« le versant a glissé car... »), les preuves référencées (photo 1, tableau 2) et les savoirs du cours (rôle des racines, rôle de l'eau).
\end{exemplebox}

\courssub{Rédiger une conclusion argumentée}

La conclusion d'une analyse d'accident doit faire le lien entre \textbf{chaque cause} et \textbf{la mesure qui la combat}. C'est la logique du diagnostic et du remède.

\begin{center}
\begin{tabularx}{\linewidth}{|X|X|}
\hline
\rowcolor{popblueL} \textbf{Cause identifiée} & \textbf{Mesure de prévention / protection} \\
\hline
Déboisement du versant & Reboisement (les racines fixent le sol) \\
\hline
Sol gorgé d'eau (fortes pluies) & Drainage, canaux d'évacuation des eaux \\
\hline
Pente raide et roches fracturées & Murs de soutènement, filets de protection, interdiction de construire \\
\hline
Constructions sur le versant & Choisir des zones planes et stables pour bâtir ; démolir ou déplacer les constructions menacées \\
\hline
Risque persistant malgré tout & Surveillance du versant, système d'alerte, plan d'évacuation des populations \\
\hline
\end{tabularx}
\end{center}

\begin{exoresolu}[Analyse d'un glissement de terrain dans les hauts plateaux de l'Ouest]
\textbf{Énoncé.} Après une semaine de pluies abondantes, un versant cultivé et totalement déboisé situé au-dessus d'un village de la région de l'Ouest du Cameroun s'est effondré : une coulée de boue a enseveli plusieurs maisons construites au pied de la pente. Les documents montrent : (doc. 1) une photo du versant sans arbres, avec des cultures en rangées le long de la pente ; (doc. 2) un tableau indiquant \SI{180}{mm} de pluie en 6~jours ; (doc. 3) une coupe montrant des roches traversées par de nombreuses fissures.
\begin{enumerate}
\item Identifie les causes naturelles et humaines de l'accident.
\item Rédige une explication justifiée de l'accident.
\item Propose trois mesures pour éviter qu'un tel accident se reproduise.
\end{enumerate}
\tcblower
\textbf{Solution détaillée.}
\begin{enumerate}
\item \textbf{Causes naturelles} : la pente du versant (doc. 1), les pluies abondantes : \SI{180}{mm} en 6~jours (doc. 2), les roches fracturées (doc. 3). \textbf{Causes humaines} : le déboisement total du versant et les cultures qui ne retiennent pas le sol (doc. 1), les maisons construites au pied de la pente (doc. 1).
\item \textbf{Explication justifiée} : « Le versant s'est effondré car il était déboisé (doc. 1 : absence d'arbres) et gorgé d'eau après \SI{180}{mm} de pluie en 6~jours (doc. 2) ; or les racines des arbres retiennent le sol, et l'eau alourdit les matériaux et réduit leur cohésion. De plus, les roches fracturées (doc. 3) ont offert des surfaces de glissement, car les fissures fragilisent la roche et laissent pénétrer l'eau. »
\item \textbf{Mesures} (chaque mesure reliée à sa cause) : reboiser le versant (contre le déboisement) ; creuser des canaux de drainage pour évacuer l'eau de pluie (contre l'engorgement du sol) ; interdire les constructions au pied du versant et déplacer les habitations menacées (protection des populations).
\end{enumerate}
\end{exoresolu}

\courssub{Pour aller plus loin : vulnérabilité et gestion des risques}

\begin{savaistu}[Notion complémentaire (hors strict programme)]
Les spécialistes distinguent deux idées : l'\cle{aléa} (le phénomène dangereux lui-même, par exemple le glissement de terrain) et la \cle{vulnérabilité} (la fragilité des personnes et des biens exposés : maisons mal construites, absence de plan d'évacuation). Le \textbf{risque}, c'est la rencontre des deux : risque = aléa $\times$ vulnérabilité. La \cle{gestion des risques naturels} consiste donc à agir sur les deux : surveiller l'aléa (études des versants, cartes de zones dangereuses) et réduire la vulnérabilité (construire ailleurs, informer et entraîner les populations). Cette idée te sera utile en classe de seconde.
\end{savaistu}

\courssub{Auto-évaluation}

Vérifie que tu sais :

\begin{itemize}
\item[$\square$] définir un mouvement de terrain et citer des exemples (glissement, éboulement, coulée de boue) ;
\item[$\square$] distinguer une déformation souple (pli) d'une déformation cassante (faille, diaclase) ;
\item[$\square$] expliquer le rôle de l'eau (action mécanique et chimique) et le rôle de l'homme (déboisement, constructions) dans les mouvements de terrain ;
\item[$\square$] citer des techniques de prévention (reboisement, drainage, soutènement) et de protection (alerte, évacuation) ;
\item[$\square$] analyser un accident en 4 étapes : observer $\rightarrow$ identifier les facteurs $\rightarrow$ expliquer en justifiant $\rightarrow$ proposer des mesures ;
\item[$\square$] rédiger une réponse justifiée : affirmation + preuve tirée du document + savoir du cours.
\end{itemize}

\begin{exercice}[À toi de jouer]
Un quartier de Yaoundé est construit sur une colline raide. Après de fortes pluies, plusieurs parcelles de terrain se sont affaissées et des maisons se sont fissurées. On observe que la colline a été déboisée pour construire et qu'aucun canal d'évacuation des eaux n'existe. Rédige un court paragraphe (5 à 8 lignes) expliquant l'accident et proposant deux mesures de prévention, en utilisant la méthode « affirmation + preuve + savoir ».
\end{exercice}

\begin{corrige}[Exercice]
Éléments attendus : l'affaissement s'explique par la combinaison d'une pente raide, de pluies abondantes qui ont gorgé le sol d'eau (or l'eau réduit la cohésion des matériaux) et du déboisement de la colline (or les racines des arbres retiennent le sol). Mesures : reboiser les parties non construites de la colline et creuser des canaux de drainage pour évacuer l'eau de pluie ; on peut aussi accepter : interdire toute nouvelle construction sur les pentes raides. La copie doit relier explicitement chaque mesure à la cause qu'elle combat.
\end{corrige}

Tu es maintenant capable d'analyser n'importe quel accident lié à un mouvement de terrain et de proposer des solutions concrètes : c'est exactement la démarche d'un citoyen éclairé... et celle que le BEPC attend de toi !

\newpage

\coursec{Activité d'intégration}

\courssub{Objectifs de l'activité}

À l'issue de cette activité d'intégration, tu seras capable de :
\begin{itemize}
\item extraire et interpréter des informations à partir de documents variés (article de presse, photographies, tableau de données, coupe géologique) ;
\item identifier et hiérarchiser les causes naturelles et humaines d'un mouvement de terrain réel ;
\item mobiliser les notions de \cle{déformation souple} et \cle{cassante} des roches, d'\cle{action mécanique et chimique de l'eau} et d'\cle{action de l'homme} pour expliquer un accident ;
\item proposer des mesures de \cle{prévention} et de \cle{protection} adaptées, en justifiant chacune par la cause qu'elle combat ;
\item rédiger une démarche argumentée et la présenter oralement avec clarté.
\end{itemize}

\courssub{Situation}

\textbf{Le drame de Bafoussam (Ouest-Cameroun), octobre 2019.}

Dans la nuit du 28 au 29 octobre 2019, vers 22 h, un pan entier de colline s'est détaché dans le quartier Gouache, à Bafoussam, chef-lieu de la région de l'Ouest. Une masse de terre et de roches a dévalé la pente et enseveli plusieurs habitations construites à flanc de versant. Le bilan officiel a fait état de plus de 40 morts et de nombreuses maisons détruites. Ce drame est l'un des glissements de terrain les plus meurtriers de l'histoire récente du Cameroun.

Tu es membre d'une équipe d'élèves scientifiques chargée par le conseil municipal de rédiger un rapport expliquant les causes de l'accident et proposant des mesures pour éviter qu'il ne se reproduise. Tu disposes du dossier documentaire suivant.

\textbf{Document 1 : extrait d'article de presse (résumé).}

« Il pleuvait abondamment depuis plusieurs jours sur Bafoussam. Dans la nuit du 28 au 29 octobre, les habitants ont entendu un grondement : la colline a glissé d'un coup, emportant une dizaine de maisons construites sur la pente. Les rescapés racontent que des fissures étaient apparues dans les murs de certaines maisons et sur le sol, plusieurs jours avant la catastrophe. »

\textbf{Document 2 : photographies du quartier avant et après la catastrophe.}

Avant : la pente, raide (environ $35^{\circ}$), est presque totalement déboisée ; les arbres ont été coupés pour construire des maisons et cultiver. Aucune construction de soutènement, aucun canal d'évacuation des eaux de pluie. Après : une large cicatrice de terre brune dénudée descend du sommet jusqu'aux maisons détruites ; les arbres restants, en contrebas, ont été couchés par la coulée.

\textbf{Document 3 : pluies tombées à Bafoussam les 10 jours précédant l'accident.}

\begin{center}
\begin{tabularx}{\linewidth}{|l|*{10}{>{\centering\arraybackslash}X|}}
\hline
\rowcolor{popblueL} Jour & J-10 & J-9 & J-8 & J-7 & J-6 & J-5 & J-4 & J-3 & J-2 & J-1 \\
\hline
Pluie (mm) & 2 & 0 & 5 & 12 & 18 & 25 & 30 & 42 & 55 & 48 \\
\hline
\end{tabularx}
\end{center}

Rappel : la hauteur totale de pluie sur ces 10 jours est de $\SI{237}{mm}$, dont plus de $\SI{175}{mm}$ sur les 4 derniers jours. La moyenne mensuelle d'octobre à Bafoussam est d'environ $\SI{280}{mm}$.

\textbf{Document 4 : coupe géologique simplifiée du versant.}

\begin{popfigure}
\begin{tikzpicture}[scale=1.0]
% socle fracturé
\fill[poppurple!35] (0,0) -- (8,0) -- (8,2.2) -- (0,1.2) -- cycle;
\draw[poppurple!35] (0,0) -- (8,0) -- (8,2.2) -- (0,1.2) -- cycle;
% fractures
\draw[popdark, thick] (1.5,0.2) -- (2.1,1.35);
\draw[popdark, thick] (3.5,0.3) -- (4.0,1.6);
\draw[popdark, thick] (5.5,0.4) -- (6.0,1.9);
\draw[popdark, thick] (7.0,0.5) -- (7.4,2.1);
% altérites
\fill[poporange!40] (0,1.2) -- (8,2.2) -- (8,3.6) -- (0,2.6) -- cycle;
\draw[poporange!40] (0,1.2) -- (8,2.2) -- (8,3.6) -- (0,2.6) -- cycle;
% surface
\draw[popdark, very thick] (0,2.6) -- (8,3.6);
% maisons
\draw[popdark] (5.6,3.28) rectangle (6.2,3.75);
\draw[popdark] (5.5,3.75) -- (5.9,4.0) -- (6.3,3.75);
\draw[popdark] (6.6,3.38) rectangle (7.2,3.85);
\draw[popdark] (6.5,3.85) -- (6.9,4.1) -- (7.3,3.85);
% flèche de glissement
\draw[-{Stealth[length=3mm]}, poppink, very thick] (4.2,3.2) .. controls (5.0,2.9) .. (6.2,2.6);
% pluie
\foreach \x in {1,2,3,4,5,6,7} {\draw[-{Stealth[length=1.5mm]}, popblue] (\x,5.2) -- (\x-0.15,4.6);}
% légendes internes
\node[popdark, font=\small] at (4,0.8) {Socle fracturé (roche dure fissurée)};
\node[popdark, font=\small] at (3.2,2.05) {Altérites argileuses (couche meuble, 3 à 5 m)};
\node[popblue, font=\small] at (1.2,5.4) {Pluies abondantes};
\node[popink, font=\small] at (6.9,4.4) {Maisons};
\node[poppink, font=\small] at (3.6,3.6) {Glissement};
\end{tikzpicture}
\legende{Coupe simplifiée du versant de Gouache : une couche d'altérites argileuses (terre meuble issue de l'altération des roches) repose sur un socle rocheux fracturé. Les maisons sont construites sur la pente.}
\end{popfigure}

\courssub{Consignes}

Travail en groupes de 4 élèves, puis restitution orale (5 minutes par groupe) et correction collective.

\begin{enumerate}
\item \textbf{Document 3 :} calcule la hauteur totale de pluie tombée pendant les 10 jours, puis pendant les 4 derniers jours. Compare ces valeurs à la moyenne mensuelle d'octobre ($\SI{280}{mm}$). Que remarques-tu ?
\item \textbf{Documents 1 et 3 :} quel signe annonciateur de l'accident est mentionné dans l'article ? À quel type de déformation des roches et du sol (souple ou cassante) ces fissures correspondent-elles ? Justifie ta réponse.
\item \textbf{Document 4 :} décris les deux couches du sous-sol du versant. Laquelle est susceptible de glisser ? Sur quelle surface peut-elle s'écouler ?
\item Explique, en t'appuyant sur tes connaissances, comment l'eau de pluie a pu déclencher le glissement (action mécanique et action chimique de l'eau sur les argiles).
\item \textbf{Document 2 :} relève dans les photographies les indices montrant l'action de l'homme sur ce versant. Explique comment chacun a pu favoriser l'accident.
\item Dresse la liste de toutes les causes de l'accident, en les classant en deux colonnes : \emph{causes naturelles} et \emph{causes humaines}. Entoure la cause déclenchante.
\item Rédige en 8 à 12 lignes une explication argumentée de l'accident, en utilisant au moins quatre mots-clés de la leçon : \cle{glissement de terrain}, \cle{déformation cassante}, \cle{altérites}, \cle{eau}, \cle{déboisement}.
\item Propose au maire de Bafoussam \textbf{trois mesures de prévention} adaptées à ce quartier. Pour chacune, précise la cause qu'elle combat et justifie ton choix.
\end{enumerate}

\courssub{Ressources à mobiliser}

\begin{aretenir}[Ce que tu dois réutiliser]
\begin{itemize}
\item Un \cle{mouvement de terrain} est le déplacement d'une masse de roches ou de terre sous l'effet de la pesanteur (glissement, éboulement, coulée de boue).
\item Les roches subissent des \cle{déformations souples} (plis, sans rupture) ou \cle{cassantes} (failles, fissures, avec rupture).
\item L'eau agit \textbf{mécaniquement} : elle alourdit le sol, remplit les fissures, rend la couche meuble plus lourde et plus glissante. Elle agit \textbf{chimiquement} : elle altère les roches et transforme certains minéraux en argiles (formation des \cle{altérites}), qui deviennent plastiques et glissantes quand elles sont mouillées.
\item L'homme aggrave les risques par le \cle{déboisement} (les racines ne retiennent plus le sol), les constructions sur les pentes, les terrassements et le détournement des eaux.
\item Prévention : reboisement, interdiction de construire sur les pentes raides, canaux de drainage, murs de soutènement, surveillance des fissures, sensibilisation des populations.
\end{itemize}
\end{aretenir}

\courssub{Démarche et corrigé commenté}

\begin{exoresolu}[Rapport d'expertise sur le glissement de Bafoussam]
Énoncé : à partir des quatre documents, expliquer les causes de l'accident de Gouache et proposer trois mesures de prévention justifiées.
\tcblower
\textbf{Étape 1 : analyse des pluies (question 1).}

Total sur 10 jours : $2+0+5+12+18+25+30+42+55+48 = \SI{237}{mm}$.

Total sur les 4 derniers jours : $30+42+55+48 = \SI{175}{mm}$.

Or $\dfrac{175}{237} \approx 0,74$ : environ $74\,\%$ de la pluie est tombée en seulement 4 jours. De plus, les $\SI{237}{mm}$ en 10 jours représentent près de $85\,\%$ de la moyenne mensuelle ($\SI{280}{mm}$). \emph{Conclusion : des pluies exceptionnellement abondantes et concentrées ont saturé le sol en eau.}

\textbf{Étape 2 : les signes annonciateurs (question 2).}

L'article signale des fissures dans les murs et sur le sol plusieurs jours avant. Une fissure est une rupture : il s'agit d'une \cle{déformation cassante}. Elle montre que la masse de terre commençait déjà à se fracturer et à se déplacer lentement avant la rupture finale.

\textbf{Étape 3 : lecture de la coupe (question 3).}

Le sous-sol comprend : en profondeur, un \textbf{socle rocheux dur mais fracturé} ; en surface, une couche de \textbf{3 à 5 m d'altérites argileuses}, terre meuble issue de l'altération chimique de la roche. C'est la couche d'altérites qui a glissé, en s'écoulant sur la surface inclinée du socle, qui joue le rôle de « toboggan ».

\textbf{Étape 4 : rôle de l'eau (question 4).}

\emph{Action mécanique :} l'eau s'est infiltrée dans le sol et les fissures du socle ; elle a alourdi la couche d'altérites et formé, au contact du socle, une surface boueuse très glissante. \emph{Action chimique :} depuis longtemps, l'eau a altéré la roche en argiles ; mouillées, ces argiles deviennent plastiques et perdent leur cohésion. La couche saturée, trop lourde, a fini par glisser sur le socle incliné.

\textbf{Étape 5 : rôle de l'homme (question 5).}

Les photographies montrent : (a) un \textbf{déboisement} quasi total — sans racines, le sol n'est plus retenu et l'eau s'infiltre plus vite ; (b) des \textbf{maisons construites sur une pente raide} ($35^{\circ}$) — leur poids surcharge la couche meuble ; (c) \textbf{aucun canal d'évacuation des eaux} ni mur de soutènement.

\textbf{Étape 6 : hiérarchisation des causes (question 6).}

\begin{center}
\begin{tabularx}{\linewidth}{|X|X|}
\hline
\rowcolor{popblueL} \textbf{Causes naturelles} & \textbf{Causes humaines} \\
\hline
Pente raide ($35^{\circ}$) ; couche d'altérites argileuses sur socle fracturé ; \fbox{pluies exceptionnelles (cause déclenchante)} & Déboisement du versant ; constructions sur la pente ; absence de drainage et de soutènement \\
\hline
\end{tabularx}
\end{center}

\textbf{Étape 7 : explication rédigée (question 7) — exemple de réponse.}

« Le drame de Gouache est un \cle{glissement de terrain} : une couche d'\cle{altérites} argileuses, formée par l'altération chimique de la roche par l'\cle{eau}, reposait sur un socle fracturé incliné. Les pluies exceptionnelles ($\SI{237}{mm}$ en 10 jours) ont saturé cette couche, l'ont alourdie et ont rendu son contact avec le socle glissant. Le \cle{déboisement} du versant et le poids des maisons ont supprimé les éléments qui retenaient le sol. Des fissures, signes d'une \cle{déformation cassante}, annonçaient la rupture. Dans la nuit du 28 au 29 octobre, la couche saturée a glissé d'un coup, ensevelissant les habitations. »

\textbf{Étape 8 : trois mesures de prévention (question 8).}

\begin{enumerate}
\item \textbf{Reboiser le versant} avec des arbres à racines profondes : combat le déboisement ; les racines retiennent le sol et absorbent une partie de l'eau.
\item \textbf{Interdire toute nouvelle construction sur la pente et reloger les familles} en zone sûre : combat la surcharge du versant et protège les vies.
\item \textbf{Creuser des canaux de drainage} en haut et sur les flancs du versant (et éventuellement un mur de soutènement au pied) : combat l'accumulation de l'eau dans le sol, cause déclenchante du glissement.
\end{enumerate}
\end{exoresolu}

\begin{attention}[Grille d'évaluation de la restitution]
Chaque exposé est noté sur trois critères : \textbf{exactitude} (les causes citées sont correctes et appuyées sur les documents), \textbf{justification} (chaque mesure proposée est reliée à une cause précise), \textbf{clarté} (vocabulaire scientifique exact : glissement, altérites, déformation cassante, drainage). Erreur fréquente à éviter : affirmer que « la pluie seule » explique tout — un accident résulte presque toujours de la combinaison de causes naturelles et humaines.
\end{attention}

\begin{savaistu}[Les mouvements de terrain au Cameroun]
La région de l'Ouest, avec ses reliefs accidentés, ses sols argileux issus de l'altération des roches volcaniques et ses pluies abondantes, est l'une des zones du Cameroun les plus exposées aux glissements de terrain. Après le drame de Bafoussam en 2019, les autorités ont fait évacuer plusieurs quartiers construits sur des pentes raides et lancé des campagnes de sensibilisation : reconnaître les fissures dans les murs et le sol, c'est parfois pouvoir évacuer à temps.
\end{savaistu}

\courssub{Bilan de l'activité}

Cette enquête sur la catastrophe de Bafoussam t'a montré qu'un mouvement de terrain n'est jamais une fatalité inexplicable : il résulte de la rencontre entre un \textbf{terrain fragile} (pente raide, altérites argileuses sur socle fracturé), un \textbf{élément déclencheur naturel} (des pluies exceptionnelles, dont l'action mécanique et chimique fragilise le sol) et des \textbf{actions humaines aggravantes} (déboisement, constructions mal placées, absence de drainage). Tu as appris à lire des documents variés, à hiérarchiser des causes et à rédiger une explication argumentée — la démarche complète du scientifique. Surtout, tu as constaté que des mesures simples de prévention (reboisement, drainage, contrôle des constructions, surveillance des fissures) peuvent sauver des vies : comprendre les mouvements de terrain, c'est déjà s'en protéger.

\newpage

\coursec{Méthodes et exercices résolus}

\coursec{Leçon 13 : Accidents en zones montagneuses ou liés aux mouvements de terrain}

\courssub{Les mouvements de terrain : définition et exemples}

\begin{definition}[Mouvement de terrain]
    Un \textbf{mouvement de terrain} est un déplacement spontané, naturel ou induit, d'une masse de roches, de débris ou de sols le long d'un versant, sous l'effet de la gravité. La vitesse varie de quelques millimètres par an (fluage) à plusieurs mètres par seconde (chute de blocs).
\end{definition}

\begin{propriete}[Classification des mouvements de terrain (Critères Varnes simplifiés pour le BEPC)]
    On classe les mouvements selon trois critères combinés :
    \begin{enumerate}
        \item \textbf{Le matériau :} Roche intacte (blocs), Roches meubles (argiles, sables, marnes), Débris hétérogènes, Boue/Argile fluide.
        \item \textbf{La cinématique (vitesse) :}
            \begin{itemize}
                \item \textbf{Très lente} ($<$ 1 cm/an) : \emph{Fluage} (\emph{creep}), \emph{Solifluxion}.
                \item \textbf{Lente} ($<$ 1 m/jour) : \textbf{Glissement} (rotatif ou plan), \textbf{Coulée} lente.
                \item \textbf{Rapide} ($>$ 1 m/min) : \textbf{Coulée de boue/débris}, \textbf{Éboulement}.
                \item \textbf{Très rapide} ($>$ 1 m/s) : \textbf{Chute de blocs/pierres}.
            \end{itemize}
        \item \textbf{La géométrie du déplacement :}
            \begin{itemize}
                \item \textbf{Chute} : Trajectoire balistique (chute libre, rebonds, roulement).
                \item \textbf{Éboulement} : Désagrégation d'une masse rocheuse en blocs multiples.
                \item \textbf{Glissement} : Déplacement sur une surface de rupture nette (plan \emph{plan} ou courbe \emph{rotatif}).
                \item \textbf{Coulée} : Écoulement visqueux/fluide, comportement rhéologique (type Bingham).
                \item \textbf{Affaissement} : Tassement vertical (doline, marnière, tassement minier).
            \end{itemize}
    \end{enumerate}
\end{propriete}

\begin{exemplebox}[Exemples camerounais]
    \begin{itemize}
        \item \textbf{Chute/Éboulement de blocs basaltiques} : Route du Mont Cameroun (RN3), falaise de Dschang (RN6).
        \item \textbf{Glissements rotatifs/plans en argiles latéritiques/marnes} : Versants de Bafoussam, quartiers pentus de Yaoundé (Mvog-Ada, Nkolbisson), marnes de l'Éocène sur RN11 (Douala-Bafoussam).
        \item \textbf{Coulées de débris/boues} : Fond de vallée volcanique (Mont Cameroun), ravines de Douala (érosion régressive), cônes de déjection de l'Adamaoua.
        \item \textbf{Affaissements (Dolines/Karst)} : Région de Ngaoundal (calcaires), anciennes carrières (Yaoundé, Douala).
    \end{itemize}
\end{exemplebox}

\begin{popfigure}
    \begin{tikzpicture}[scale=0.7]
        % Chute
        \begin{scope}[xshift=0cm]
            \draw[popdark, thick] (0,0) -- (0,3) -- (1,3.5) -- (1,0); % Falaise
            \foreach \i in {1,2,3} {
                \draw[poporange, fill=poporange!30] (0.2+0.2*\i, 3.2) rectangle ++(0.3,0.3);
                \draw[->, popred, thick] (0.35+0.2*\i, 3.2) -- ++(0, -1.5);
            }
            \node at (0.5, -0.5) {\textbf{Chute de blocs}};
        \end{scope}
        % Glissement rotatif
        \begin{scope}[xshift=4cm]
            \draw[popdark, thick] (0,0) .. controls (1,0.5) and (2,1) .. (3,0); % Surface topographique
            \draw[popblue, thick, dashed] (0,0) .. controls (1,-0.5) and (2,-1) .. (3,0); % Surface rupture
            \draw[poporange, fill=poporange!20] (0,0) .. controls (1,0.5) and (2,1) .. (3,0) .. controls (2,-0.5) and (1,-0.8) .. (0,0);
            \draw[->, popred, thick] (1.5, 0.5) -- ++(-0.5, -0.3);
            \node at (1.5, -1) {\textbf{Glissement rotatif}};
        \end{scope}
        % Glissement plan
        \begin{scope}[xshift=8.5cm]
            \draw[popdark, thick] (0,0) -- (3,1); % Pente
            \draw[popblue, thick, dashed] (0,-0.5) -- (3,0.5); % Plan rupture
            \draw[poporange, fill=poporange!20] (0,0) -- (3,1) -- (3,0.5) -- (0,-0.5) -- cycle;
            \draw[->, popred, thick] (1.5, 0.5) -- ++(0.5, 0.15);
            \node at (1.5, -1) {\textbf{Glissement plan}};
        \end{scope}
        % Coulée
        \begin{scope}[xshift=13cm]
            \draw[popdark, thick] (0,0) -- (1,2) -- (2,2) -- (3,0) -- cycle; % Vallon
            \draw[popblue!70!popgreen, fill=popblue!20!popgreen!20] (0.5,1.8) .. controls (1,1.5) and (1.5,1) .. (2,0.5) -- (2,0) -- (0.5,0) -- cycle;
            \draw[->, popred, thick] (1.2, 1.2) -- ++(0.5, -0.3);
            \node at (1.5, -1) {\textbf{Coulée de boue/débris}};
        \end{scope}
    \end{tikzpicture}
    \legende{Schémas simplifiés des 4 types majeurs de mouvements de terrain (géométrie et cinématique).}
\end{popfigure}

\courssub{Les déformations des roches : souples et cassantes}

\begin{definition}[Déformation souple (ductile) et cassante (fragile)]
    Sous l'effet des contraintes tectoniques (compression, extension), les roches se déforment :
    \begin{itemize}
        \item \textbf{Déformation souple :} Les roches se plient sans rompre $\rightarrow$ formation de \textbf{plis} (anticlinaux, synclinaux). Les couches restent continues. Favorise la fracturation en charnière (diaclasage) et l'infiltration d'eau.
        \item \textbf{Déformation cassante :} Les roches rompent $\rightarrow$ formation de \textbf{failes} (déplacement relatif des compartiments) et de \textbf{diaclases} (fractures sans rejet). Créent des \textbf{plans de faiblesse} majeurs (plans de rupture potentiels, voies de circulation de l'eau).
    \end{itemize}
\end{definition}

\begin{propriete}[Influence de la structure géologique sur la stabilité des versants (Structuralisme)]
    L'orientation des plans structuraux (stratification, schistosité, faille, diaclase) par rapport à la pente topographique contrôle le type de mouvement possible :
    \begin{center}
        \begin{tabularx}{\linewidth}{|l|X|X|}\hline
        \rowcolor{popblueL} \textbf{Configuration} & \textbf{Description} & \textbf{Risque principal} \\ \hline
        \textbf{Versant \emph{envers} (Dip slope)} & Strates parallèles à la pente, pendage vers l'aval. Pendage $>$ Pente. & \textbf{Glissement plan} le long des plans de stratification (surtout si intercalailles tendres : argiles, marnes, schistes). \\ \hline
        \textbf{Versant \emph{revers} (Escarpé)} & Strates perpendiculaires à la pente (pendage vers l'amont). & \textbf{Éboulement / Basculement (Toppling)} par sous-cavage du pied ou rupture traversante. \\ \hline
        \textbf{Faille / Diaclase jourdant dans la pente} & Plan de fracture intersectant la topographie vers l'aval. & \textbf{Glissement plan} le long du plan de faille/diaclase (plan de rupture "tout fait"). \\ \hline
        \textbf{Diaclasage dense / Roches altérées} & Massif très fracturé, saprolite épaisse. & \textbf{Coulées, Glissements diffuses, Affaissements}. \\ \hline
        \end{tabularx}
    \end{center}
    \textbf{Vocabulaire clé :} \cle{Stratification}, \cle{Pendage} (direction + valeur en degrés), \cle{Jour} (ligne d'intersection strate/surface topographique), \cle{Plan de rupture}.
\end{propriete}

\begin{popfigure}
    \begin{tikzpicture}[scale=0.8]
        % Versant Envers
        \begin{scope}[xshift=0cm]
            \draw[popdark, very thick] (0,0) -- (5,-2.5); % Pente 26°
            \foreach \y in {0,-0.8,-1.6,-2.4} {
                \draw[popblue, thick] (-0.5,\y) -- (5.5,\y-1.5); % Strates pendage 30°
            }
            \draw[->, poporange, thick] (2.5,-0.5) -- ++(0.5,-0.3) node[midway, right] {Pendage 30°};
            \draw[->, poporange, thick] (2.5,0.2) -- ++(0.5,-0.25) node[midway, above] {Pente 26°};
            \node[popblue] at (5.5, -1.5) {Strates};
            \node[popdark] at (2.5, 0.6) {Versant \emph{Envers}};
            \node[align=center, font=\small] at (2.5, -3.5) {Risque : Glissement plan\\sur intercalaire tendre};
        \end{scope}
        % Versant Revers
        \begin{scope}[xshift=7cm]
            \draw[popdark, very thick] (0,0) -- (5,-2.5); % Pente
            \foreach \y in {0,-0.8,-1.6,-2.4} {
                \draw[popblue, thick] (-0.5,\y) -- (5.5,\y+1.5); % Strates pendage -30° (vers amont)
            }
            \draw[->, poporange, thick] (2.5,-0.5) -- ++(-0.5,-0.3) node[midway, left] {Pendage -30°};
            \draw[->, poporange, thick] (2.5,0.2) -- ++(0.5,-0.25) node[midway, above] {Pente 26°};
            \node[popblue] at (-0.5, -1.5) {Strates};
            \node[popdark] at (2.5, 0.6) {Versant \emph{Revers}};
            \node[align=center, font=\small] at (2.5, -3.5) {Risque : Éboulement,\\Basculement (Toppling)};
        \end{scope}
    \end{tikzpicture}
    \legende{Influence de l'orientation des strates (pendage) par rapport à la pente : Versant \emph{envers} (glissement plan) vs Versant \emph{revers} (éboulement/basculement).}
\end{popfigure}

\courssub{Le rôle de l'eau dans les mouvements de terrain}

\begin{propriete}[Action MÉCANIQUE de l'eau (Facteur déclencheur et aggravant immédiat)]
    \begin{enumerate}
        \item \textbf{Saturation et Pression interstitielle ($u$) :} L'eau remplit les pores/fissures. La pression de l'eau ($u$) pousse les grains l'un de l'autre $\rightarrow$ diminue la \textbf{contrainte efficace} ($\sigma' = \sigma - u$) $\rightarrow$ diminue la \textbf{résistance au cisaillement} ($\tau_f = c' + \sigma' \tan \phi'$). Le sol passe de l'état solide à l'état liquide (liquéfaction, déclenchement coulées).
        \item \textbf{Érosion et Sous-cavage :} L'eau courante (ruissellement concentré, rivière en pied de versant) enlève le matériau au pied $\rightarrow$ perte de la \textbf{butée} naturelle $\rightarrow$ rupture en traction en tête de versant.
        \item \textbf{Surcharge hydraulique :} Poids de l'eau stockée dans les fissures ou sur le versant (ruissellement) $\rightarrow$ augmentation de la contrainte de cisaillement motrice ($P_{\parallel}$).
    \end{enumerate}
\end{propriete}

\begin{propriete}[Action CHIMIQUE de l'eau (Facteur de préparation à long terme)]
    \begin{enumerate}
        \item \textbf{Altération chimique (Roche $\rightarrow$ Sol meuble) :} \textbf{Hydrolyse} des silicates (feldspaths, micas) $\rightarrow$ formation d'\textbf{argiles} (kaolinite, smectite). \textbf{Dissolution} (calcaires, gypse). \textbf{Oxydation} (pyrite, minéraux ferreux). Résultat : perte de cohésion, formation de \textbf{saprolite} / \textbf{latérite} instable.
        \item \textbf{Gonflement / Retrait des argiles :} Argiles gonflantes (smectites) : absorption d'eau $\rightarrow$ augmentation de volume $\rightarrow$ fissuration $\rightarrow$ perte de cohésion ($c' \downarrow$). Retrait au séchage $\rightarrow$ fissures de dessiccation = voies d'infiltration préférentielles.
        \item \textbf{Lessivage / Piping :} Élimination des ciments (carbonates, oxydes de fer) liant les grains $\rightarrow$ effondrement de la structure (suffosion, dolines, cheminées d'érosion).
    \end{enumerate}
\end{propriete}

\begin{methode}[Analyser le rôle de l'eau dans un accident]
    \begin{enumerate}
        \item \textbf{Identifier la source :} Pluie intense/longue (infiltration), crue rivière (érosion pied), fuite réseau (recharge artificielle), variation nappe.
        \item \textbf{Séparer Mécanique / Chimique :}
            \begin{itemize}
                \item \textit{Mécanique (court terme) :} Saturation $\rightarrow$ Pression interstitielle $u \uparrow$ $\rightarrow$ Frottement $\downarrow$ $\rightarrow$ Rupture.
                \item \textit{Chimique (long terme) :} Altération $\rightarrow$ Cohésion $c' \downarrow$ $\rightarrow$ Matériau pré-fragilisé.
            \end{itemize}
        \item \textbf{Quantifier si possible :} Seuil pluviométrique (ex. 100 mm/24h, 200 mm cumulés), niveau piézométrique critique.
    \end{enumerate}
    \textbf{Distinction BEPC :} \textbf{Déclencheur} = événement brutal (orage, séisme) sur versant \emph{préparé}. \textbf{Préparation} = altération, fracturation, géométrie défavorable. L'eau est souvent les deux.
\end{methode}

\courssub{Le rôle de l'action de l'homme}

\begin{propriete}[Activités anthropiques déstabilisantes et mécanismes]
    \begin{center}
        \begin{tabularx}{\linewidth}{|l|X|X|}\hline
        \rowcolor{popblueL} \textbf{Action humaine} & \textbf{Mécanisme de déstabilisation} & \textbf{Exemple Cameroun} \\ \hline
        \textbf{Déforestation / Déboisement} & Perte \textbf{renforcement racinaire} ($c' \downarrow$), perte \textbf{évapotranspiration} (sol + humide), érosion superficielle (battance). & Pentes de Bafoussam, Yaoundé, Mont Cameroun. \\ \hline
        \textbf{Agriculture sur brûlis / Labour sens pente} & Destruction structure sol, création sillons préférentiels ruissellement, perte cohésion. & Hauts-Plateaux (Bamileke), Adamaoua. \\ \hline
        \textbf{Urbanisation anarchique (Pied/Tête versant)} & \textbf{Surcharge} en tête (poids $\uparrow$), \textbf{Suppression butée} en pied (excavation), \textbf{Imperméabilisation} (ruissellement concentré). & Yaoundé (quartiers non lotis), Douala (ravines). \\ \hline
        \textbf{Terrassements / Routes / Barrages} & Création pentes artificielles raides ($>$ angle frottement), remblais mal compactés, déblais pied. & RN6 (Dschang), RN11 (Marnes), Barrage Lom Pangar. \\ \hline
        \textbf{Fuites réseaux (Eau/Égouts/Fosses septiques)} & \textbf{Recharge artificielle chronique} nappe perchée $\rightarrow$ $u \uparrow$ permanent. & Yaoundé (fosses septiques non étanches), Douala. \\ \hline
        \textbf{Exploitation minière / Carrières} & Création vides, surcharges stériles, vibrations (tassement dynamique/liquéfaction). & Carrières granit (Yaoundé), Mines (Est, Nord). \\ \hline
    \end{tabularx}
\end{center}
\end{propriete}

\courssub{Prévention et protection contre les mouvements de terrain}

\begin{definition}[Prévention vs Protection]
    \begin{itemize}
        \item \textbf{Prévention (Réduit l'ALÉA) :} Agit sur la \textbf{CAUSE} pour diminuer la probabilité/intensité du phénomène. Ex : Drainage (baisse $u$), Reprofilage (allège tête/charge pied), Revégétalisation (racines, évapotranspiration), Règlementation (POS, PPRI - Zone Rouge inconstructible).
        \item \textbf{Protection (Réduit la VULNÉRABILITÉ/ENJEUX) :} Agit sur les \textbf{CONSEQUENCES} pour limiter les dégâts. Ex : Ouvrages de retenue (filets, merlons, bassins), Déviation trajet, Surveillance/Alerte/Évacuation (PCS), Adaptation bâti (pilotis).
    \end{itemize}
\end{definition}

\begin{propriete}[Tableau de choix des techniques (Mémo BEPC)]
    \begin{center}
        \begin{tabularx}{\linewidth}{|l|X|X|}\hline
        \rowcolor{popblueL} \textbf{Type de mouvement} & \textbf{Prévention (Cause / Aléa)} & \textbf{Protection (Effet / Enjeux)} \\ \hline
        \textbf{Chute / Éboulement} (Blocs, Rapide) & Purge manuelle, Grillage plaqué ancré, Ancrage actif (tirants), Reboisement pied falaise. & \textbf{Filets pare-blocs} (haute énergie), \textbf{Merlons / Bassins de réception}, Galeries pare-blocs, Déviation trajet. \\ \hline
        \textbf{Glissement Plan / Rotatif} (Lent, Argiles/Marnes/Roches tendres) & \textbf{Drainage profond} (Drains horizontaux, Puits de drainage), \textbf{Reprofilage} (Alléger tête, Charger pied/butée), \textbf{Ancrage} (Tirants, Micropieux, Clous), Revégétalisation (Vetiver, Bambous). & Murs de soutènement (poids, console, berlinoise), Palplanches, \textbf{Surveillance} (Inclinomètres, Piézomètres, Topographie, Radar satellite), Alertes. \\ \hline
        \textbf{Coulée Boueuse / Débris} (Rapide, Chenal) & Reboisement bassin versant, \textbf{Seuils de torrent} (freiner vitesse, piéger solide), Bassins de décantation amont. & \textbf{Bassins de rétention / Dépôt} (exutoire), Digues de déviation, \textbf{Système d'alerte crues / Évacuation} (sirènes, SMS). \\ \hline
        \textbf{Affaissement / Doline} (Karst, Minier) & Comblement cavités (injection coulis), Contrôle pompages nappe, Remblaiement contrôlé. & Surveillance tassement (extensomètres), Renforcement fondations (pieux, micropieux), Évacuation zone. \\ \hline
        \end{tabularx}
    \end{center}
    \textbf{Critère de choix :} \textbf{Coût / Efficacité / Délai / Acceptabilité / Durabilité.} Ex : Filet = rapide/coûteux/visuel. Drainage = lent/durable/invisible. Reprofilage = lourd/définitif.
\end{propriete}

\begin{savaistu}[Le Vetiver (Chrysopogon zizanioides) : Allié végétal au Cameroun]
    Le \textbf{Vetiver} est une graminée tropicale à système racinaire \textbf{fasciculé, profond (3 à 4 m) et dense}. Utilisé en \textbf{génie végétal} (haies en lignes de niveau) :
    \begin{itemize}
        \item \textbf{Renforcement mécanique :} Racines = "clous" naturels traversant le plan de rupture superficiel $\rightarrow$ augmentent cohésion $c'$ de 5 à 15 kPa.
        \item \textbf{Régulation hydrologique :} Forte évapotranspiration $\rightarrow$ assèche le sol $\rightarrow$ baisse pression interstitielle $u$.
        \item \textbf{Piégeage sédimentaire :} Haies denses filtrent ruissellement, créent terrasses naturelles.
    \end{itemize}
    Succès documentés sur RN11 (marnes), pentes de Bafoussam, berges de rivières. Peu consommateur d'eau une fois établi, non invasif (stérile).
\end{savaistu}

\courssub{Synthèse et démarche d'analyse d'un accident}

\begin{methode}[Rédiger une analyse de risque complète (Scénario $\rightarrow$ Aléa $\rightarrow$ Enjeux $\rightarrow$ Risque $\rightarrow$ Mesures)]
    \textbf{Objectif :} Structurer une réponse argumentée type "synthèse d'unité" ou "analyse de document" pour le BEPC.
    \begin{enumerate}
        \item \textbf{Décrire le SITE (Contexte géologique/géomorphologique) :} Nature des roches, structure (pendage, failles), morphologie (pente, forme), hydrogéologie (nappe, sources). $\rightarrow$ \textbf{Susceptibilité} (intrinsèque).
        \item \textbf{Identifier les FACTEURS DÉCLENCHEURS :} Pluie (seuils), séisme, action humaine (travaux, fuite), érosion pied. $\rightarrow$ \textbf{Aléa} (probabilité d'occurrence dans un délai donné + intensité/vitesse).
        \item \textbf{Inventorier les ENJEUX (Éléments exposés) :} Population (nombre, vulnérabilité), habitats, infrastructures (routes, écoles, hôpitaux), activités économiques, patrimoine. Localiser précisément (pied, tête, corps du glissement).
        \item \textbf{Qualifier le RISQUE :} Risque = Aléa $\times$ Enjeux $\times$ Vulnérabilité. Niveau : \textbf{Fort / Modéré / Faible / Nul}. Justifier.
        \item \textbf{Proposer une STRATÉGIE GLOBALE (Gestion intégrée) :}
            \begin{itemize}
                \item \textbf{Connaissance :} Cartographie aléa (Zonage), surveillance (instrumentation).
                \item \textbf{Prévention (Réduction aléa) :} Drainage, reprofilage, revégétalisation, règlementation (POS, PPRI).
                \item \textbf{Protection (Réduction vulnérabilité) :} Ouvrages (filets, merlons), adaptation bâti.
                \item \textbf{Gestion de crise :} Plan Communal de Sauvegarde (PCS), exercices évacuation, système alerte, secours.
            \end{itemize}
    \end{enumerate}
    \textbf{Conseil rédaction :} Utiliser connecteurs logiques : \emph{"Du fait de..."}, \emph{"Ce qui entraîne..."}, \emph{"Par conséquent..."}, \emph{"Mesure adaptée car..."}. Chiffrer (volume, vitesse, pluviométrie, population). Citer vocabulaire précis (\cle{pression interstitielle}, \cle{contrainte efficace}, \cle{versant envers}, \cle{plan de rupture}).
\end{methode}

\begin{methode}[Interpréter une carte d'aléa / de zonage géotechnique (PPRI type)]
    \textbf{Objectif :} Lire une carte réglementaire (légende, zonage) pour prendre une décision d'aménagement.
    \begin{enumerate}
        \item \textbf{Analyser la LÉGENDE :} Codes couleurs/hachures.
            \begin{itemize}
                \item \textbf{Zone Rouge (R) / Aléa Fort :} Interdiction construction (sauf exceptions drastiques). Ex : axe principal glissement, chute blocs active, lit torrent.
                \item \textbf{Zone Bleue (B) / Aléa Modéré :} Construction soumise prescriptions (Étude G2 obligatoire, mesures imposées : drainage, pieux, interdiction sous-sol/fosse septique/remblai $>$50cm).
                \item \textbf{Zone Blanche / Aléa Faible/Nul :} Pas de prescription spécifique aléa MT.
            \end{itemize}
        \item \textbf{Localiser le PROJET :} Situer emprise par rapport aux limites.
        \item \textbf{Appliquer le RÈGLEMENT :} Lire règlement joint. Vérifier constructibilité, obligations, interdictions.
        \item \textbf{Conclure :} "Projet en zone [Couleur] $\rightarrow$ [Autorisé / Autorisé sous conditions / Interdit]. Mesures imposées : ..."
    \end{enumerate}
    \textbf{Piège BEPC :} Ne pas confondre \textbf{Carte d'Aléa} (phénomène physique seul) et \textbf{Carte de Zonage / PPRI} (réglementaire, intègre enjeux). L'énoncé précise laquelle est fournie.
\end{methode}

\courssub{Méthodes et exercices résolus}

\begin{methode}[Identifier et classifier les types de mouvements de terrain]
    \textbf{Objectif :} Distinguer les principaux types de mouvements de terrain selon leur vitesse, le matériau mis en jeu et la géométrie du déplacement.
    \begin{enumerate}
        \item \textbf{Observer le matériau :} roche intacte (blocs), roches meubles (argiles, sables, marnes), débris hétérogènes, boue.
        \item \textbf{Analyser la cinématique (vitesse) :} très lente (fluage, \emph{creep} $<$ 1 cm/an), lente (glissement, coulée $<$ 1 m/jour), rapide (éboulement, coulée de boue $>$ 1 m/min), très rapide (chute de blocs $>$ 1 m/s).
        \item \textbf{Repérer la géométrie du mouvement :} rotation (glissement rotatif), translation (glissement plan), écoulement (coulée), chute (chute de blocs, éboulement), affaissement (tassement, doline).
        \item \textbf{Nommer le type :} associer matériau + cinématique + géométrie (ex. : \emph{glissement rotatif d'argiles}, \emph{coulée de boue}, \emph{chute de blocs calcaires}).
    \end{enumerate}
    \textbf{Repères pour le BEPC :} Connaître les 5 grands types : \textbf{chute}, \textbf{éboulement}, \textbf{glissement}, \textbf{coulée}, \textbf{affaissement}. Savoir donner un exemple local (ex. : éboulements sur la route du Mont Cameroun, glissements de terrain à Bafoussam ou Yaoundé).
\end{methode}

\begin{exoresolu}[Classification de mouvements de terrain observés au Cameroun]
    \textbf{Énoncé :} Lors d'une sortie géologique dans la région de l'Ouest, trois phénomènes sont observés :
    \begin{enumerate}
        \item Sur une pente raide en bordure de route (RN6), des blocs de basalte de 0,5 à 2 m$^3$ se détachent brutalement de la falaise et roulent jusqu'au bas du talus. Vitesse estimée $>$ 50 km/h.
        \item Sur une pente douce (15\%) recouverte d'argiles latéritiques saturées d'eau après de fortes pluies, une masse de terrain de 200 m de long glisse sur une surface de rupture concave (en cuillère), en conservant sa structure interne. Vitesse : 2 cm/jour.
        \item Dans un fond de vallée alluvionnaire, après un orage violent, un mélange eau + sable + graviers + débris végétaux s'écoule rapidement dans le lit d'un ruisseau, suivant la topographie. Vitesse : 5 m/min.
    \end{enumerate}
    Pour chaque phénomène, précise : (a) le type de mouvement, (b) le matériau, (c) la cinématique, (d) la géométrie.
    \tcblower
    \textbf{Solution détaillée :}
    \begin{enumerate}
        \item \textbf{Phénomène 1 :}
            \begin{itemize}
                \item (a) \textbf{Chute de blocs} (volume $>$ 1 m$^3$).
                \item (b) Matériau : \textbf{roche dure intacte} (basalte).
                \item (c) Cinématique : \textbf{très rapide} ($>$ 1 m/s).
                \item (d) Géométrie : \textbf{chute libre} puis roulement/rebond (trajectoire balistique).
            \end{itemize}
            \textit{Justification :} Détachement brutal de blocs individuels sur falaise, vitesse extrême, pas de plan de rupture continu dans la masse.
        \item \textbf{Phénomène 2 :}
            \begin{itemize}
                \item (a) \textbf{Glissement rotatif}.
                \item (b) Matériau : \textbf{roches meubles} (argiles latéritiques).
                \item (c) Cinématique : \textbf{lente} (cm/jour).
                \item (d) Géométrie : \textbf{rotation} autour d'un axe horizontal (surface de rupture concave, "en cuillère"), la masse bascule vers l'arrière.
            \end{itemize}
            \textit{Justification :} Surface de rupture concave (pas plane), conservation de la structure, saturation en eau (facteur déclencheur), vitesse lente.
        \item \textbf{Phénomène 3 :}
            \begin{itemize}
                \item (a) \textbf{Coulée de débris} (\emph{debris flow}) ou \textbf{coulée boueuse} si matrice argileuse dominante.
                \item (b) Matériau : \textbf{débris hétérogènes} (sable, graviers, blocs, boue, bois) en suspension dans l'eau.
                \item (c) Cinématique : \textbf{rapide} (m/min).
                \item (d) Géométrie : \textbf{écoulement} chenalisé (suit le thalweg), comportement fluide (rhéologie de Bingham).
            \end{itemize}
            \textit{Justification :} Mélange eau/sédiments, écoulement rapide dans un chenal, déclenché par pluie intense (ruissellement concentré).
    \end{enumerate}
    \textbf{Conclusion :} L'identification précise du type de mouvement conditionne le choix des mesures de protection (filets pare-blocs pour la chute, drainage pour le glissement, seuils/bassins de rétention pour la coulée).
\end{exoresolu}

\begin{methode}[Relier déformations des roches (souples/cassantes) et risque de mouvement de terrain]
    \textbf{Objectif :} Expliquer comment la structure géologique (plis, failles, diaclases) issue des déformations contrôle la stabilité des versants.
    \begin{enumerate}
        \item \textbf{Identifier le type de déformation :}
            \begin{itemize}
                \item \textbf{Souple (ductile) :} les roches se plient sans rompre $\rightarrow$ \textbf{plis} (anticlinaux, synclinaux). Favorise l'infiltration de l'eau dans les charnières fracturées.
                \item \textbf{Cassante (fragile) :} les roches rompent $\rightarrow$ \textbf{failes} (déplacement relatif des compartiments) et \textbf{diaclases} (fractures sans déplacement). Créent des plans de faiblesse, des plans de rupture potentiels et des voies de circulation pour l'eau.
            \end{itemize}
        \item \textbf{Analyser l'orientation par rapport à la pente (structuralisme) :}
            \begin{itemize}
                \item \textbf{Pente \emph{envers} (versant dip slope) :} strates parallèles à la pente, trempage vers l'aval. \textbf{Risque fort} de glissement plan le long des plans de stratification (surtout si couches tendres intercalées : argiles, marnes).
                \item \textbf{Pente \emph{revers} (versant escarpé) :} strates perpendiculaires à la pente. Risque d'éboulement/basculement par sous-cavage ou basculement de colonnes (toppling).
                \item \textbf{Faille/Diaclase jourdant dans la pente :} plan de rupture prêt à l'emploi pour un glissement.
            \end{itemize}
        \item \textbf{Conclure sur le risque :} Une zone très fracturée (diaclasage dense) ou une faille active est un site privilégié pour l'initiation de mouvements de terrain.
    \end{enumerate}
    \textbf{Vocabulaire clé :} \cle{stratification}, \cle{pendage} (direction + valeur), \cle{jour} (intersection strate/topographie), \cle{plan de rupture}.
\end{methode}

\begin{exoresolu}[Analyse structurale d'un versant instable]
    \textbf{Énoncé :} On observe un versant constitué d'une série sédimentaire alternant couches de grès (dures, perméables) et couches d'argile (tendres, imperméables). Les strates ont un pendage de 30° vers le Nord. La pente topographique fait face au Nord (versant \emph{envers}), avec une inclinaison de 25°.
    \begin{enumerate}
        \item Représenter schématiquement la coupe géologique (strates, pente, pendage).
        \item Quel type de mouvement de terrain est le plus probable ? Justifier par la géométrie.
        \item Quel rôle joue l'eau de pluie dans ce contexte structural ?
        \item Proposer une mesure de protection adaptée à la cause structurelle.
    \end{enumerate}
    \tcblower
    \textbf{Solution détaillée :}
    \begin{enumerate}
        \item \textbf{Schéma (description) :} Tracer une ligne de pente inclinée à 25° vers la droite (Nord). Tracer des lignes parallèles inclinées à 30° vers la droite (strates). Les strates "plongent" sous la pente. Alterner traits épais (grès) et traits hachurés (argile).
        \begin{popfigure}
            \begin{tikzpicture}[scale=0.8]
                % Topography
                \draw[popdark, very thick] (0,0) -- (10,-4.66); % slope 25 deg
                \draw[popdark, thick] (0,0) -- (10,0); % horizontal top
                \draw[popdark, thick] (10,-4.66) -- (10,-6); % vertical base
                \draw[popdark, thick] (0,0) -- (0,-6); % vertical back
                % Strata (dip 30 deg)
                \foreach \y in {0,-1,-2,-3,-4,-5} {
                    \draw[popblue, thick] (0,\y) -- (11.55,\y-6.66); % dip 30 deg
                }
                % Labels
                \node[popblue] at (11, -2) {Strates (30°)};
                \node[popdark] at (5, 0.5) {Surface topographique (25°)};
                \draw[->, poporange, thick] (5,-1) -- (5,-2.5) node[midway, right] {Pendage};
                \draw[->, poporange, thick] (5,-0.5) -- (6.5,-1.5) node[midway, above] {Pente};
            \end{tikzpicture}
            \legende{Coupe schématique : Versant \emph{envers} (strates et pente même sens, pendage $>$ pente).}
        \end{popfigure}
        \item \textbf{Glissement plan (ou glissement translatif) le long des couches d'argile.}
            \textit{Justification :} C'est une configuration classique de \textbf{versant \emph{envers}} (\emph{dip slope}). Les strates jourdent dans le sens de la pente. Le pendage (30°) est supérieur à la pente (25°), ce qui signifie que les plans de stratification affleurent en pente. Les couches d'argile imperméables agissent comme des plans de glissement potentiels (faible cohésion, lubrification par l'eau). La masse de grès au-dessus peut glisser sur l'argile.
        \item \textbf{Rôle de l'eau :}
            \begin{itemize}
                \item \textit{Mécanique :} L'eau s'infiltre dans les grès (perméables) et s'accumule sur les argiles (imperméables) $\rightarrow$ augmentation de la \textbf{pression interstitielle} ($u$) $\rightarrow$ diminution de la contrainte efficace ($\sigma' = \sigma - u$) $\rightarrow$ diminution du frottement de Coulomb ($\tau = c' + \sigma' \tan \phi'$) le long du plan d'argile.
                \item \textit{Chimique :} Gonflement des argiles (smectites) au contact de l'eau $\rightarrow$ perte de résistance au cisaillement.
            \end{itemize}
        \item \textbf{Mesure de protection :} \textbf{Drainage profond} (drains horizontaux forés dans le sens de la pente pour abaisser la nappe perchée sur l'argile) + \textbf{Ancrage} (tirants d'ancrage scellés dans le grès sain sous la couche d'argile pour "clouer" le versant). Éviter le surcharge en tête de versant.
    \end{enumerate}
    \textbf{Conclusion :} La concordance entre pendage des strates et pente topographique (versant \emph{envers}) avec des intercalailles tendres est un critère majeur de susceptibilité aux glissements plans.
\end{exoresolu}

\begin{methode}[Analyser le rôle déclencheur et aggravant de l'eau]
    \textbf{Objectif :} Expliquer les mécanismes physiques et chimiques par lesquels l'eau (pluie, nappe, cours d'eau) initie ou aggrave l'instabilité.
    \begin{enumerate}
        \item \textbf{Identifier la source d'eau :} pluviométrie intense/longue (infiltration), crue de rivière (érosion du pied / mise en charge), rupture de canalisation, fonte des neiges (peu pertinent en zone tropicale hors très haute altitude), variation niveau lac/nappe.
        \item \textbf{Décrire l'action \textbf{MÉCANIQUE} :}
            \begin{itemize}
                \item \textbf{Saturation / Pression interstitielle ($u$) :} L'eau remplit les pores $\rightarrow$ pousse les grains l'un de l'autre $\rightarrow$ annule les contraintes effectives $\rightarrow$ le sol passe de l'état solide à l'état liquide (liquéfaction, coulée).
                \item \textbf{Érosion / Sous-cavage :} L'eau courante (ruissellement, rivière) enlève le matériau au pied du versant $\rightarrow$ perte de butée $\rightarrow$ rupture en traction en haut du versant.
                \item \textbf{Surcharge :} Poids de l'eau stockée dans les fissures / sur le versant.
            \end{itemize}
        \item \textbf{Décrire l'action \textbf{CHIMIQUE} :}
            \begin{itemize}
                \item \textbf{Altération :} Dissolution (calcaire, gypse), hydrolyse (feldspaths $\rightarrow$ argiles), oxydation. Transformant la roche saine en matériau meuble instable (saprolite, latérite).
                \item \textbf{Gonflement / Retrait :} Argiles gonflantes (smectites) : augmentation volume $\rightarrow$ fissuration $\rightarrow$ perte cohésion.
                \item \textbf{Lessivage :} Élimination des ciments (carbonates, oxydes de fer) liant les grains $\rightarrow$ effondrement structure (piping, dolines).
            \end{itemize}
        \item \textbf{Quantifier si possible :} Seuil de pluviométrie (ex. : 100 mm en 24h, ou 200 mm cumulés), niveau piézométrique critique.
    \end{enumerate}
    \textbf{Distinction BEPC :} \textbf{Déclencheur} = événement brutal (orage, séisme) sur un versant \emph{préparé}. \textbf{Facteur de préparation} = altération chimique longue, fracturation, géométrie défavorable. L'eau est souvent les deux.
\end{methode}

\begin{exoresolu}[Rôle de l'eau dans un glissement de terrain à Bafoussam (Hauts-Plateaux)]
    \textbf{Énoncé :} En septembre 2023, après trois semaines de pluies continues (cumul 450 mm), un glissement de terrain se produit sur un versant en latérite argileuse (épaisseur 15 m) reposant sur un socle granitique frais. Le glissement mobilise 50 000 m$^3$. Une citerne d'eau fuyait en haut du versant depuis 6 mois.
    \begin{enumerate}
        \item Distinguer dans ce scénario les facteurs de \textbf{préparation} et le(s) facteur(s) \textbf{déclencheur(s)}.
        \item Expliquer le mécanisme mécanique de la rupture par l'eau (pression interstitielle). Faire un schéma de forces sur un élément de sol sur le plan de rupture.
        \item Expliquer le rôle chimique de l'eau sur le long terme (altération du granite).
        \item Pourquoi la fuite de la citerne est-elle un facteur aggravant majeur ?
    \end{enumerate}
    \tcblower
    \textbf{Solution détaillée :}
    \begin{enumerate}
        \item \textbf{Facteurs de préparation (long terme) :}
            \begin{itemize}
                \item Nature du matériau : \textbf{latérite argileuse} (faible perméabilité, forte plasticité, sensibilité à l'eau).
                \item Structure : \textbf{contact latérite/granite frais} = surface de rupture potentielle plane (pédiment).
                \item Altération chimique ancienne : transformation du granite en latérite (perte de cohésion).
                \item Fuite citerne (6 mois) : \textbf{recharge artificielle chronique} de la nappe perchée, maintien d'une humidité élevée, gonflement des argiles.
            \end{itemize}
            \textbf{Facteur déclencheur (court terme) :} \textbf{Pluies exceptionnelles} (450 mm en 3 semaines) $\rightarrow$ saturation rapide de la colonne de latérite, montée brutale de la nappe perchée jusqu'à la surface de rupture.
        \item \textbf{Mécanisme mécanique (Schéma de forces) :}
            \begin{popfigure}
                \begin{tikzpicture}[scale=1.2]
                    % Plan de rupture incliné
                    \draw[popdark, very thick] (0,0) -- (6,1.5);
                    % Bloc de sol
                    \draw[poporange!70!popdark, thick, fill=poporange!20] (0,0) -- (6,1.5) -- (6,3) -- (0,1.5) -- cycle;
                    % Forces
                    \coordinate (G) at (3,2.25);
                    \draw[->, popblue, very thick] (G) -- (3,0.5) node[midway, right] {$P$ (Poids)};
                    \draw[->, popgreen, very thick] (G) -- (4.5,3.5) node[midway, right] {$R$ (Résistance frottement)};
                    \draw[->, popred, very thick] (G) -- (1.5,3.5) node[midway, left] {$F_u$ (Poussée eau / Pression interstitielle)};
                    % Decomposition de P
                    \draw[dashed, popblue] (G) -- (3,0.5);
                    \draw[dashed, popblue] (G) -- (5.5,2.25);
                    \draw[->, popblue] (G) -- (5.5,2.25) node[midway, above] {$P_{\parallel}$ (Moteur)};
                    \draw[->, popblue] (G) -- (3,0.5) node[midway, right] {$P_{\perp}$ (Normal)};
                    % Angle
                    \draw (5.5,1.5) arc (0:14:0.5) node[midway, right, scale=0.8] {$\alpha$};
                \end{tikzpicture}
                \legende{Forces sur une tranche de sol : $P_{\parallel}$ entraine le glissement, $R = (P_{\perp} - u A) \tan \phi$ s'oppose. L'eau ($u$) diminue $R$.}
            \end{popfigure}
            \textit{Explication :} Le poids $P$ se décompose en $P_{\parallel}$ (moteur, parallèlement à la pente) et $P_{\perp}$ (normal, créant le frottement). L'eau dans les pores exerce une pression $u$ (force $F_u = u \times A$ vers le haut, normale au plan). La contrainte efficace normale devient $\sigma' = \frac{P_{\perp}}{A} - u$. La résistance au cisaillement $\tau_f = c' + \sigma' \tan \phi'$ diminue. Si $\tau_f < \frac{P_{\parallel}}{A}$, rupture.
        \item \textbf{Rôle chimique (long terme) :} \textbf{Hydrolyse} des feldspaths et micas du granite $\rightarrow$ formation de \textbf{kaolinite/smectite} (argiles). \textbf{Lessivage} de la silice et des bases (Ca, Mg, K, Na). Perte de la structure rocheuse cohérente $\rightarrow$ formation de la latérite meuble (saprolite). C'est la "préparation" géologique du désastre.
        \item \textbf{Fuite citerne :} Elle a maintenu la nappe perchée à un niveau anormalement haut \textbf{avant} les grosses pluies. Le versant était "pré-chargé" en eau. Les pluies n'ont eu qu'à combler le faible déficit résiduel pour atteindre le seuil critique de pression interstitielle. C'est une cause \textbf{anthropique} directe.
    \end{enumerate}
    \textbf{Conclusion :} L'accident résulte de la convergence d'une préparation géologique (altération), d'une aggravation anthropique (fuite) et d'un déclencheur météorologique (pluie extrême).
\end{exoresolu}

\begin{methode}[Évaluer l'impact de l'action de l'homme sur la stabilité des versants]
    \textbf{Objectif :} Identifier les activités humaines qui déstabilisent les pentes et expliquer le mécanisme (modification de la géométrie, de l'hydrologie, de la charge, de la végétation).
    \begin{enumerate}
        \item \textbf{Lister l'action :} Déforestation / déboisement, agriculture sur brûlis / labour dans le sens de la pente, urbanisation anarchique (construction en pied ou tête de versant), terrassements / remblais / déblais (routes, barrages), exploitation minière / carrières, fuites réseaux (eau, égouts), surpâturage.
        \item \textbf{Analyser le mécanisme de déstabilisation :}
            \begin{itemize}
                \item \textbf{Végétation :} Perte du \textbf{renforcement racinaire} (racines = "clous" naturels traversant le plan de rupture), perte de l'\textbf{évapotranspiration} (sol plus humide), perte de protection contre l'érosion superficielle (ruissellement).
                \item \textbf{Géométrie :} \textbf{Déblai} en pied (suppression butée) $\rightarrow$ rupture. \textbf{Remblai} en tête (surcharge) $\rightarrow$ augmentation contrainte de cisaillement. \textbf{Terrassement} créant des pentes artificielles trop raides ($>$ angle de frottement).
                \item \textbf{Hydrologie :} Imperméabilisation (urbanisation) $\rightarrow$ ruissellement concentré $\rightarrow$ érosion/gouffres. Fuites réseaux $\rightarrow$ recharge nappe. Drains mal conçus $\rightarrow$ saturation locale.
                \item \textbf{Vibrations :} Trafic lourd, carrières, travaux publics $\rightarrow$ tassement dynamique / liquéfaction sols meubles saturés.
            \end{itemize}
        \item \textbf{Quantifier l'aggravation :} Comparer facteur de sécurité $F_s$ avant/après travaux (si données numériques). Sinon, argumenter qualitativement.
    \end{enumerate}
    \textbf{Exemples Cameroun :} Glissements à Yaoundé (quartiers non lotis, pentes raides, fosses septiques), éboulements route Mont Cameroun (défrichement, vibrations), érosion régressive Douala (drainage défaillant).
\end{methode}

\begin{exoresolu}[Étude de cas : Urbanisation et glissements à Yaoundé]
    \textbf{Énoncé :} Dans un quartier non loti en pente forte (35°) de Yaoundé, construit sur des schistes micasteux altérés en argiles sableuses, on observe : (1) des maisons construites en pied de versant sur remblais, (2) des fosses septiques non étanches au-dessus des habitations, (3) absence de réseau d'évacuation des eaux pluviales, (4) déboisement total. Lors d'une forte pluie, un glissement détruit 3 maisons en pied de versant.
    \begin{enumerate}
        \item Pour chacune des 4 observations, expliquer le mécanisme de déstabilisation (géométrique, hydraulique, mécanique).
        \item Classer ces actions en "modification de la charge", "modification de l'hydrologie", "modification de la résistance du sol".
        \item Proposer 3 mesures d'urbanisme préventives applicables \emph{avant} construction.
    \end{enumerate}
    \tcblower
    \textbf{Solution détaillée :}
    \begin{enumerate}
        \item \textbf{Analyse des 4 facteurs anthropiques :}
            \begin{enumerate}
                \item \textbf{Maisons en pied sur remblais :} \textbf{Surcharge en pied} (butée) \textbf{MAS} \textbf{suppression de la butée naturelle} si excavation pour fondation. Le remblai, souvent mal compacté, est un matériau faible qui s'écrase ou glisse. \textit{Mécanisme : Géométrique (suppression butée) + Charge (surcharge mal maîtrisée).}
                \item \textbf{Fosses septiques fuyantes en tête :} \textbf{Recharge artificielle massive} de la nappe perchée. L'eau s'infiltre verticalement, rejoint le plan de rupture (contact altérites/roche saine) $\rightarrow$ hausse pression interstitielle $u$ $\rightarrow$ chute résistance. \textit{Mécanisme : Hydraulique (augmentation $u$).}
                \item \textbf{Absence drainage eaux pluviales :} \textbf{Ruissellement concentré} sur pente nue $\rightarrow$ \textbf{érosion régressive} (ravines) qui entament le versant, créent des surplombs qui s'effondrent. Infiltration massive en tête de versant. \textit{Mécanisme : Hydraulique (infiltration + érosion) + Géométrique (création ravines).}
                \item \textbf{Déboisement total :} Perte \textbf{renforcement racinaire} (cohesion $c'$ diminuée). Perte \textbf{évapotranspiration} (sol reste saturé plus longtemps). \textbf{Erosion de surface} (croûtes de battance, ruissellement). \textit{Mécanisme : Résistance (perte $c'$) + Hydraulique (saturation).}
            \end{enumerate}
        \item \textbf{Classification :}
            \begin{center}
                \begin{tabularx}{\linewidth}{|l|X|}\hline
                \rowcolor{popblueL} \textbf{Type de modification} & \textbf{Actions concernées} \\ \hline
                Charge (Géométrie/Poids) & (1) Remblais en pied / Excavations \\ \hline
                Hydrologie (Eau) & (2) Fosses septiques, (3) Absence drainage, (4) Déboisement (évapotranspiration) \\ \hline
                Résistance du sol (Mécanique) & (4) Déboisement (perte racines), (3) Érosion (affaiblissement structure) \\ \hline
                \end{tabularx}
            \end{center}
        \item \textbf{Mesures préventives (Aménagement du territoire) :}
            \begin{enumerate}
                \item \textbf{Zonage / Plan d'Occupation des Sols (POS) :} Interdire la construction dans les zones à forte pente ($>$ 30°) et en pied de versant instable (zone rouge). Délimiter des zones "non aedificandi".
                \item \textbf{Assainissement obligatoire :} Imposer le raccordement à un réseau d'égouts ou des fosses septiques \textbf{étanches} avec épandage en zone stable (hors versant), + réseau d'évacuation des eaux pluviales (caniveaux, collecteurs) dimensionné pour les pluies centennales.
                \item \textbf{Gestion de la végétation / Génie végétal :} Interdire le déboisement total. Imposer la conservation de bandes boisées (rideaux brise-vent, stabilisation rives). Utiliser des techniques de \textbf{génie végétal} (fascines, plantations de vétiver, bambous) pour renforcer le sol par les racines.
            \end{enumerate}
    \end{enumerate}
    \textbf{Conclusion :} L'accident n'est pas "naturel", il est \textbf{construit} par l'accumulation de mauvaises pratiques d'occupation du sol. La prévention passe par la planification urbaine (POS) et l'assainissement, bien plus efficaces que les murs de soutènement curatifs.
\end{exoresolu}

\begin{methode}[Proposer des techniques de prévention (aléa) et de protection (enjeux) adaptées]
    \textbf{Objectif :} Choisir la technique adaptée au type de mouvement, au matériau, à la vitesse et au contexte (coût, urgence, enjeux).
    \begin{enumerate}
        \item \textbf{Distinction fondamentale :}
            \begin{itemize}
                \item \textbf{Prévention (réduit l'aléa) :} Agit sur la \textbf{cause} (drainage, reprofilage, revégétalisation, interdiction construction). Ex : Drainage profond pour baisser nappe $\rightarrow$ augmente $F_s$.
                \item \textbf{Protection (réduit la vulnérabilité) :} Agit sur les \textbf{conséquences} (ouvrages de retenue, filets, déviation, alarmes, évacuation). Ex : Merlon, filet pare-blocs, bassin de réception.
            \end{itemize}
        \item \textbf{Tableau de choix (Mémo BEPC) :}
            \begin{center}
                \begin{tabularx}{\linewidth}{|l|X|X|}\hline
                \rowcolor{popblueL} \textbf{Type de mouvement} & \textbf{Prévention (Cause)} & \textbf{Protection (Effet)} \\ \hline
                \textbf{Chute / Éboulement} (blocs, rapide) & Purge manuelle, grillage plaqué (ancré), ancrage actif (tirants) & \textbf{Filets pare-blocs} (haute énergie), \textbf{Merlons / Bassins de réception}, déviation trajet \\ \hline
                \textbf{Glissement plan/rotatif} (lent, argile/roc) & \textbf{Drainage profond} (drains horizontaux, puits), \textbf{Reprofilage} (alléger tête, charger pied), \textbf{Ancrage} (tirants, micropieux), Revégétalisation & Murs de soutènement (poids, en console), Palplanches, Surveillance (inclinomètres, topographie) \\ \hline
                \textbf{Coulée boueuse / Débris} (rapide, chenal) & Reboisement bassin versant, Seuils de torrent (freiner vitesse), Bassins de décantation en amont & \textbf{Bassins de rétention / Dépôt} (exutoire), \textbf{Digues de déviation}, Alertes crues / Évacuation \\ \hline
                \textbf{Affaissement / Doline} (karst, minier) & Comblement cavités (injection coulis), Contrôle pompages & Remblaiement contrôlé, Surveillance tassement \\ \hline
            \end{tabularx}
        \end{center}
        \item \textbf{Surveillance / Alerte :} Inclinomètres (profondeur glissement), Extensomètres (fissures), Piézomètres (niveau eau), Radar interférométrique (satellite), Observation visuelle (fissures murs, arbres penchés). Système d'alerte précoce (seuils pluviométriques).
    \end{enumerate}
    \textbf{Critère de choix :} \textbf{Coût / Efficacité / Délai / Acceptabilité.} Ex : Filet pare-blocs = rapide, coûteux, visuel. Drainage = lent à mettre en œuvre, durable, invisible. Reprofilage = lourd, définitif.
\end{methode}

\begin{exoresolu}[Dimensionnement simplifié d'un ouvrage de protection : Bassin de réception pour coulée de débris]
    \textbf{Énoncé :} Un torrent en montagne (bassin versant 2 km$^2$) génère des coulées de débris estimées à un volume solide maximal $V_s = 15\,000$ m$^3$ (période de retour 100 ans). La concentration volumique en solide $C_v$ est de 0,60. On veut dimensionner un bassin de réception en pied de versant pour protéger un village en aval.
    \begin{enumerate}
        \item Calculer le volume total de la coulée (solide + eau) $V_{tot}$.
        \item Si le bassin a une surface de fond $S = 800$ m$^2$, quelle doit être sa hauteur utile $H$ pour stocker l'événement centennal avec un coefficient de sécurité 1,2 sur le volume ?
        \item Citer 2 autres fonctions de ce bassin outre le stockage.
        \item Pourquoi le reboisement du bassin versant est-il une mesure de prévention complémentaire indispensable ?
    \end{enumerate}
    \tcblower
    \textbf{Solution détaillée :}
    \begin{enumerate}
        \item \textbf{Volume total $V_{tot}$ :}
            La concentration $C_v = \frac{V_s}{V_{tot}} \Rightarrow V_{tot} = \frac{V_s}{C_v} = \frac{15\,000}{0,60} = 25\,000 \text{ m}^3$.
            (Volume eau $V_e = V_{tot} - V_s = 10\,000$ m$^3$).
        \item \textbf{Hauteur utile $H$ du bassin :}
            Volume de conception $V_{conception} = 1,2 \times V_{tot} = 1,2 \times 25\,000 = 30\,000 \text{ m}^3$.
            $V_{conception} = S \times H \Rightarrow H = \frac{V_{conception}}{S} = \frac{30\,000}{800} = 37,5 \text{ m}$.
            \textit{Remarque :} Une hauteur de 37,5 m est considérable pour un bassin de fond de vallée. En pratique, on augmenterait la surface $S$ (élargir le bassin) ou on accepterait un débordement contrôlé par un déversoir pour des événements plus rares, ou on combinerait avec des seuils en amont pour réduire $V_s$.
        \item \textbf{Autres fonctions :}
            \begin{enumerate}
                \item \textbf{Dissipation d'énergie / Décantation :} Le bassin ralentit la coulée, permet aux gros blocs de se déposer au fond et à l'eau de s'évacuer par un déversoir de fond (déshydratation du dépôt).
                \item \textbf{Protection de l'aval :} Évite l'envasement du lit de la rivière principale, protège les ouvrages d'art (ponts) et les zones urbanisées en aval.
            \end{enumerate}
        \item \textbf{Rôle du reboisement (Prévention) :}
            \begin{itemize}
                \item \textbf{Interception pluviométrique :} La canopée réduit l'énergie cinétique des gouttes (moins d'érosion splash) et évapore une partie de l'eau (moins d'infiltration/ruissellement).
                \item \textbf{Renforcement racinaire :} Les racines augmentent la cohésion $c'$ du sol superficiel (facteur 1,5 à 5), stabilisant les couches minces sources de coulées.
                \item \textbf{Régulation hydrologique :} L'évapotranspiration abaisse l'humidité antécédente du sol, retardant la saturation lors des orages.
                \item \textbf{Piégeage sédimentaire :} La litière et la végétation de sous-bois filtrent le ruissellement diffus.
            \end{itemize}
            \textit{Conclusion :} L'ouvrage de protection (bassin) gère le symptôme (le volume qui arrive). La prévention (reboisement, seuils) agit sur la cause (production du volume solide). Les deux sont nécessaires (stratégie "défense en profondeur").
    \end{enumerate}
\end{exoresolu}

\begin{exoresolu}[Synthèse : Analyse de risque d'un glissement menaçant une route nationale (RN11)]
    \textbf{Énoncé :} Un glissement rotatif lent (2 cm/jour) affecte un versant en marnes argileuses (Éocène) surplombant la RN11 (axe Douala-Bafoussam). Le glissement a une longueur de 300 m, largeur 150 m, épaisseur moyenne 10 m. La route passe en pied de versant. Des fissures apparaissent sur la chaussée. La pluviométrie moyenne est 2000 mm/an. Le pendage des marnes est de 20° vers la route.
    \begin{enumerate}
        \item Rédiger l'analyse de risque structurée (Site, Aléa, Enjeux, Risque).
        \item Proposer 2 mesures de prévention et 2 mesures de protection prioritaires, justifiées par la cinématique (lente) et le matériau (marnes).
        \item Expliquer pourquoi le "traitement à la source" (drainage) est préférable au "traitement au pied" (mur de soutènement) seul.
    \end{enumerate}
    \tcblower
    \textbf{Solution détaillée :}
    \begin{enumerate}
        \item \textbf{Analyse de risque structurée :}
            \begin{itemize}
                \item \textbf{Site (Susceptibilité forte) :} Géologie : \textbf{Marnes argileuses} (roche tendre, imperméable, gonflante, sensible à l'eau, faible résistance au cisaillement résiduelle). Structure : \textbf{Pendage 20° vers la route} (versant \emph{envers} modéré, plans de stratification jourdant dans la pente). Morphologie : Versant convexe ? (non précisé, mais glissement rotatif suggère rupture concave). Hydrogéologie : Nappe perchée probable dans les marnes altérées, alimentée par pluviométrie forte (2000 mm).
                \item \textbf{Aléa (Probabilité forte) :} Phénomène : \textbf{Glissement rotatif lent} (cinématique active : 2 cm/jour = 7 m/an !). Déclencheur : Pluies saisonnières (saturation), érosion pied par route (sous-cavage). Récurrence : Saisonnière (chaque saison des pluies), aggravation continue.
                \item \textbf{Enjeux (Très forts) :} \textbf{RN11} : Axe structurant international (Douala-Bafoussam-Ngaoundéré), trafic poids lourds intense, liaison port/hinterland. Coupure = paralysie économique majeure. Risque direct pour usagers (chute blocs, affaissement chaussée). Habitations / activités riveraines possibles.
                \item \textbf{Risque : \textbf{TRÈS FORT}.} Aléa actif + Enjeux vitaux + Vulnérabilité route non protégée. Urgence absolue.
            \end{itemize}
        \item \textbf{Mesures prioritaires :}
            \begin{itemize}
                \item \textbf{Prévention 1 (Urgence) : Drainage profond par forage horizontal (drains Tiba/Micro-drains) depuis la berme de route vers l'amont du glissement.} \textit{Justification :} Marnes = imperméables $\rightarrow$ drainage lent mais seul moyen de baisser $u$ durablement. Adapté au mouvement lent (temps pour forer).
                \item \textbf{Prévention 2 : Reprofilage / Allègement de la tête de glissement} (si accessible) + \textbf{Revégétalisation} (vetiver, arbres à racines pivotantes) sur la surface du glissement. \textit{Justification :} Réduire la charge motrice $P_{\parallel}$ et augmenter cohésion $c'$ par racines. Évapotranspiration.
                \item \textbf{Protection 1 : Merlon / Butée en pied de versant (enrochements calibrés + géotextile) au niveau de la route.} \textit{Justification :} Reconstitue une butée mécanique pour stopper/ralentir l'avancée du glissement vers la route. Rapide à mettre en œuvre.
                \item \textbf{Protection 2 : Surveillance instrumentée (Inclinomètres dans forages, topographie de surface, piézomètres) + Système d'alerte (feux tricolores, limitation vitesse/poids, fermeture préventive si seuil vitesse dépassé).} \textit{Justification :} Mouvement lent = temps pour réagir. Surveillance = gestion du risque résiduel.
            \end{itemize}
        \item \textbf{Pourquoi Drainage (Source) > Mur (Pied) seul ?}
            \begin{itemize}
                \item \textbf{Mur de soutènement (Pied) :} S'oppose à la poussée du glissement ($P_{\parallel}$). Poussée énorme : volume $\approx 300 \times 150 \times 10 = 450\,000$ m$^3$ $\approx 1$ million de tonnes. Mur devrait être massif (coût exorbitant), risque de rupture par dépassement (glissement passe par-dessus) ou rupture de fondation (marnes molles). Ne traite pas la cause (eau).
                \item \textbf{Drainage (Source) :} Traite la \textbf{cause} (pression interstitielle $u$). Baisse $u$ $\rightarrow$ augmente résistance $\tau_f$ $\rightarrow$ peut stabiliser le versant \textbf{sans ouvrage massif} (Facteur de sécurité $F_s > 1$). Coût bien inférieur. Durable si entretien drains.
                \item \textbf{Approche optimale :} \textbf{Drainage (Traitement cause) + Merlon léger (Protection pied / Sécurité usagers) + Surveillance.} C'est la stratégie "Traitement à la source + Protection passive".
            \end{itemize}
    \end{enumerate}
    \textbf{Conclusion :} Face à un glissement lent dans des marnes sur axe routier majeur, la priorité est le drainage profond (cause) couplé à une protection physique du pied (effet) et une surveillance rigoureuse.
\end{exoresolu}

\begin{exoresolu}[Lecture de carte de zonage (PPRI Glissement) pour un projet de collège]
    \textbf{Énoncé :} On souhaite construire un collège (600 élèves) sur un terrain de 3 ha. La carte de zonage PPRI "Mouvements de terrain" de la commune montre :
    \begin{itemize}
        \item Le terrain est traversé par une \textbf{Zone Bleue} (Aléa Modéré) sur sa partie haute (pente 15\%, marnes).
        \item La partie basse (plat, alluvions) est en \textbf{Zone Blanche}.
        \item Une \textbf{Zone Rouge} (Aléa Fort) correspond à un ancien glissement rotatif stabilisé, situé juste au-dessus de la limite haute du terrain.
    \end{itemize}
    Le règlement de la Zone Bleue impose : "Étude géotechnique G2 obligatoire. Interdiction des fosses septiques. Fondations sur pieux jusqu'au substratum sain. Drainage périphérique du bâtiment. Pas de remblai $>$ 50 cm."
    Le règlement de la Zone Rouge impose : "Inconstructible. Surveillance inclinométrique obligatoire."
    \begin{enumerate}
        \item Le projet est-il réalisable sur ce terrain ? Justifier en localisant les bâtiments.
        \item Lister les 4 mesures techniques obligatoires issues du règlement pour la partie en Zone Bleue.
        \item Quel risque résiduel majeur persiste lié à la Zone Rouge voisine ? Quelle mesure de protection ajouter ?
    \end{enumerate}
    \tcblower
    \textbf{Solution détaillée :}
    \begin{enumerate}
        \item \textbf{Faisabilité : OUI, conditionnellement.}
            \textit{Stratégie d'implantation :} Implanter \textbf{tous les bâtiments (salles de classe, admin, réfectoire, internat) en ZONE BLANCHE} (partie basse plate). Réserver la ZONE BLEUE (partie haute) aux \textbf{espaces extérieurs non construits} : cour de récréation, terrains de sport, parkings, espaces verts, bassins de rétention eaux pluviales.
            \textit{Justification :} La Zone Blanche n'a pas de prescription aléa mouvement de terrain. La Zone Bleue autorise des aménagements légers sous prescriptions, mais interdit les constructions lourdes (collège = enjeu vie humaine très fort, vulnérabilité élevée). Placer 600 élèves en zone aléa modéré est inacceptable au regard du principe de précaution (PPRI).
        \item \textbf{Mesures obligatoires (si locaux techniques/parkings en Zone Bleue) :}
            \begin{enumerate}
                \item \textbf{Essai géotechnique G2 (AVP/PRO) :} Définir la profondeur du substratum sain, paramètres mécaniques marnes, niveau piézométrique.
                \item \textbf{Fondations sur pieux (profonds) :} Transmettre les charges sous la zone glissante/marnes altérées, dans le substratum rocheux sain (grès/calcaire sous-jacent).
                \item \textbf{Drainage périphérique :} Drain périphérique en pied de murs/radier + regard de visite + rejet en zone stable (exutoire) $\rightarrow$ Baisser nappe, éviter pression eau sur murs.
                \item \textbf{Interdiction fosses septiques + Gestion eaux pluviales :} Raccordement réseau collectif obligatoire ou micro-station étanche avec rejet hors zone. Pas d'infiltration sur place (évite recharge nappe).
            \end{enumerate}
            \textit{Note :} Pas de remblai $>$ 50 cm respecté (terrain naturel en Zone Blanche, éventuel léger remblai pour plateforme sport en Zone Bleue).
        \item \textbf{Risque résiduel Zone Rouge (Glissement ancien au-dessus) :}
            \begin{itemize}
                \item \textbf{Risque :} Réactivation du glissement ancien (Zone Rouge) par forte pluie ou séisme $\rightarrow$ \textbf{Éboulement / Glissement secondaire} atteignant la limite haute du terrain (Zone Bleue/Blanche). Ruée de blocs/boue sur les bâtiments implantés en contrebas.
                \item \textbf{Mesure de protection ajoutée :} \textbf{Merlon / Butée de protection} (enrochements + géogrille) à la limite haute du terrain (pied de la Zone Rouge) + \textbf{Filet pare-blocs} ancré sur la Zone Rouge si falaises/blocs instables. Ou \textbf{Déviation du trajet} par une digue en amont.
                \item \textbf{Surveillance :} Exiger la mise en place d'inclinomètres dans la Zone Rouge (obligation règlement) avec transmission données à la mairie/éducation nationale pour alerte précoce.
            \end{itemize}
    \end{enumerate}
    \textbf{Conclusion :} L'aménagement du territoire (choix de l'emprise bâtie en Zone Blanche) est la première et meilleure mesure de prévention. Le respect strict du règlement PPRI (pieux, drainage, assainissement) pour les ouvrages annexes en Zone Bleue et la protection passive contre la Zone Rouge voisine complètent la sécurité.
\end{exoresolu}

\begin{attention}[Les 5 erreurs qui coûtent des points au BEPC]
    \begin{enumerate}
        \item \textbf{Confondre "Aléa" et "Risque".} \textit{Erreur :} "Le risque est fort car la pente est raide." \textit{Correct :} "L'\textbf{aléa} est fort (pente raide, marnes, pluviométrie). Le \textbf{risque} est fort \textbf{car} des enjeux (route, maisons) sont exposés à cet aléa." $\rightarrow$ \textbf{Définitions strictes :} Aléa = probabilité/intensité phénomène. Risque = Aléa $\times$ Enjeux $\times$ Vulnérabilité.
        \item \textbf{Oublier le rôle de l'EAU (mécanique + chimique) dans l'explication de la rupture.} \textit{Erreur :} "La pluie mouille le sol qui glisse." \textit{Correct :} "L'infiltration \textbf{sature} le matériau $\rightarrow$ crée une \textbf{pression interstitielle ($u$)} qui \textbf{diminue la contrainte efficace ($\sigma' = \sigma - u$)} et donc la \textbf{résistance au cisaillement ($\tau_f = c' + \sigma' \tan \phi'$)}. Parallèlement, l'eau \textbf{altère chimiquement} la roche (hydrolyse) et \textbf{gonfle les argiles}." $\rightarrow$ Mots-clés obligatoires : \textbf{Saturation, Pression interstitielle, Contrainte efficace, Résistance au cisaillement, Altération, Gonflement}.
        \item \textbf{Ne pas distinguer Prévention (cause) et Protection (effet) dans les mesures proposées.} \textit{Erreur :} "Pour prévenir le glissement, on construit un mur en bas." \textit{Correct :} "Le mur en bas est une mesure de \textbf{PROTECTION} (réduit vulnérabilité/conséquences). La \textbf{PRÉVENTION} (réduit aléa) serait le \textbf{drainage} pour baisser la nappe ou le \textbf{reprofilage} pour alléger la tête." $\rightarrow$ Classer systématiquement les mesures dans les deux colonnes.
        \item \textbf{Ignorer la structure géologique (pendage, stratification, failles) dans l'analyse de la stabilité.} \textit{Erreur :} "C'est une pente en argile, ça glisse." \textit{Correct :} "C'est un \textbf{versant \emph{envers}** (pendage strates = 25° vers la pente, pente topographique = 20°) dans des \textbf{argiles intercalées dans des grès}. Les plans de stratification jourdant dans la pente servent de \textbf{plans de rupture préférentiels} pour un \textbf{glissement plan}." $\rightarrow$ Le vocabulaire structural (pendage, jour, versants \emph{envers}/\emph{revers}, diaclases) est nota bene au programme.}
        \item \textbf{Réponses trop générales sans s'appuyer sur les données chiffrées ou le contexte local de l'énoncé.} \textit{Erreur :} "Il faut planter des arbres." \textit{Correct :} "Au vu de la pluviométrie de 2000 mm/an et de la nature marneuse, il faut planter du \textbf{Vetiver (Chrysopogon zizanioides)} en lignes de niveau tous les 2 m sur le corps du glissement pour son système racinaire profond (3-4 m) qui augmente la cohésion $c'$ de 5 à 10 kPa et assure l'évapotranspiration estivale." $\rightarrow$ \textbf{Contextualiser :} Utiliser les chiffres de l'énoncé (volume, pluviométrie, pente, nature roche). Citer des exemples/espèces camerounais (Vetiver, Bambou, Eucalyptus - avec prudence pour l'eau).
    \end{enumerate}
\end{attention}

\newpage

\coursec{Exercices}

\coursec{Accidents en zones montagneuses ou liés aux mouvements de terrain}

\courssub{Je vérifie mes connaissances}

\begin{exercice}[QCM : Types de déformations et mouvements]
Parmi les affirmations suivantes, une seule est correcte. La recopier en la complétant.
\begin{enumerate}
\item Un pli est une déformation \textbf{souple} qui se produit lorsque les roches sont soumises à des contraintes \textbf{lentes} et \textbf{faibles} à température \textbf{élevée}.
\item Une faille est une déformation \textbf{cassante} caractérisée par une rupture et un déplacement relatif des compartiments rocheux.
\item Les éboulements et chutes de blocs sont favorisés par la présence de \textbf{diaclases} qui découpent la roche en blocs instables.
\item La lubrification des plans de faille ou de stratification par l'eau de pluie diminue le \textbf{coefficient de frottement} et favorise le \textbf{glissement}.
\end{enumerate}
\end{exercice}

\begin{exercice}[Vrai / Faux justifié : Rôle de l'eau et de l'homme]
Dire si les affirmations sont vraies ou fausses. Justifier chaque réponse par une phrase précise.
\begin{enumerate}
\item \textbf{Faux.} L'infiltration de l'eau dans les fissures \textbf{diminue} la cohésion des roches tendres (argiles, marnes) par gonflement et lubrification, ce qui réduit leur résistance au cisaillement.
\item \textbf{Vrai.} L'action chimique de l'eau (dissolution du calcaire, altération des feldspaths en argile) transforme les minéraux résistants en minéraux tendres, ce qui \textbf{diminue la résistance mécanique} du massif rocheux.
\item \textbf{Faux.} Le déboisement d'un versant \textbf{expose} le sol à l'érosion (ruissellement, battance) et \textbf{déstabilise} les pentes par la perte du renforcement racinaire et du rôle de butée des arbres.
\item \textbf{Vrai.} La construction en bas de pente (recharge/surcharge) augmente la contrainte de cisaillement au niveau du plan de rupture potentiel, tandis que la construction en haut de pente (décharge/entaillage) réduit le contrefort naturel ; les deux peuvent déclencher ou réactiver un glissement.
\end{enumerate}
\end{exercice}

\begin{exercice}[Définitions et vocabulaire technique]
Définir précisément les termes suivants en utilisant le vocabulaire scientifique du cours :
\begin{enumerate}
\item \textbf{Mouvement de terrain (ou mouvement de masse) :} Déplacement spontané, plus ou moins rapide, d'une masse de roches ou de sols le long d'une pente, sous l'effet de la gravité, lorsque les contraintes de cisaillement dépassent la résistance du matériau.
\item \textbf{Déformation souple :} Déformation réversible ou irréversible sans rupture, se produisant sous l'effet de contraintes lentes et continues (ex: \textbf{pli} anticlinal ou synclinal).
\item \textbf{Déformation cassante :} Déformation par rupture brutale du matériau sous l'effet de contraintes rapides ou fortes (ex: \textbf{faille} avec rejet, \textbf{diaclase} sans rejet).
\item \textbf{Diaclase :} Fracture de la roche sans déplacement relatif des deux compartiments (rejet nul), souvent organisée en réseaux (diaclases de stratification, de retrait, tectoniques).
\item \textbf{Angle de repos (ou angle de talus naturel) :} Angle maximal d'inclinaison par rapport à l'horizontale qu'un matériau meuble ou éboulé peut maintenir sans s'effondrer, caractérisant son équilibre statique naturel.
\item \textbf{Mesure de protection active :} Ouvrage ou action qui agit \textbf{sur la cause} du risque pour stabiliser durablement le massif (ex: \textbf{drainage profond}, \textbf{soutènement par clouage}, \textbf{reboisement}). \textbf{Mesure de protection passive :} Ouvrage qui \textbf{protège les biens et les personnes} sans stabiliser le massif lui-même (ex: \textbf{meridienne pare-blocs}, \textbf{filet de protection}, \textbf{déviation de route}).
\end{enumerate}
\end{exercice}

\begin{exercice}[Calcul direct : Angle de pente et stabilité]
Un talus routier a été décapé à la verticale (pente = \SI{90}{\degree}) dans une formation argileuse. Après quelques jours de pluie, la pente s'effondre et se stabilise naturellement avec une inclinaison de \SI{35}{\degree} par rapport à l'horizontale.
\begin{enumerate}
\item Cet angle de \SI{35}{\degree} atteint après l'éboulis est l'\textbf{angle de repos} (ou angle de talus naturel) de cette argile saturée.
\item Donnée : $\mu = 0,70$. Relation : $\mu = \tan\varphi$.
  \[ \varphi = \arctan(0,70) \approx \SI{35,0}{\degree} \]
  L'angle de frottement calculé (\SI{35}{\degree}) est \textbf{égal} à l'angle de repos observé (\SI{35}{\degree}).
\item Comme l'angle de la pente stabilisée (\SI{35}{\degree}) est \textbf{égal} à l'angle de frottement interne $\varphi$, la pente est en \textbf{équilibre limite} (stable pour l'instant, mais tout surcroît de contrainte ou d'eau la rendra instable). Pour un matériau cohérent (argile), la cohésion résiduelle assure une marge de sécurité nulle ici.
\end{enumerate}
\end{exercice}

\courssub{Je m'entraîne}

\begin{exercice}[Analyse d'un affleurement : Plis, failles et risque d'éboulement]
On observe sur une photographie d'affleurement (schéma ci-dessous) des couches sédimentaires affectées par une déformation.
\begin{popfigure}
\begin{tikzpicture}[scale=0.85]
% Couches plissées (anticlinal puis synclinal)
\draw[popblue, thick] plot[smooth, tension=0.7] coordinates {(0,0) (1,0.6) (2,0) (3,-0.6) (4,0) (5,0.6) (6,0)};
\draw[popblue, thick] plot[smooth, tension=0.7] coordinates {(0,-0.6) (1,0) (2,-0.6) (3,-1.2) (4,-0.6) (5,0) (6,-0.6)};
\draw[popblue, thick] plot[smooth, tension=0.7] coordinates {(0,-1.2) (1,-0.6) (2,-1.2) (3,-1.8) (4,-1.2) (5,-0.6) (6,-1.2)};
\draw[popblue, thick] plot[smooth, tension=0.7] coordinates {(0,-1.8) (1,-1.2) (2,-1.8) (3,-2.4) (4,-1.8) (5,-1.2) (6,-1.8)};
% Faille normale traversant (verticale)
\draw[poporange, very thick, decorate, decoration={zigzag, segment length=4, amplitude=1.5}] (3,0.8) -- (3,-2.6);
% Blocs instables en haut à gauche (versant pentu)
\draw[popdark, fill=poporange!20] (0.5,0.8) -- (1.5,1.3) -- (1.2,0.1) -- cycle;
\draw[popdark, fill=poporange!20] (1.8,0.8) -- (2.5,1.2) -- (2.2,0.0) -- cycle;
% Flèches mouvement faille : FAILLE NORMALE (bloc gauche = toit = descend)
\draw[->, popdark, thick] (2.6,0.2) -- (2.9,-0.2) node[right, pos=1.2, font=\small] {Bloc toit $\downarrow$};
\draw[->, popdark, thick] (3.4,-0.2) -- (3.7,0.2) node[right, pos=1.2, font=\small] {Bloc mur $\uparrow$};
\node[popdark, font=\small] at (1.5,1.6) {Blocs instables};
\node[poporange, font=\small] at (3.5,1.0) {Faille};
\node[popblue, font=\small] at (5.5,1.0) {Couches plissées};
\end{tikzpicture}
\legende{Schéma simplifié d'un affleurement montrant des plis (bleu) et une faille normale (orange à zigzag). Les flèches indiquent le mouvement relatif des compartiments.}
\end{popfigure}
\begin{enumerate}
\item \textbf{Structures bleues (plis) :} Déformation \textbf{souple} (courbure continue des couches sans rupture). \textbf{Structure orange (faille) :} Déformation \textbf{cassante} (rupture nette avec rejet visible via le zigzag et les flèches).
\item La structure bleue centrale (au niveau de la faille, $x \approx 3$) présente une forme en $\cup$ (creux vers le bas) : c'est un \textbf{synclinal}. Les structures de part et d'autre (formes en $\cap$) sont des \textbf{anticlinaux}.
\item Les blocs hachurés (haut gauche) risquent de s'effondrer par \textbf{éboulement/chute de blocs} car : (1) ils sont situés sur un versant pentu (pente > angle de repos) ; (2) ils sont découpés et isolés par les \textbf{diaclases} (fissures non représentées mais implicites dans le massif) et la \textbf{faille} qui tronque les couches ; (3) la faille a créé une surface de rupture affaiblie (gougeon, roche pulvérisée) et a surélevé le bloc mur (droite), créant un dévers ou une surcharge instable côté toit (gauche).
\item C'est une \textbf{faille normale} (extensif). Le compartiment \textbf{toit} (bloc de gauche, au-dessus du plan de faille) s'est déplacé \textbf{vers le bas} (flèche descendante) par rapport au compartiment \textbf{mur} (bloc de droite) qui est remonté relativement. Cela correspond à un jeu extensif.
\end{enumerate}
\end{exercice}

\begin{exercice}[Interprétation de courbes : Pluviométrie et glissements de terrain]
Le tableau ci-dessous donne les cumuls mensuels de pluviométrie (en mm) et le nombre de glissements de terrain recensés dans la localité de Mbankolo (périphérie de Yaoundé) sur une année.

\begin{center}
\begin{tabularx}{\linewidth}{|c|*{12}{X|}}\hline
Mois & Janv & Fév & Mars & Avr & Mai & Juin & Juil & Août & Sept & Oct & Nov & Déc \\ \hline
Pluie (mm) & 15 & 30 & 120 & 180 & 210 & 150 & 60 & 80 & 200 & 250 & 100 & 20 \\ \hline
Nb glissements & 0 & 0 & 2 & 5 & 8 & 4 & 1 & 1 & 6 & 9 & 3 & 0 \\ \hline
\end{tabularx}
\end{center}

\begin{enumerate}
\item \textbf{Représentation graphique :} Sur un même graphique (abscisses : mois), tracer l'histogramme des pluviométries (axe Y gauche, mm) et la courbe en points reliés du nombre de glissements (axe Y droite, nombre).
\item \textbf{Évolution parallèle :} Les deux grandeurs évoluent globalement en phase. Les glissements débutent en mars (après le début des pluies), culminent en \textbf{mai} (8 glissements) et \textbf{octobre} (9 glissements), correspondant aux deux pics de pluviométrie (mai : 210 mm, octobre : 250 mm). L'activité est nulle en saison sèche (janv, fév, déc).
\item \textbf{Décalage :} Le pic de pluviométrie maximal est en \textbf{octobre} (250 mm). Le pic de glissements maximal est aussi en \textbf{octobre} (9 glissements). Le délai est d'environ \textbf{0 à 1 mois} (quasi-simultané pour le second pic). Pour le premier pic (mai), les pluies montent dès mars/avril, les glissements suivent avec 1-2 mois de retard (inertie de la saturation). \textbf{Explication hydrogéologique :} Ce délai correspond au temps nécessaire pour que l'eau de pluie \textbf{infiltre}, \textbf{sature la zone non saturée}, fasse \textbf{monter la nappe phréatique} jusqu'au plan de rupture potentiel, et génère une \textbf{pression interstitielle ($u$)} suffisante pour réduire la contrainte effective ($\sigma' = \sigma - u$) et déclencher la rupture.
\item \textbf{Rôle déclencheur de l'eau :} L'eau est le \textbf{facteur déclencheur} principal : par sa pression interstitielle, elle diminue le frottement interne (terme $\sigma' \tan\varphi$ dans la résistance au cisaillement $\tau = c' + \sigma' \tan\varphi$) et assure la lubrification des plans de glissement, faisant basculer le massif de l'équilibre stable à l'instabilité.
\end{enumerate}
\end{exercice}

\begin{exercice}[Étude de cas comparée : Versant boisé vs Versant urbanisé]
Deux versants voisins (même géologie : schistes altérés en argiles, même pente \SI{30}{\degree}, même exposition aux pluies) présentent des comportements différents après un orage violent (\SI{80}{mm} en 2h) :
\begin{itemize}
\item \textbf{Versant A :} Couvert par une forêt dense (arbres à racines profondes, litière épaisse). Aucun dommage.
\item \textbf{Versant B :} Déboisé il y a 5 ans pour lotissement (maisons, routes, dalles béton, canalisations d'eaux usées fuyantes). Apparition de fissures dans les murs, affaissement de la route, glissement de terrain détruisant deux maisons en bas de pente.
\end{itemize}
\begin{enumerate}
\item \textbf{Trois facteurs anthropiques sur le Versant B :}
      \begin{enumerate}
      \item \textbf{Imperméabilisation} (toitures, routes, dalles) $\rightarrow$ augmentation du ruissellement de surface, diminution de l'infiltration naturelle $\rightarrow$ saturation plus rapide et plus intense des argiles en pied de pente.
      \item \textbf{Fuites des réseaux d'assainissement/eaux usées} $\rightarrow$ apport d'eau direct et concentré dans le sous-sol $\rightarrow$ montée anormale de la nappe / pression interstitielle.
      \item \textbf{Surcharge (recharge) en haut de pente} (maisons, remblais) et \textbf{décharge (entaillage)} pour les routes $\rightarrow$ modification de la géométrie de la pente (augmentation contrainte cisaillement, diminution contrefort).
      \item (Bonus) \textbf{Suppression de la couverture végétale} $\rightarrow$ perte de l'interception pluviométrique, de l'évapotranspiration et du renforcement racinaire.
      \end{enumerate}
\item \textbf{Rôle mécanique des racines (Versant A) :}
      \begin{itemize}
      \item \textbf{Renforcement racinaire :} Les racines traversent les plans de rupture potentiels et agissent comme des \textbf{ancrages} (clous naturels) qui augmentent la cohésion apparente du sol ($c'_{app} = c'_{sol} + c'_{racines}$).
      \item \textbf{Effet de butée :} Les souches et gros bois agissent comme des \textbf{butées} passives qui s'opposent au mouvement des masses superficielles par leur enfoncement dans le substratum stable.
      \end{itemize}
\item \textbf{Rôle hydrologique :}
      \begin{itemize}
      \item \textbf{Forêt (Versant A) :} \textbf{Interception} (litière/feuillage $\approx$ 20-30\% pluie), \textbf{Infiltration contrôlée} (porosité structurée par racines/faune), \textbf{Évapotranspiration} forte (pompage racinaire) $\rightarrow$ \textbf{teneur en eau modérée}, \textbf{pression interstitielle basse}.
      \item \textbf{Urbanisation (Versant B) :} \textbf{Imperméabilisation} $\rightarrow$ ruissellement concentré (rues, gouttières) $\rightarrow$ infiltration localisée massive en pied de pente / fuites réseaux $\rightarrow$ \textbf{saturation rapide}, \textbf{montée nappe}, \textbf{pression interstitielle élevée} $\rightarrow$ perte de résistance.
      \end{itemize}
\item \textbf{Conclusion :} La couverture végétale est une \textbf{mesure de prévention active} essentielle, gratuite et durable : elle agit simultanément sur les causes hydrauliques (régulation eau) et mécaniques (renforcement/butée) de l'instabilité, contrairement aux ouvrages de génie civil passifs ou curatifs.
\end{enumerate}
\end{exercice}

\begin{exercice}[QCM de synthèse : Techniques de prévention / protection]
Associer chaque technique (colonne de gauche) à la cause principale du risque qu'elle combat (colonne de droite). Une technique peut correspondre à plusieurs causes.

\begin{center}
\begin{tabularx}{\linewidth}{|X|X|}\hline
\textbf{Techniques} & \textbf{Causes / Mécanismes combattus} \\ \hline
1. Drainage profond (drains, puits de décompression) & C. Excès d'eau interstitielle / Pression de pore / Lubrification plans de glissement \\ \hline
2. Soutènement (mur de soutènement, palplanches, clous, ancrages) & D. Sous-charge / Érosion du pied (décharge) \quad \text{et} \quad A. Surcharge en haut de pente (recharge) \\ \hline
3. Reboisement / Revégétalisation (herbacées, arbres) & B. Érosion superficielle / Ruissellement \quad \text{et} \quad C. Excès d'eau interstitielle (évapotranspiration) \\ \hline
4. Purge / Déchargement en haut de pente & A. Surcharge en haut de pente (recharge) \\ \hline
5. Bénétage (création de paliers) & E. Pente trop raide / Angle > angle de repos \\ \hline
6. Méridiennes / Rideaux pare-blocs (filets, barrières) & F. Chute de blocs / Éboulements / Propagation \\ \hline
\end{tabularx}
\end{center}
\begin{enumerate}
\item \textbf{Correspondances :} 1-C, 2-A/D, 3-B/C, 4-A, 5-E, 6-F.
\item \textbf{Classification :}
      \begin{itemize}
      \item \textbf{Protection Active (stabilise le massif) :} 1 (Drainage), 2 (Soutènement), 3 (Reboisement), 4 (Purge), 5 (Bénétage).
      \item \textbf{Protection Passive (protège biens/people) :} 6 (Méridiennes/Pare-blocs).
      \end{itemize}
\end{enumerate}
\end{exercice}

\courssub{Je me prépare au BEPC}

\begin{exercice}[Sujet type BEPC : Rapport d'expertise pour un quartier instable à Yaoundé]
\textbf{Contexte :} Le quartier \emph{Nkolbisson-Nord} est construit sur un versant en pente douce (\SI{15}{\degree}) constitué d'argiles plastiques (marnes) surmontées de latérite. Après les fortes pluies de septembre-octobre, de nombreuses maisons présentent des fissures diagonales aux angles des fenêtres, le sol du terrain de football s'est fissuré en arc de cercle, et des arbres penchent vers l'aval. Le sous-préfet demande un rapport d'expertise géotechnique rapide.

\begin{enumerate}
\item \textbf{Diagnostic :} Il s'agit d'un \textbf{glissement de terrain rotatif (ou en cuillère / \emph{slump})}.
      \textbf{Justification :} La pente douce (\SI{15}{\degree}) et la nature du terrain (argiles plastiques homogènes sur roof imperméable) favorisent une surface de rupture courbe (cylindrique). Les \textbf{fissures en arc de cercle} (tête de glissement), les \textbf{fissures diagonales (cisaillement)} aux angles des fenêtres des maisons (rigides) et les \textbf{arbres penchés vers l'aval} (pistolets) sont les indices morphologiques caractéristiques d'une rotation du massif autour d'un axe horizontal transversal.
\item \textbf{Causes (trois probables) :}
      \begin{enumerate}
      \item \textbf{Saturation en eau des argiles (Naturelle) :} Pluies intenses $\rightarrow$ infiltration $\rightarrow$ montée de la nappe phréatique dans la couche argileuse $\rightarrow$ augmentation de la \textbf{pression interstitielle ($u$)} $\rightarrow$ diminution de la contrainte effective ($\sigma' = \sigma - u$) $\rightarrow$ chute de la résistance au cisaillement ($\tau_f = c' + \sigma' \tan\varphi'$).
      \item \textbf{Surcharge anthropique en tête de versant (Recharge) :} Constructions (maisons, terrain de football remblayé) en haut de pente $\rightarrow$ augmentation de la contrainte verticale $\sigma$ et de la contrainte de cisaillement $\tau$ mobilisée sur le plan de rupture $\rightarrow$ rapprochement de l'état limite.
      \item \textbf{Affaiblissement mécanique par fissuration / Décharge au pied :} Fissures de retrait-gonflement des argiles (alternance sec/humide) créant des plans de faiblesse ; éventuels petits entaillages au pied pour routes/caniveaux $\rightarrow$ perte de contrefort.
      \end{enumerate}
\item \textbf{Mesures d'urgence (Protection Passive) :}
      \begin{enumerate}
      \item \textbf{Évacuation immédiate} des habitations les plus menacées (fissures traversantes, ouverture portes/fenêtres bloquée) et mise en sécurité de la population (hébergement d'urgence).
      \item \textbf{Interdiction d'accès} et mise en place d'un périmètre de sécurité (rubalise, garde) autour de la zone de rupture (tête, corps, pied) pour éviter tout surcroît de charge ou victime secondaire.
      \end{enumerate}
\item \textbf{Mesures de stabilisation définitive (Protection Active) :}
      \begin{enumerate}
      \item \textbf{Drainage vertical par puits de décompression / drains horizontaux forés :} Abaissement durable de la nappe phréatique dans la couche argileuse $\rightarrow$ \textbf{diminution de la pression interstitielle $u$} $\rightarrow$ \textbf{augmentation de la contrainte effective $\sigma'$} $\rightarrow$ \textbf{augmentation du frottement $\sigma'\tan\varphi'$ et de la résistance au cisaillement} sur le plan de rupture.
      \item \textbf{Purge / Reprofilage de la tête de glissement :} Enlèvement des terres instables et des surcharges (remblais, constructions légères) en zone de tête $\rightarrow$ \textbf{diminution de la poussée active} et de la contrainte de cisaillement mobilisée $\rightarrow$ rétablissement de l'équilibre (Facteur de Sécurité $F_s > 1,3$).
      \item \textbf{Soutènement par paroi clouée / berlinoise au pied (si place) ou butage :} Création d'un \textbf{contrefort rigide} au pied du glissement $\rightarrow$ \textbf{augmentation de la résistance d'appui (butée)} $\rightarrow$ compensation de la perte de contrefort naturel. (Alternative : \textbf{Reboisement profond} en tête et corps du glissement pour renforcement racinaire et évapotranspiration long terme).
      \end{enumerate}
\item \textbf{Recommandation administrative :}
      \og \textbf{Interdiction formelle de toute nouvelle construction dans la zone instable délimitée ; mise en place immédiate d'un Plan de Prévention des Risques (PPR) mouvements de terrain prescrivant les travaux de drainage et de purge ci-dessus, et l'installation d'un suivi topométrique (inclinomètres, piézomètres) pour valider la stabilisation avant toute réoccupation.} \fg
\end{enumerate}
\end{exercice}

\begin{exercice}[Problème de synthèse : Aménagement d'un talus routier éboulé sur la route Ngaoundéré-Garoua]
\textbf{Contexte :} Sur la RN1, un talus de déblai haut de \SI{12}{m}, taillé à \SI{70}{\degree} dans des grès fracturés (diaclases verticales espacées de \SI{0,5}{m} à \SI{1}{m}) intercalés de niveaux argileux, s'est partiellement effondré sur la chaussée après les pluies d'août. La coupe géologique schématique du talus instable est donnée ci-dessous.

\begin{popfigure}
\begin{tikzpicture}[scale=0.7]
% Profil initial (ligne pointillée)
\draw[popmuted, dashed, thick] (0,0) -- (0,12) -- (15,12);
% Profil éboulé (éboulis au pied)
\draw[popdark, very thick] (0,0) -- (0,12);
\draw[popdark, very thick] (0,12) -- (2,12); % Replat haut
% Face éboulée
\draw[popdark, very thick] (2,12) -- (4,6) -- (6,0) -- (15,0);
% Éboulis
\draw[poporange, fill=poporange!30] (4,6) -- (6,0) -- (8,0) -- (8,4) -- cycle;
% Couches : Grès (hachures) / Argile (points) - Uniquement sur la partie stable (gauche) pour lisibilité
\foreach \y in {1,3,5,7,9,11} {
  \draw[popblue, pattern=north east lines, pattern color=popblue] (0,\y) rectangle (2,\y+1.5); % Grès haut
}
\foreach \y in {0,2,4,6,8,10} {
  \draw[poporange, pattern=dots, pattern color=poporange] (0,\y) rectangle (2,\y+1); % Argile
}
% Diaclases verticales
\draw[popdark, thick] (0.5,0) -- (0.5,12);
\draw[popdark, thick] (1.5,0) -- (1.5,12);
% Nappe phréatique
\draw[popblue, thick, decorate, decoration={snake, amplitude=1, segment length=4}] (0,3) -- (4,3);
\node[popblue, right] at (4,3) {Nappe};
% Flèches poussée argile
\draw[->, popred, thick] (1,2.5) -- (2.5,2.5);
\draw[->, popred, thick] (1,4.5) -- (2.5,4.5);
\node[popred, font=\small] at (3,3.5) {Poussée argile gonflée};
% Route
\draw[popdark, fill=popdark!20] (6,0) rectangle (15,-1);
\node at (10,-0.5) {Chauss\'ee RN1};
\node at (1,13) {Haut : 12 m};
\node at (1,6) {Gr\`es fractur\'e / Argile};
\node at (7,2) {\'Eboulis};
\end{tikzpicture}
\legende{Coupe schématique du talus routier éboulé (non à l'échelle). Grès (hachures bleues), Argiles (points orange), Diaclases verticales (traits noirs), Nappe (ligne bleue ondulée), Poussée argile (flèches rouges).}
\end{popfigure}

\begin{enumerate}
\item \textbf{Analyse des causes (Mécanisme de rupture combiné) :}
      \begin{itemize}
      \item \textbf{Rôle des diaclases (Grès) :} Les diaclases verticales découpent le grès en \textbf{blocs prismatiques} (largeur 0,5-1 m). La pente taillée à \SI{70}{\degree} libère ces blocs qui ne sont plus maintenus par le confinement latéral. Ils basculent ou glissent par \textbf{basculement (toppling)} ou \textbf{glissement plan} sur les niveaux argileux intercalés (jour des diaclases en face de pente).
      \item \textbf{Rôle des niveaux argileux :} Les argiles gonflantes (smectites) absorbent l'eau $\rightarrow$ \textbf{gonflement} $\rightarrow$ génération d'une \textbf{poussée de gonflement} latérale (flèches rouges) qui écarte les blocs de grès et lubrifie les plans de contact grès/argile. Elles constituent des \textbf{plans de glissement} préférentiels (faible $\varphi'$, $c'$ faible une fois saturées).
      \item \textbf{Rôle de l'eau (Nappe) :} La nappe phréatique intersecte la face du talus. L'eau infiltre les diaclases, sature les argiles $\rightarrow$ \textbf{pression interstitielle $u$ élevée} dans les argiles $\rightarrow$ \textbf{perte de cohésion $c'$ et de frottement $\tan\varphi'$} sur les plans de stratification. L'eau lubrifie aussi les diaclases.
      \item \textbf{Rôle de la géométrie :} Pente taillée à \SI{70}{\degree} $\gg$ Angle de repos des éboulis (\SI{35}{\degree}--\SI{40}{\degree}) et $\gg$ Angle de frottement des argiles saturées (\SI{18}{\degree}--\SI{25}{\degree}). La pente est \textbf{structurellement instable} dès la taille (surcreusement).
      \end{itemize}
\item \textbf{Projet d'aménagement complet (Description du croquis à réaliser) :}
      Sur une copie du schéma, dessiner et légender :
      \begin{enumerate}
      \item \textbf{Reprofilage / Purge :} Nouvelle pente adoucie à \SI{45}{\degree} (ou angle de repos + marge) en haut (recharge) et/ou bénétage.
      \item \textbf{Bénétage :} Création de \textbf{3 à 4 banquettes} (paliers de \SI{2}{m} larg.) tous les \SI{3}{m} de hauteur pour casser la pente, intercepter les blocs et permettre l'entretien/drainage.
      \item \textbf{Drainage :} \textbf{Drains horizontaux forés} (long. \SI{10}{--15}{m}, pente \SI{5}{\%}) depuis les banquettes vers l'intérieur du massif pour abaisser la nappe ; \textbf{Fossé de pied} étanche + caniveau pour évacuer l'eau.
      \item \textbf{Soutènement :} \textbf{Paroi clouée} (clous passifs \SI{6}{--9}{m}, treillis, béton projeté) sur la face reprofilée, ou \textbf{Mur en L} au pied si emprise limitée.
      \item \textbf{Revégétalisation :} Hydro-ensemencement sur banquettes + plantation d'arbres à racines pivotantes (ex: \emph{Acacia}, \emph{Eucalyptus}) en haut de talus.
      \item \textbf{Protection passive :} \textbf{Méridienne pare-blocs} (filet câblé \SI{2000}{kJ} ou barrière rigide) au pied du talus, devant la chaussée, pour arrêter les blocs résiduels.
      \end{enumerate}
\item \textbf{Justification technique (exemples pour 4 techniques) :}
      \begin{enumerate}
      \item \textbf{Bénétage :} Réduit la hauteur de pente libre et crée des banquettes qui \textbf{interceptent les blocs détachés des diaclases}, réduisant l'énergie cinétique d'éventuels éboulements secondaires et facilitant le drainage.
      \item \textbf{Drains horizontaux forés :} Abaissent la nappe phréatique au cœur du massif $\rightarrow$ \textbf{désaturent les niveaux argileux} $\rightarrow$ \textbf{suppriment la poussée de gonflement} et \textbf{rétablissent la résistance au cisaillement} ($c', \varphi'$) sur les plans de stratification.
      \item \textbf{Paroi clouée / Mur de soutènement :} Apporte une \textbf{réaction d'appui (butée)} au pied du talus qui compense la décharge due à l'éboulis et \textbf{confine les blocs de grès} instables, empêchant leur basculement.
      \item \textbf{Méridienne pare-blocs :} \textbf{Intercepte et stoppe} les blocs de grès (taille \SI{0,5}{--1}{m}) qui pourraient encore se détacher des diaclases, \textbf{protégeant la chaussée et les usagers} (protection passive) en attendant la stabilisation active complète.
      \end{enumerate}
\item \textbf{Surveillance post-travaux :}
      \begin{enumerate}
      \item \textbf{Inclinometre(s) (forage(s) incliné(s) dans le massif) :} Mesure le \textbf{profil de déplacement horizontal en profondeur} (vitesse, direction, profondeur du plan de rupture) $\rightarrow$ valide l'arrêt du mouvement (stabilité).
      \item \textbf{Piézomètre(s) (tube perforé en pointe dans les argiles/nappe) :} Mesure la \textbf{charge piézométrique (niveau d'eau / pression interstitielle $u$)} $\rightarrow$ valide l'efficacité du drainage (baisse durable de la nappe en dessous des plans de rupture).
      \end{enumerate}
      (Complément possible : \textbf{Topométrie de surface} (prismes sur paroi) pour suivi déplacements surfaciques ; \textbf{Extensomètre} pour ouverture diaclases).
\end{enumerate}
\end{exercice}

\begin{exercice}[Sujet type BEPC : Synthèse documentaire - Le glissement de terrain de Mbankolo (Yaoundé), octobre 2023]
\textbf{Documents fournis :}

\textbf{Document 1 : Extrait du bulletin météo (Météo Cameroun).}
\og Octobre 2023 : Pluviométrie exceptionnelle à Yaoundé. Cumul mensuel : \SI{420}{mm} (moyenne 1981-2010 : \SI{280}{mm}). Intensités horaires maximales : \SI{45}{mm/h} le 12/10. Sol saturé en eau dès le 05/10. \fg

\textbf{Document 2 : Coupe géologique simplifiée du site (Forage S1).}
\begin{center}
\begin{tabular}{|c|c|c|}\hline
\textbf{Profondeur (m)} & \textbf{Lithologie} & \textbf{Paramètres géotechniques} \\ \hline
0 -- 3 & Latérite dure (ferrugineuse), fissurée & Perméable, rigide \\ \hline
3 -- 12 & Argiles plastiques grises (altération schistes) & $w = 35\%$, $w_L = 65\%$, $I_p = 30\%$, $c' = 5$ kPa, $\varphi' = \SI{18}{\degree}$ \\ \hline
12 -- 15 & Schistes micacés altérés (Saprolite) & Fissurés, perméabilité moyenne \\ \hline
$> 15$ & Schistes sains (Roof) & Imperméable, dur \\ \hline
\end{tabular}
\end{center}
\emph{Légende : $w$ = teneur en eau naturelle, $w_L$ = limite de liquidité, $I_p$ = indice de plasticité, $c'$ = cohésion effective, $\varphi'$ = angle de frottement effectif.}

\textbf{Document 3 : Photographie aérienne du quartier (2010 vs 2023).}
2010 : Zone boisée (eucalyptus), quelques cases traditionnelles en haut de pente.
2023 : Zone dense, maisons en parpaings, routes non revêtues, réseaux d'assainissement inexistants (fossés à ciel ouvert souvent bouchés), décharge sauvage en haut de pente.

\textbf{Document 4 : Témoignage d'un riverain.}
\og L'eau sortait du sol au milieu de la cour, le mur de clôture a craqué, puis la maison a glissé de \SI{2}{m} vers le bas en une nuit. \fg

\begin{enumerate}
\item \textbf{Contexte géologique :}
      \begin{itemize}
      \item \textbf{Roche mère (Roof) :} Schistes sains ($> \SI{15}{m}$), imperméables, base rigide du système.
      \item \textbf{Formation superficielle :} Altérites schisteuses : Saprolite fissurée (\SI{12}{--15}{m}) surmontée d'une épaisse couche d'\textbf{argiles plastiques grises (\SI{3}{--12}{m})} (haute plasticité $I_p=30\%$, $w_L=65\%$), elle-même coiffée d'une \textbf{latérite dure fissurée (\SI{0}{--3}{m})} (perméable, rigide).
      \item \textbf{Plan de rupture probable :} Au \textbf{toit des argiles plastiques (vers \SI{12}{m} profondeur)} ou \textbf{à l'intérieur de la couche d'argile} (zone la plus faible).
      \item \textbf{Justification :} Les argiles ont les paramètres les plus faibles : \textbf{$c' = \SI{5}{kPa}$ (très faible cohésion)} et \textbf{$\varphi' = \SI{18}{\degree}$ (faible frottement)}. De plus, elles sont \textbf{saturées} ($w \approx w_L$ ? $w=35\%, w_L=65\%$, mais Doc 1 dit "sol saturé"), situées juste au-dessus du roof imperméable qui barre l'écoulement vertical $\rightarrow$ \textbf{piège hydraulique} : l'eau s'accumule sur le roof, saturant la base de l'argile $\rightarrow$ pression interstitielle maximale au contact argile/schiste.
      \end{itemize}
\item \textbf{Facteur déclencheur (Pluviométrie Oct 2023) :}
      Le cumul exceptionnel (\SI{420}{mm} vs \SI{280}{mm} moy.) et l'intensité (\SI{45}{mm/h}) ont provoqué une \textbf{infiltration massive et rapide}. La latérite fissurée (perméable) a transmis l'eau vers les argiles. La saturation précoce (dès le 05/10, Doc 1) a entraîné la \textbf{montée de la nappe phréatique} à l'intérieur de la couche d'argile jusqu'au contact avec le roof imperméable.
      \textbf{Argumentation contrainte effective :} La contrainte totale $\sigma$ (poids du sol) reste constante. La pression interstitielle $u$ augmente fortement (nappe montante, saturation). Selon la loi des contraintes effectives de Terzaghi : $\sigma' = \sigma - u$. L'augmentation de $u$ entraîne une \textbf{diminution drastique de $\sigma'$}. La résistance au cisaillement $\tau_f = c' + \sigma' \tan\varphi'$ chute donc brutalement (car $c'=5$ kPa est déjà très faible). Lorsque $\tau_{mobilisée} > \tau_f$, la rupture se produit.
\item \textbf{Facteurs aggravants anthropiques (Doc 3) :}
      \begin{enumerate}
      \item \textbf{Imperméabilisation (toitures, routes, dalles) :} Réduit l'infiltration diffuse naturelle, concentre les eaux de ruissellement vers les bas-fonds/fossés bouchés $\rightarrow$ \textbf{recharge localisée massive} de la nappe en pied de pente / zones basses.
      \item \textbf{Absence/insuffisance d'assainissement (fossés bouchés, pas de réseau) :} Eaux pluviales et usées stagnent, s'infiltrent localement $\rightarrow$ \textbf{apport d'eau direct} dans le sous-sol (recharge artificielle) $\rightarrow$ maintien saturation.
      \item \textbf{Surcharge en tête de pente (décharge sauvage, constructions) :} Ajout de poids (contrainte verticale $\sigma$ augmentée) $\rightarrow$ augmentation de la contrainte de cisaillement $\tau$ mobilisée sur le plan de rupture $\rightarrow$ rapprochement de la rupture.
      \item \textbf{Déboisement (disparition eucalyptus) :} Perte de l'\textbf{évapotranspiration} (pompage racinaire $\approx$ \SI{500}{--1000}{mm/an}) qui maintenait la nappe basse ; perte du \textbf{renforcement racinaire} dans la latérite/argile superficielle.
      \end{enumerate}
\item \textbf{Typologie du mouvement :}
      \begin{itemize}
      \item \textbf{Vitesse :} \textbf{Rapide} (\og en une nuit \fg, déplacement \SI{2}{m}).
      \item \textbf{Cinématique :} \textbf{Glissement translationnel (planaire) ou rotatif peu profond}. La présence d'un plan de rupture net (contact argile/roof ou dans l'argile) et le témoignage (\og glissé de 2m \fg, \og eau sortait du sol \fg = exutoire de nappe au pied) suggèrent un glissement sur un plan préexistant (stratification argile/schiste) plutôt qu'une coulée de boue fluide. La latérite rigide (0-3m) forme une \textbf{carapace} qui se fracture (mur clôture craque) et glisse en bloc sur l'argile plastifiée.
      \end{itemize}
\item \textbf{Mesures correctives priorisées :}
      \begin{center}
      \begin{tabularx}{\linewidth}{|l|X|X|}\hline
      \textbf{Étape} & \textbf{Technique précise} & \textbf{Justification (Documents)} \\ \hline
      \textbf{Urgence (J+0 à J+30)} & \textbf{Évacuation + Méridienne pare-blocs / Filets au pied + Surveillance inclinométrique/piezométrique.} & Doc 4 (maison glissée, danger vie) ; Doc 2 (bloc latérite rigide qui peut basculer). Protège vies/biens immédiatement. \\ \hline
      \textbf{Consolidation (Mois 1--6)} & \textbf{Drainage profond : Puits de décompression / Drains horizontaux forés dans l'argile (niveau 3-12m) pour abaisser nappe sous le plan de rupture.} & Doc 1 (saturation cause déclenchante) ; Doc 2 (nappe perchée sur roof, $c'=5$ kPa). Agit sur la cause hydraulique principale ($u \downarrow \rightarrow \sigma' \uparrow \rightarrow \tau_f \uparrow$). \\ \hline
      \textbf{Aménagement durable (An 1+)} & \textbf{1. Reprofilage/Purge tête de pente (enlèvement décharge/surcharge). 2. Réseau d'assainissement complet (collecteurs étanches, bassins rétention). 3. Reboisement espèces locales racines profondes + Interdiction construction zone (PPR).} & Doc 3 (causes anthropiques : surcharge, fuites, déboisement, imperméabilisation). Traite les causes racines pour $F_s$ durable $> 1,3$. \\ \hline
      \end{tabularx}
      \end{center}
\end{enumerate}
\end{exercice}

\newpage

\coursec{Corrigés des exercices}

\begin{corrige}[Exercice 1 — QCM : Types de déformations et mouvements]
\tcblower
\begin{enumerate}
\item \textbf{Correcte.} Un pli est bien une déformation \cle{souple} : la roche se courbe sans se rompre lorsque les contraintes sont lentes et prolongées, surtout lorsque la température est élevée (la roche est alors plus plastique).
\item \textbf{Correcte.} Une faille est une déformation \cle{cassante} : il y a rupture de la roche et déplacement relatif des deux compartiments (rejet non nul).
\item \textbf{Correcte.} Les diaclases (fractures sans rejet) découpent la roche en blocs ; sur un versant pentu, ces blocs peuvent se détacher : ce sont les éboulements et chutes de blocs.
\item \textbf{Correcte.} L'eau de pluie qui s'infiltre dans les plans de faille ou de stratification joue un rôle de \cle{lubrifiant} : elle diminue le frottement entre les surfaces rocheuses et favorise le glissement.
\end{enumerate}
\end{corrige}
\begin{attention}[Points de vigilance]
Ne pas confondre \textbf{faille} (fracture \emph{avec} déplacement) et \textbf{diaclase} (fracture \emph{sans} déplacement). Ne pas écrire qu'un pli se forme sous contrainte \og brutale \fg{} : c'est le contraire.
\end{attention}

\begin{corrige}[Exercice 2 — Vrai / Faux justifié : Rôle de l'eau et de l'homme]
\tcblower
\begin{enumerate}
\item \textbf{Faux.} L'infiltration de l'eau dans les fissures \textbf{diminue} la cohésion des roches tendres (argiles, marnes) par gonflement et lubrification, ce qui réduit leur résistance et favorise les glissements.
\item \textbf{Vrai.} L'action chimique de l'eau (dissolution du calcaire, altération des minéraux en argile) transforme des minéraux résistants en minéraux tendres : la résistance mécanique du massif \textbf{diminue}.
\item \textbf{Faux.} Le déboisement \textbf{expose} le sol à l'érosion par le ruissellement et \textbf{déstabilise} les pentes : les racines, qui retenaient le sol comme un maillage naturel, disparaissent.
\item \textbf{Vrai.} Construire en haut de pente ajoute une \textbf{surcharge} qui pousse le terrain vers le bas ; creuser en bas de pente supprime le \textbf{contrefort} naturel qui retenait le versant. Les deux actions peuvent déclencher un glissement.
\end{enumerate}
\end{corrige}
\begin{attention}[Points de vigilance]
La justification est exigée : une réponse \og vrai \fg{} ou \og faux \fg{} seule ne rapporte aucun point. Toujours citer le mécanisme (gonflement, lubrification, érosion, surcharge).
\end{attention}

\begin{corrige}[Exercice 3 — Définitions et vocabulaire technique]
\tcblower
\begin{enumerate}
\item \textbf{Mouvement de terrain :} déplacement, plus ou moins rapide, d'une masse de roches ou de sols le long d'une pente, sous l'effet de la gravité, lorsque la pente devient instable.
\item \textbf{Déformation souple :} déformation des roches \textbf{sans rupture} (courbure), se produisant sous l'effet de contraintes lentes et prolongées ; exemple : le \cle{pli}.
\item \textbf{Déformation cassante :} déformation des roches \textbf{avec rupture}, sous l'effet de contraintes rapides ou fortes ; exemples : la \cle{faille} (avec déplacement) et la \cle{diaclase} (sans déplacement).
\item \textbf{Diaclase :} fracture de la roche \textbf{sans déplacement} relatif des deux compartiments ; les réseaux de diaclases découpent la roche en blocs.
\item \textbf{Angle de repos :} angle maximal par rapport à l'horizontale qu'un matériau éboulé peut conserver sans s'effondrer à nouveau ; il caractérise l'équilibre naturel du matériau.
\item \textbf{Protection active :} technique qui agit \textbf{sur la cause} du risque pour stabiliser le versant (drainage, soutènement, reboisement). \textbf{Protection passive :} technique qui \textbf{protège les personnes et les biens} sans stabiliser le versant (filets pare-blocs, méridiennes, déviation de route).
\end{enumerate}
\end{corrige}
\begin{attention}[Points de vigilance]
Une définition doit comporter l'idée essentielle soulignée ici en gras : par exemple, pour la diaclase, l'absence de déplacement est le critère qui la distingue de la faille.
\end{attention}

\begin{corrige}[Exercice 4 — Calcul direct : Angle de pente et stabilité]
\tcblower
\begin{enumerate}
\item L'angle de \SI{35}{\degree} atteint après l'effondrement est l'\cle{angle de repos} (angle de talus naturel) de cette argile humide : c'est l'inclinaison pour laquelle le matériau éboulé est de nouveau stable.
\item Calcul de l'angle de frottement interne :
\[ \mu = \tan\varphi \quad \Longrightarrow \quad \varphi = \arctan(0.70) \approx \SI{35.0}{\degree} \]
L'angle de frottement calculé (\SI{35}{\degree}) est \textbf{égal} à l'angle de repos observé (\SI{35}{\degree}).
\item La pente stabilisée (\SI{35}{\degree}) étant égale à l'angle de frottement $\varphi$, le talus est en \textbf{équilibre limite} : il est stable pour l'instant, mais tout facteur supplémentaire (nouvelles pluies, surcharge, creusement du pied) peut le faire glisser à nouveau. Il faut donc adoucir la pente en dessous de \SI{35}{\degree} pour obtenir une marge de sécurité.
\end{enumerate}
\end{corrige}
\begin{attention}[Points de vigilance]
Vérifier que la calculatrice est en mode \textbf{degrés} pour $\arctan$. Retenir la règle : pente $> \varphi$ $\Rightarrow$ instable ; pente $< \varphi$ $\Rightarrow$ stable.
\end{attention}

\begin{corrige}[Exercice 5 — Analyse d'un affleurement : Plis, failles et risque d'éboulement]
\tcblower
\begin{enumerate}
\item \textbf{Structures bleues (plis) :} déformation \textbf{souple} — les couches sont courbées de façon continue, sans rupture. \textbf{Structure orange (faille) :} déformation \textbf{cassante} — il y a rupture nette et déplacement relatif des deux compartiments (rejet visible).
\item La structure en creux (forme en $\cup$) au niveau de la faille est un \textbf{synclinal} ; les structures en bosse (forme en $\cap$) de part et d'autre sont des \textbf{anticlinaux}.
\item Les blocs hachurés risquent un \textbf{éboulement (chute de blocs)} car : (1) ils reposent sur un versant pentu ; (2) le massif est découpé par la faille et les diaclases, ce qui isole les blocs ; (3) la faille a créé une zone de roche broyée et affaiblie qui constitue une surface de glissement potentielle.
\item C'est une \textbf{faille normale} : le compartiment du \textbf{toit} (situé au-dessus du plan de faille) est \textbf{descendu} par rapport au compartiment du \textbf{mur}. Ce jeu correspond à un étirement (extension) de la roche.
\end{enumerate}
\end{corrige}
\begin{attention}[Points de vigilance]
Pour identifier toit et mur : le toit est le compartiment situé \emph{au-dessus} du plan de faille incliné. Faille normale = toit descend ; faille inverse = toit monte.
\end{attention}

\begin{corrige}[Exercice 6 — Interprétation de courbes : Pluviométrie et glissements de terrain]
\tcblower
\begin{enumerate}
\item \textbf{Graphique attendu :} en abscisses les 12 mois ; histogramme des pluies (axe de gauche, en mm) ; courbe en points reliés du nombre de glissements (axe de droite). Les deux axes doivent être légendés avec leurs unités.
\item \textbf{Évolution parallèle :} les deux grandeurs évoluent globalement ensemble : aucun glissement en saison sèche (janvier, février, décembre), début en mars avec les premières pluies, maxima en \textbf{mai} (8 glissements pour \SI{210}{mm}) et \textbf{octobre} (9 glissements pour \SI{250}{mm}).
\item \textbf{Décalage :} le maximum de glissements suit le maximum de pluie avec un retard d'environ \textbf{0 à 1 mois}. \textbf{Explication :} l'eau de pluie doit d'abord s'\textbf{infiltrer} dans le sol et \textbf{saturer} les couches profondes ; ce n'est que lorsque les plans de glissement sont mouillés et lubrifiés que le terrain se met en mouvement. D'où le délai entre la pluie et le glissement.
\item \textbf{Rôle déclencheur de l'eau :} l'eau est le \textbf{facteur déclencheur} principal : en mouillant les argiles (gonflement) et en lubrifiant les plans de faille et de stratification, elle \textbf{diminue le frottement} qui retenait la masse de terrain ; le versant bascule alors de l'équilibre à l'instabilité.
\end{enumerate}
\end{corrige}
\begin{attention}[Points de vigilance]
Dans l'interprétation d'un graphique, toujours citer des \textbf{valeurs chiffrées} (pics, dates) pour appuyer la conclusion. Ne pas écrire \og l'eau fait glisser le terrain \fg{} sans préciser le mécanisme (infiltration, saturation, lubrification).
\end{attention}

\begin{corrige}[Exercice 7 — Étude de cas comparée : Versant boisé vs Versant urbanisé]
\tcblower
\begin{enumerate}
\item \textbf{Trois facteurs anthropiques sur le versant B :}
      \begin{enumerate}
      \item \textbf{Imperméabilisation} (toitures, routes, dalles béton) : l'eau ne peut plus s'infiltrer ; elle ruisselle et s'accumule en bas de pente, saturant rapidement les argiles.
      \item \textbf{Fuites des canalisations d'eaux usées} : apport d'eau permanent et concentré dans le sous-sol, qui maintient le terrain saturé même hors saison des pluies.
      \item \textbf{Surcharge et modifications de la pente} : poids des maisons et des remblais en haut de versant, entaillage du pied pour les routes (perte du contrefort naturel).
      \item (Complément) \textbf{Déboisement} : disparition des arbres qui retenaient le sol.
      \end{enumerate}
\item \textbf{Rôle mécanique des racines (versant A) :}
      \begin{itemize}
      \item \textbf{Renforcement :} les racines traversent le sol comme un maillage et le \og cousent \fg{} : elles augmentent la cohésion du terrain, à la manière de clous naturels.
      \item \textbf{Effet de butée :} les troncs et grosses racines, enfoncés dans le sous-sol stable, s'opposent au glissement des couches superficielles.
      \end{itemize}
\item \textbf{Rôle hydrologique :}
      \begin{itemize}
      \item \textbf{Forêt (A) :} le feuillage et la litière \textbf{interceptent} une partie de la pluie ; les racines \textbf{pompent} l'eau du sol et la rejettent dans l'air (évapotranspiration) $\Rightarrow$ le sol reste modérément humide.
      \item \textbf{Urbanisation (B) :} imperméabilisation $\Rightarrow$ ruissellement concentré ; fuites des réseaux $\Rightarrow$ infiltration localisée $\Rightarrow$ \textbf{saturation rapide} des argiles et perte de résistance.
      \end{itemize}
\item \textbf{Conclusion :} la couverture végétale est une \textbf{mesure de prévention active} efficace, peu coûteuse et durable : elle agit à la fois sur les causes hydrauliques (régulation de l'eau) et mécaniques (renforcement par les racines) de l'instabilité.
\end{enumerate}
\end{corrige}
\begin{attention}[Points de vigilance]
Bien distinguer dans la rédaction les deux rôles de la forêt : \textbf{mécanique} (racines = armature) et \textbf{hydrologique} (interception + évapotranspiration). C'est cette double action qui explique la différence entre les deux versants.
\end{attention}

\begin{corrige}[Exercice 8 — QCM de synthèse : Techniques de prévention / protection]
\tcblower
\begin{enumerate}
\item \textbf{Correspondances :}
      \begin{center}
      \begin{tabular}{|c|l|}\hline
      \textbf{Technique} & \textbf{Cause combattue} \\ \hline
      1. Drainage profond & C (excès d'eau dans le sol) \\ \hline
      2. Soutènement & A et D (surcharge en haut ; érosion du pied) \\ \hline
      3. Reboisement & B et C (érosion, ruissellement ; excès d'eau) \\ \hline
      4. Purge en haut de pente & A (surcharge en haut de pente) \\ \hline
      5. Bénétage (paliers) & E (pente trop raide) \\ \hline
      6. Méridiennes / pare-blocs & F (chute de blocs, éboulements) \\ \hline
      \end{tabular}
      \end{center}
\item \textbf{Classification :}
      \begin{itemize}
      \item \textbf{Protection active} (agit sur la cause, stabilise le versant) : 1 (drainage), 2 (soutènement), 3 (reboisement), 4 (purge), 5 (bénétage).
      \item \textbf{Protection passive} (protège les biens et les personnes sans stabiliser) : 6 (méridiennes et filets pare-blocs).
      \end{itemize}
\end{enumerate}
\end{corrige}
\begin{attention}[Points de vigilance]
Le critère de classement est : la technique \textbf{supprime-t-elle ou réduit-elle la cause} du mouvement (active) ou se contente-t-elle d'\textbf{arrêter les blocs / protéger} (passive) ? Un mur pare-blocs ne stabilise pas le versant.
\end{attention}

\begin{corrige}[Exercice 9 — Sujet type BEPC : Rapport d'expertise pour un quartier instable à Yaoundé]
\tcblower
\begin{enumerate}
\item \textbf{Diagnostic :} il s'agit d'un \textbf{glissement de terrain}. \textbf{Justification (indices observés) :} fissures du sol \textbf{en arc de cercle} (fissure de tête du glissement), fissures diagonales dans les murs des maisons (contraintes de cisaillement), arbres \textbf{penchés vers l'aval} (entraînés par le déplacement du sol). Le terrain argileux (marnes) sur pente douce est favorable à ce type de mouvement.
\item \textbf{Trois causes probables :}
      \begin{enumerate}
      \item \textbf{Saturation en eau des argiles (cause naturelle) :} les pluies de septembre-octobre ont infiltré le sol ; les argiles ont gonflé et perdu leur cohésion ; l'eau a lubrifié le plan de glissement.
      \item \textbf{Surcharge en tête de versant (cause humaine) :} le poids des maisons et des remblais (terrain de football) en haut de pente pousse le terrain vers le bas.
      \item \textbf{Affaiblissement du pied de pente :} creusements (routes, caniveaux) et fissures de retrait-gonflement des argiles réduisent le contrefort naturel.
      \end{enumerate}
\item \textbf{Mesures d'urgence (protection passive) :}
      \begin{enumerate}
      \item \textbf{Évacuer immédiatement} les habitations les plus fissurées et mettre la population en sécurité.
      \item \textbf{Interdire l'accès} à la zone instable par un périmètre de sécurité (rubans, surveillance), afin d'éviter toute victime et toute surcharge supplémentaire.
      \end{enumerate}
\item \textbf{Mesures de stabilisation (protection active) :}
      \begin{enumerate}
      \item \textbf{Drainage} (fossés, drains) : évacuer l'eau du sous-sol $\Rightarrow$ les argiles s'assèchent, retrouvent leur cohésion, et les plans de glissement ne sont plus lubrifiés. C'est la mesure prioritaire car l'eau est la cause déclenchante.
      \item \textbf{Purge / allègement de la tête} : enlever les terres instables et les surcharges en haut de pente $\Rightarrow$ diminution de la force qui pousse le terrain.
      \item \textbf{Soutènement ou reboisement} : construire un mur de soutènement au pied (butée) et/ou replanter des arbres dont les racines renforceront le sol et pomperont l'eau.
      \end{enumerate}
\item \textbf{Recommandation administrative :} \og Interdiction de toute nouvelle construction dans la zone instable ; réalisation des travaux de drainage et de purge ; surveillance régulière de la zone (fissures, déplacements) avant toute réoccupation des lieux. \fg
\end{enumerate}
\textbf{Barème indicatif (sur 20) :} diagnostic justifié (4 pts) ; trois causes avec mécanismes (6 pts) ; mesures d'urgence (3 pts) ; mesures de stabilisation justifiées (5 pts) ; recommandation claire (2 pts).
\end{corrige}
\begin{attention}[Points de vigilance]
Au BEPC, chaque mesure proposée doit être \textbf{reliée à une cause} identifiée : \og drainer car l'eau sature les argiles \fg. Une liste de mesures sans justification perd la moitié des points. Bien distinguer urgence (protéger les personnes) et stabilisation (traiter la cause).
\end{attention}

\begin{corrige}[Exercice 10 — Problème de synthèse : Aménagement d'un talus routier éboulé (RN1)]
\tcblower
\begin{enumerate}
\item \textbf{Analyse des causes :}
      \begin{itemize}
      \item \textbf{Diaclases dans le grès :} elles découpent la roche en blocs ; la taille du talus à \SI{70}{\degree} a libéré ces blocs qui ne sont plus maintenus latéralement : ils basculent ou glissent.
      \item \textbf{Niveaux argileux :} les argiles gonflent au contact de l'eau et poussent les blocs de grès vers l'extérieur ; une fois mouillées, elles deviennent molles et constituent des \textbf{plans de glissement} privilégiés.
      \item \textbf{Eau (nappe et pluies d'août) :} l'eau s'infiltre dans les diaclases, sature les argiles, les fait gonfler et lubrifie les contacts entre couches $\Rightarrow$ perte de frottement et de cohésion.
      \item \textbf{Géométrie :} une pente de \SI{70}{\degree} est très supérieure à l'angle de repos des éboulis (\SI{35}{\degree}--\SI{40}{\degree}) : le talus était instable dès sa construction (surcreusement).
      \end{itemize}
\item \textbf{Projet d'aménagement (croquis à légender) :}
      \begin{enumerate}
      \item \textbf{Reprofilage} : adoucir la pente (par exemple \SI{45}{\degree}) en enlevant les matériaux instables en haut.
      \item \textbf{Bénétage} : créer 3 à 4 banquettes (paliers) régulièrement espacées pour casser la pente et retenir les blocs.
      \item \textbf{Drainage} : drains dans le massif pour abaisser la nappe ; fossé au pied du talus pour évacuer l'eau.
      \item \textbf{Soutènement} : clouage de la paroi ou mur de soutènement au pied.
      \item \textbf{Revégétalisation} : planter herbes et arbres sur les banquettes et en haut du talus.
      \item \textbf{Protection passive} : filet ou barrière pare-blocs entre le talus et la chaussée.
      \end{enumerate}
\item \textbf{Justification de quatre techniques :}
      \begin{enumerate}
      \item \textbf{Bénétage :} les banquettes réduisent la hauteur de chute libre et \textbf{interceptent les blocs} détachés des diaclases ; elles facilitent aussi le drainage et l'entretien.
      \item \textbf{Drains :} en abaissant la nappe, ils \textbf{assèchent les argiles} : fin du gonflement, retour de la cohésion et du frottement sur les plans de glissement. On agit sur la cause principale (l'eau).
      \item \textbf{Soutènement :} le mur ou le clouage apporte une \textbf{butée} au pied qui maintient les blocs de grès et compense le soutien perdu lors de l'éboulement.
      \item \textbf{Pare-blocs :} il \textbf{arrête les blocs} qui se détacheraient encore et protège directement les usagers de la RN1 (protection passive) pendant et après les travaux.
      \end{enumerate}
\item \textbf{Surveillance post-travaux :}
      \begin{enumerate}
      \item \textbf{Mesure des déplacements} (repères sur la paroi, fissuromètres) : vérifier que le talus ne bouge plus.
      \item \textbf{Mesure du niveau de la nappe} (piézomètres) : vérifier que le drainage reste efficace, surtout pendant la saison des pluies.
      \end{enumerate}
\end{enumerate}
\textbf{Barème indicatif (sur 20) :} analyse des 4 causes (6 pts) ; croquis d'aménagement légendé (6 pts) ; justification de 4 techniques (6 pts) ; surveillance (2 pts).
\end{corrige}
\begin{attention}[Points de vigilance]
L'ordre logique est exigé : d'abord \textbf{sécuriser} (pare-blocs), puis \textbf{traiter la cause} (drainage, reprofilage), enfin \textbf{surveiller}. Toute technique proposée doit être justifiée par un élément de l'énoncé (diaclases, argiles, nappe, pente de \SI{70}{\degree}).
\end{attention}

\begin{corrige}[Exercice 11 — Sujet type BEPC : Synthèse documentaire — Glissement de Mbankolo (octobre 2023)]
\tcblower
\begin{enumerate}
\item \textbf{Contexte géologique :}
      \begin{itemize}
      \item \textbf{Base :} schistes sains ($> \SI{15}{m}$), durs et imperméables.
      \item \textbf{Formations superficielles :} schistes altérés fissurés (\SI{12}{--15}{m}), puis une épaisse couche d'\textbf{argiles plastiques} (\SI{3}{--12}{m}), coiffée d'une \textbf{latérite dure et fissurée} (\SI{0}{--3}{m}).
      \item \textbf{Plan de rupture probable :} à la \textbf{base de la couche d'argiles} (vers \SI{12}{m}), au contact avec les schistes.
      \item \textbf{Justification :} les argiles sont le niveau le plus faible (cohésion $c' = \SI{5}{kPa}$ très faible, angle de frottement $\varphi' = \SI{18}{\degree}$ faible). De plus, le toit des schistes sains est \textbf{imperméable} : l'eau qui s'infiltre à travers la latérite fissurée s'y accumule et sature la base des argiles, qui deviennent un plan savonneux.
      \end{itemize}
\item \textbf{Facteur déclencheur :} la pluviométrie exceptionnelle d'octobre 2023 (\SI{420}{mm} contre \SI{280}{mm} en moyenne, pointes de \SI{45}{mm/h}) a provoqué une infiltration massive à travers la latérite fissurée. Le sol étant saturé dès le 05/10 (Doc. 1), l'eau s'est accumulée sur le niveau imperméable des schistes. \textbf{Conséquence :} les argiles saturées ont perdu leur cohésion et l'eau a lubrifié le contact argile/schiste : la résistance du terrain est devenue inférieure aux forces de pesanteur sur la pente $\Rightarrow$ rupture et glissement.
\item \textbf{Facteurs aggravants humains (Doc. 3) :}
      \begin{enumerate}
      \item \textbf{Imperméabilisation} (toitures, routes) : l'eau ruisselle au lieu de s'infiltrer partout ; elle se concentre dans les fossés et les bas-fonds.
      \item \textbf{Absence d'assainissement} (fossés à ciel ouvert bouchés) : l'eau stagne et s'infiltre durablement dans le sous-sol.
      \item \textbf{Surcharge en haut de pente} : décharge sauvage et maisons en parpaings ajoutent du poids qui pousse le terrain vers le bas.
      \item \textbf{Déboisement} : la disparition des eucalyptus a supprimé le pompage de l'eau par les racines et le maintien du sol par le maillage racinaire.
      \end{enumerate}
\item \textbf{Typologie du mouvement :} glissement \textbf{rapide} (\og en une nuit \fg, déplacement de \SI{2}{m}, Doc. 4), se produisant le long d'un \textbf{plan de glissement} (contact argile/schiste). La phrase \og l'eau sortait du sol \fg{} confirme que la nappe affleurait au pied du glissement, signe d'une saturation totale. La carapace de latérite rigide s'est fracturée (mur de clôture craqué) et a glissé en blocs sur les argiles ramollies.
\item \textbf{Mesures correctives priorisées :}
      \begin{center}
      \begin{tabularx}{\linewidth}{|l|X|X|}\hline
      \textbf{Étape} & \textbf{Technique} & \textbf{Justification (documents)} \\ \hline
      Urgence & Évacuation des habitants ; interdiction d'accès ; protection pare-blocs au pied ; surveillance. & Doc. 4 : maisons déjà déplacées, vies en danger. \\ \hline
      Consolidation & \textbf{Drainage} du versant (drains, fossés étanches) pour évacuer l'eau et assécher les argiles. & Doc. 1 et 2 : la saturation des argiles est la cause déclenchante ; $c' = \SI{5}{kPa}$ très faible. \\ \hline
      Durable & Allègement de la tête (retrait de la décharge) ; réseau d'assainissement ; reboisement ; interdiction de construire dans la zone. & Doc. 3 : traite les causes humaines (surcharge, fuites, déboisement). \\ \hline
      \end{tabularx}
      \end{center}
\end{enumerate}
\textbf{Barème indicatif (sur 20) :} contexte géologique et plan de rupture justifié (5 pts) ; facteur déclencheur argumenté avec les chiffres du Doc. 1 (4 pts) ; 4 facteurs humains reliés au Doc. 3 (4 pts) ; typologie argumentée avec le Doc. 4 (3 pts) ; mesures priorisées et justifiées (4 pts).
\end{corrige}
\begin{attention}[Points de vigilance]
Dans une synthèse documentaire, \textbf{chaque affirmation doit s'appuyer sur un document cité} (\og d'après le Doc. 2... \fg). Utiliser les valeurs numériques fournies ($c'$, $\varphi'$, \SI{420}{mm}) : c'est ce qui distingue une copie excellente d'une copie moyenne. Ne pas inventer de causes absentes des documents.
\end{attention}

\newpage

\coursec{Fiche bilan}

\begin{popfigure}
\begin{center}\tcbox[colback=popblue,colframe=popblue,coltext=white,arc=9pt,boxsep=4pt,left=6pt,right=6pt]{\bfseries\small Mouvements de terrain}\end{center}
\vspace{-2pt}
\begin{tcbraster}[raster columns=3,raster equal height=rows,raster column skip=6pt,raster row skip=6pt,enhanced,blankest]
\begin{tcolorbox}[enhanced,colback=white,colframe=poporange,boxrule=0.9pt,arc=3pt,title={Définition et exemples},fonttitle=\bfseries\scriptsize,coltitle=white,colbacktitle=poporange,left=4pt,right=4pt,top=2pt,bottom=2pt,boxsep=1pt,toptitle=1.5pt,bottomtitle=1.5pt,halign=left]
\begin{itemize}[nosep,leftmargin=*,itemsep=1pt,font=\scriptsize]
\item Glissements, éboulements, coulées
\end{itemize}
\end{tcolorbox}
\begin{tcolorbox}[enhanced,colback=white,colframe=popgreen,boxrule=0.9pt,arc=3pt,title={Déformations des roches},fonttitle=\bfseries\scriptsize,coltitle=white,colbacktitle=popgreen,left=4pt,right=4pt,top=2pt,bottom=2pt,boxsep=1pt,toptitle=1.5pt,bottomtitle=1.5pt,halign=left]
\begin{itemize}[nosep,leftmargin=*,itemsep=1pt,font=\scriptsize]
\item Souples : plis
\item Cassantes : failles, fractures
\end{itemize}
\end{tcolorbox}
\begin{tcolorbox}[enhanced,colback=white,colframe=popblue!80,boxrule=0.9pt,arc=3pt,title={Rôle de l'eau},fonttitle=\bfseries\scriptsize,coltitle=white,colbacktitle=popblue!80,left=4pt,right=4pt,top=2pt,bottom=2pt,boxsep=1pt,toptitle=1.5pt,bottomtitle=1.5pt,halign=left]
\begin{itemize}[nosep,leftmargin=*,itemsep=1pt,font=\scriptsize]
\item Action mécanique : ruissellement, alourdissement
\item Action chimique : dissolution
\end{itemize}
\end{tcolorbox}
\begin{tcolorbox}[enhanced,colback=white,colframe=popgold,boxrule=0.9pt,arc=3pt,title={Action de l'homme},fonttitle=\bfseries\scriptsize,coltitle=white,colbacktitle=popgold,left=4pt,right=4pt,top=2pt,bottom=2pt,boxsep=1pt,toptitle=1.5pt,bottomtitle=1.5pt,halign=left]
\begin{itemize}[nosep,leftmargin=*,itemsep=1pt,font=\scriptsize]
\item Déboisement, constructions, cultures en pente
\end{itemize}
\end{tcolorbox}
\begin{tcolorbox}[enhanced,colback=white,colframe=poppurple,boxrule=0.9pt,arc=3pt,title={Prévention et protection},fonttitle=\bfseries\scriptsize,coltitle=white,colbacktitle=poppurple,left=4pt,right=4pt,top=2pt,bottom=2pt,boxsep=1pt,toptitle=1.5pt,bottomtitle=1.5pt,halign=left]
\begin{itemize}[nosep,leftmargin=*,itemsep=1pt,font=\scriptsize]
\item Reboisement, drainage, murs de soutènement
\end{itemize}
\end{tcolorbox}
\begin{tcolorbox}[enhanced,colback=white,colframe=poppink,boxrule=0.9pt,arc=3pt,title={Analyse d'un accident},fonttitle=\bfseries\scriptsize,coltitle=white,colbacktitle=poppink,left=4pt,right=4pt,top=2pt,bottom=2pt,boxsep=1pt,toptitle=1.5pt,bottomtitle=1.5pt,halign=left]
\begin{itemize}[nosep,leftmargin=*,itemsep=1pt,font=\scriptsize]
\item Observer, identifier les causes, proposer des solutions
\end{itemize}
\end{tcolorbox}
\end{tcbraster}

\legende{Carte mentale de la leçon : les mouvements de terrain, leurs causes (déformations des roches, action de l'eau, action de l'homme) et les moyens de prévention et de protection.}
\end{popfigure}

\begin{bilan}
\textbf{Définitions clés}
\begin{itemize}
  \item \cle{Mouvement de terrain} : déplacement de masses de roches ou de sols sur une pente, sous l'effet de la pesanteur. Exemples : \cle{glissement de terrain}, \cle{éboulement} (chute de blocs), \cle{coulée de boue}.
  \item \cle{Déformation souple} : la roche se plie sans se rompre ; elle donne des \cle{plis}. Elle a lieu lorsque la contrainte est lente et la roche assez souple.
  \item \cle{Déformation cassante} : la roche se rompt ; elle donne des \cle{fractures} (sans déplacement des compartiments) et des \cle{failles} (avec déplacement des compartiments).
  \item \cle{Faille} : fracture de la roche accompagnée d'un déplacement relatif des deux compartiments qu'elle sépare.
\end{itemize}

\textbf{Causes des mouvements de terrain}
\begin{center}
\begin{tabularx}{\linewidth}{|l|X|}
\hline
\rowcolor{popblueL} \textbf{Cause} & \textbf{Mécanisme à retenir} \\
\hline
Déformations cassantes & Les failles et fractures fragilisent les roches et créent des surfaces de glissement. \\
\hline
Action mécanique de l'eau & Le ruissellement érode les pentes ; l'eau de pluie alourdit et ramollit le sol, qui glisse plus facilement. \\
\hline
Action chimique de l'eau & L'eau dissout certaines roches (ex. calcaire), ce qui fragilise le terrain. \\
\hline
Action de l'homme & Déboisement (les racines ne retiennent plus le sol), constructions sur les pentes, cultures qui dénudent le sol, trafic et vibrations. \\
\hline
\end{tabularx}
\end{center}

\textbf{Prévention et protection}
\begin{itemize}
  \item \cle{Reboisement} des pentes : les racines des arbres fixent le sol.
  \item \cle{Drainage} : évacuer l'eau des terrains pour éviter l'alourdissement des sols.
  \item \cle{Murs de soutènement} et filets de protection au pied des pentes.
  \item Interdire les constructions dans les zones à risque ; respecter les zones protégées.
  \item Informer et sensibiliser les populations des zones montagneuses.
\end{itemize}

\textbf{Méthode express : analyser un accident de mouvement de terrain}
\begin{enumerate}
  \item Observer le document (photo, texte) et décrire les faits.
  \item Identifier le type de mouvement (glissement, éboulement, coulée).
  \item Rechercher les causes : nature des roches (failles ?), présence d'eau, action de l'homme.
  \item Proposer des mesures de prévention ou de protection adaptées.
  \item Conclure en répondant clairement à la question posée.
\end{enumerate}
\end{bilan}

\begin{aretenir}[Checklist avant le BEPC]
\begin{itemize}
  \item $\square$ Je sais définir un mouvement de terrain et citer au moins deux exemples.
  \item $\square$ Je sais distinguer une déformation souple (pli) d'une déformation cassante (fracture, faille).
  \item $\square$ Je sais définir une faille et la reconnaître sur un schéma ou une photo.
  \item $\square$ Je connais l'action mécanique de l'eau : ruissellement, alourdissement et ramollissement des sols.
  \item $\square$ Je connais l'action chimique de l'eau : dissolution des roches.
  \item $\square$ Je sais expliquer comment le déboisement favorise les glissements de terrain.
  \item $\square$ Je sais citer les actions de l'homme qui aggravent les risques (constructions, cultures en pente).
  \item $\square$ Je connais au moins trois techniques de prévention : reboisement, drainage, murs de soutènement.
  \item $\square$ Je sais rédiger une démarche d'analyse d'un accident en étapes ordonnées.
  \item $\square$ Je sais légender un schéma de pli et de faille avec un vocabulaire exact.
  \item $\square$ Je relie chaque cause à une mesure de protection adaptée dans mes réponses.
\end{itemize}
\end{aretenir}

\findecours

\end{document}
